\documentclass[12pt]{book}
% Idioma y codificación
\usepackage[spanish, es-tabla, es-notilde]{babel}       %es-tabla para que se titule "Tabla"
\usepackage[utf8]{inputenc}

% Márgenes
\usepackage[a4paper,top=3cm,bottom=2.5cm,left=3cm,right=3cm]{geometry}

% Comentarios de bloque
\usepackage{verbatim}

% Paquetes de links
\usepackage[hidelinks]{hyperref}    % Permite enlaces
\usepackage{url}                    % redirecciona a la web

% Más opciones para enumeraciones
\usepackage{enumitem}

% Personalizar la portada
\usepackage{titling}

% Paquetes de tablas
\usepackage{multirow}

% Para añadir el símbolo de euro
\usepackage{eurosym}

% Paquete para subrayar con diferentes estilos
\usepackage[normalem]{ulem}


%------------------------------------------------------------------------

%Paquetes de figuras
\usepackage{caption}
\usepackage{subcaption} % Figuras al lado de otras
\usepackage{float}      % Poner figuras en el sitio indicado H.


% Paquetes de imágenes
\usepackage{graphicx}       % Paquete para añadir imágenes
\usepackage{transparent}    % Para manejar la opacidad de las figuras

% Paquete para usar colores
\usepackage[dvipsnames, table, xcdraw]{xcolor}
\usepackage{pagecolor}      % Para cambiar el color de la página

% Habilita tamaños de fuente mayores
\usepackage{fix-cm}

% Para los gráficos
\usepackage{tikz}
\usepackage{forest}

% Para poder situar los nodos en los grafos
\usetikzlibrary{positioning}

% Para evitar problemas entre babel y tikz
\usetikzlibrary{babel}


%------------------------------------------------------------------------

% Paquetes de matemáticas
\usepackage{mathtools, amsfonts, amssymb, mathrsfs}
\usepackage[makeroom]{cancel}     % Simplificar tachando
\usepackage{polynom}    % Divisiones y Ruffini
\usepackage{units} % Para poner fracciones diagonales con \nicefrac

\usepackage{pgfplots}   %Representar funciones
\pgfplotsset{compat=1.18}  % Versión 1.18

\usepackage{tikz-cd}    % Para usar diagramas de composiciones
\usetikzlibrary{calc}   % Para usar cálculo de coordenadas en tikz

%Definición de teoremas, etc.
\usepackage{amsthm}
%\swapnumbers   % Intercambia la posición del texto y de la numeración

\theoremstyle{plain}

\makeatletter
\@ifclassloaded{article}{
  % --- Configuración para ARTICLE ---
  \newtheorem{teo}{Teorema}[section]
  \theoremstyle{definition}
  \newtheorem{definicion}{Definición}[section]
  \newtheorem{ejercicio}{Ejercicio}[section]

  \renewcommand{\theteo}{%
    \ifnum\value{section}=0 \arabic{teo}\else \thesection.\arabic{teo}\fi
  }
  \renewcommand{\thedefinicion}{%
    \ifnum\value{section}=0 \arabic{definicion}\else \thesection.\arabic{definicion}\fi
  }
  \renewcommand{\theejercicio}{%
    \ifnum\value{section}=0 \arabic{ejercicio}\else \thesection.\arabic{ejercicio}\fi
  }
}{
  % --- Configuración para BOOK / REPORT ---
  \newtheorem{teo}{Teorema}[chapter]
  \theoremstyle{definition}
  \newtheorem{definicion}{Definición}[chapter]
  \newtheorem{ejercicio}{Ejercicio}[section]

  % Si no hay sección: Capítulo.Teorema (ej. 2.1)
  % Si hay sección:    Capítulo.Sección.Teorema (ej. 2.3.1)
  \renewcommand{\theteo}{%
    \ifnum\value{section}=0
      \thechapter.\arabic{teo}%
    \else
      \thesection.\arabic{teo}% (\thesection ya contiene el número de capítulo)
    \fi
  }
  \renewcommand{\thedefinicion}{%
    \ifnum\value{section}=0
      \thechapter.\arabic{definicion}%
    \else
      \thesection.\arabic{definicion}%
    \fi
  }
  \renewcommand{\theejercicio}{%
    \ifnum\value{section}=0
      \arabic{ejercicio}%
    \else
      \thesection.\arabic{ejercicio}% Formato por defecto: ej. 2.2.1
    \fi
  }
}
\makeatother

% Teoremas dependientes o con contador compartido
\newtheorem{coro}{Corolario}[teo]
\newtheorem{prop}[teo]{Proposición}
\newtheorem{lema}[teo]{Lema}
\newtheorem*{lema*}{Lema}

\theoremstyle{remark}
\newtheorem*{observacion}{Observación}

\theoremstyle{definition}
\newtheorem*{notacion}{Notación}
\newtheorem*{ejemplo}{Ejemplo}
\newtheorem*{ejercicio*}{Ejercicio}


% \renewcommand\qedsymbol{$\blacksquare$}         % Cambiar símbolo QED
%------------------------------------------------------------------------

% Paquetes para encabezados
\usepackage{fancyhdr}
\pagestyle{fancy}
\fancyhf{}

\newcommand{\helv}{ % Modificación tamaño de letra
\fontfamily{}\fontsize{12}{12}\selectfont}
\setlength{\headheight}{15pt} % Amplía el tamaño del índice


%\usepackage{lastpage}   % Referenciar última pag   \pageref{LastPage}
%\fancyfoot[C]{%
%  \begin{minipage}{\textwidth}
%    \centering
%    ~\\
%    \thepage\\
%    \href{https://losdeldgiim.github.io/}{\texttt{\footnotesize losdeldgiim.github.io}}
%  \end{minipage}
%}
\fancyfoot[C]{\thepage}
\fancyfoot[R]{\href{https://losdeldgiim.github.io/}{\texttt{\footnotesize losdeldgiim.github.io}}}

%------------------------------------------------------------------------

% Conseguir que no ponga "Capítulo 1". Sino solo "1."
\makeatletter
\@ifclassloaded{book}{
  \renewcommand{\chaptermark}[1]{\markboth{\thechapter.\ #1}{}} % En el encabezado
    
  \renewcommand{\@makechapterhead}[1]{%
  \vspace*{50\p@}%
  {\parindent \z@ \raggedright \normalfont
    \ifnum \c@secnumdepth >\m@ne
      \huge\bfseries \thechapter.\hspace{1em}\ignorespaces
    \fi
    \interlinepenalty\@M
    \Huge \bfseries #1\par\nobreak
    \vskip 40\p@
  }}
}
\makeatother

%------------------------------------------------------------------------
% Paquetes de cógido
\usepackage{minted}
\renewcommand\listingscaption{Código fuente}

\usepackage{fancyvrb}
% Personaliza el tamaño de los números de línea
\renewcommand{\theFancyVerbLine}{\small\arabic{FancyVerbLine}}

% Estilo para C++
\newminted{cpp}{
    frame=lines,
    framesep=2mm,
    baselinestretch=1.2,
    linenos,
    escapeinside=||
}

% para minted
\definecolor{LightGray}{rgb}{0.95,0.95,0.92}
\setminted{
    linenos=true,
    stepnumber=5,
    numberfirstline=true,
    autogobble,
    breaklines=true,
    breakautoindent=true,
    breaksymbolleft=,
    breaksymbolright=,
    breaksymbolindentleft=0pt,
    breaksymbolindentright=0pt,
    breaksymbolsepleft=0pt,
    breaksymbolsepright=0pt,
    fontsize=\footnotesize,
    bgcolor=LightGray,
    numbersep=10pt
}


\usepackage{listings} % Para incluir código desde un archivo

\renewcommand\lstlistingname{Código Fuente}
\renewcommand\lstlistlistingname{Índice de Códigos Fuente}

% Definir colores
\definecolor{vscodepurple}{rgb}{0.5,0,0.5}
\definecolor{vscodeblue}{rgb}{0,0,0.8}
\definecolor{vscodegreen}{rgb}{0,0.5,0}
\definecolor{vscodegray}{rgb}{0.5,0.5,0.5}
\definecolor{vscodebackground}{rgb}{0.97,0.97,0.97}
\definecolor{vscodelightgray}{rgb}{0.9,0.9,0.9}

% Configuración para el estilo de C similar a VSCode
\lstdefinestyle{vscode_C}{
  backgroundcolor=\color{vscodebackground},
  commentstyle=\color{vscodegreen},
  keywordstyle=\color{vscodeblue},
  numberstyle=\tiny\color{vscodegray},
  stringstyle=\color{vscodepurple},
  basicstyle=\scriptsize\ttfamily,
  breakatwhitespace=false,
  breaklines=true,
  captionpos=b,
  keepspaces=true,
  numbers=left,
  numbersep=5pt,
  showspaces=false,
  showstringspaces=false,
  showtabs=false,
  tabsize=2,
  frame=tb,
  framerule=0pt,
  aboveskip=10pt,
  belowskip=10pt,
  xleftmargin=10pt,
  xrightmargin=10pt,
  framexleftmargin=10pt,
  framexrightmargin=10pt,
  framesep=0pt,
  rulecolor=\color{vscodelightgray},
  backgroundcolor=\color{vscodebackground},
}

%------------------------------------------------------------------------
% La etiqueda IA sirve para indicar a partir de donde le pasamos el preambulo a la IA para el escaneo
% IA

% Comandos definidos
\newcommand{\bb}[1]{\mathbb{#1}}
\newcommand{\cc}[1]{\mathcal{#1}}

% I prefer the slanted \leq
\let\oldleq\leq % save them in case they're every wanted
\let\oldgeq\geq
\renewcommand{\leq}{\leqslant}
\renewcommand{\geq}{\geqslant}

% Si y solo si
\newcommand{\sii}{\iff}

% MCD y MCM
\DeclareMathOperator{\mcd}{mcd}
\DeclareMathOperator{\mcm}{mcm}

% Signo
\DeclareMathOperator{\sgn}{sgn}

% Imagen
\DeclareMathOperator{\Ima}{Im}

% Aplicación identidad
\newcommand{\id}{\mathrm{Id}}

% Letras griegas
\newcommand{\eps}{\epsilon}
\newcommand{\veps}{\varepsilon}
\newcommand{\lm}{\lambda}

\newcommand{\ol}{\overline}
\newcommand{\ul}{\underline}
\newcommand{\wt}{\widetilde}
\newcommand{\wh}{\widehat}

\let\oldvec\vec
\renewcommand{\vec}{\overrightarrow}

% Derivadas parciales
\newcommand{\del}[2]{\frac{\partial #1}{\partial #2}}
\newcommand{\Del}[3]{\frac{\partial^{#1} #2}{\partial #3^{#1}}}
\newcommand{\deld}[2]{\dfrac{\partial #1}{\partial #2}}
\newcommand{\Deld}[3]{\dfrac{\partial^{#1} #2}{\partial #3^{#1}}}


\newcommand{\AstIg}{\stackrel{(\ast)}{=}}
\newcommand{\Hop}{\stackrel{L'H\hat{o}pital}{=}}

\newcommand{\red}[1]{{\color{red}#1}} % Para integrales, destacar los cambios.

% Método de integración
\newcommand{\MetInt}[2]{
    \left[\begin{array}{c}
        #1 \\ #2
    \end{array}\right]
}

% Declarar aplicaciones
% 1. Nombre aplicación
% 2. Dominio
% 3. Codominio
% 4. Variable
% 5. Imagen de la variable
\newcommand{\Func}[5]{
    \begin{equation*}
        \begin{array}{rrll}
            \displaystyle #1:& \displaystyle  #2 & \longrightarrow & \displaystyle  #3\\
               & \displaystyle  #4 & \longmapsto & \displaystyle  #5
        \end{array}
    \end{equation*}
}

%------------------------------------------------------------------------


\newcommand{\myhref}[2]{\href{#1}{\color{blue}\uline{#2}}}

% Para poder incluir árboles
\usepackage{forest}
\usepackage{booktabs}

% Para poder hacer graficas en polares (espiral de arquimedes)
\usepackage{pgfplots}
\usepgfplotslibrary{polar}

\newcommand{\showimages}{} % Descomentar esta línea para mostrar las imágenes
\newcommand{\R}{\mathbb{R}}
\newcommand{\N}{\mathbb{N}}

\newcommand{\di}{\text{div}}
\newcommand{\rot}{\text{rot}}
\newcommand{\norma}[1]{\left\lVert #1 \right\rVert}
\newcommand{\normagenerica}{\lVert \cdot \rVert}
\newcommand{\prodescalar}[2]{\langle #1, #2 \rangle}
\newcommand{\prodescalargenerico}{\langle \cdot, \cdot \rangle}

\newcommand{\Rmas}{\mathbb{R}^{+}}
\newcommand{\sop}{\text{sop }}
\newcommand{\tr}{\text{tr}}


\begin{document}

    % 1. Foto de fondo
    % 2. Título
    % 3. Encabezado Izquierdo
    % 4. Color de fondo
    % 5. Coord x del titulo
    % 6. Coord y del titulo
    % 7. Fecha
    % 8. Autor

    
    % 1. Foto de fondo
% 2. Título
% 3. Encabezado Izquierdo
% 4. Color de fondo
% 5. Coord x del titulo
% 6. Coord y del titulo
% 7. Fecha
% 8. Autor

\newcommand{\portada}[8]{

    \portadaBase{#1}{#2}{#3}{#4}{#5}{#6}{#7}{#8}
    \portadaBook{#1}{#2}{#3}{#4}{#5}{#6}{#7}{#8}
}

\newcommand{\portadaExamen}[8]{

    \portadaBase{#1}{#2}{#3}{#4}{#5}{#6}{#7}{#8}
    \portadaArticle{#1}{#2}{#3}{#4}{#5}{#6}{#7}{#8}
}




\newcommand{\portadaBase}[8]{

    % Tiene la portada principal y la licencia Creative Commons
    
    % 1. Foto de fondo
    % 2. Título
    % 3. Encabezado Izquierdo
    % 4. Color de fondo
    % 5. Coord x del titulo
    % 6. Coord y del titulo
    % 7. Fecha
    % 8. Autor
    
    
    \thispagestyle{empty}               % Sin encabezado ni pie de página
    \newgeometry{margin=0cm}        % Márgenes nulos para la primera página
    
    
    % Encabezado
    \fancyhead[L]{\helv #3}
    \fancyhead[R]{\helv \nouppercase{\leftmark}}
    
    
    \pagecolor{#4}        % Color de fondo para la portada
    
    \begin{figure}[p]
        \centering
        \transparent{0.3}           % Opacidad del 30% para la imagen
        
        \includegraphics[width=\paperwidth, keepaspectratio]{../../_assets/#1}
    
        \begin{tikzpicture}[remember picture, overlay]
            \node[anchor=north west, text=white, opacity=1, font=\fontsize{60}{90}\selectfont\bfseries\sffamily, align=left] at (#5, #6) {#2};
            
            \node[anchor=south east, text=white, opacity=1, font=\fontsize{12}{18}\selectfont\sffamily, align=right] at (9.7, 3) {\href{https://losdeldgiim.github.io/}{\textbf{Los Del DGIIM}, \texttt{\footnotesize losdeldgiim.github.io}}};
            
            \node[anchor=south east, text=white, opacity=1, font=\fontsize{12}{15}\selectfont\sffamily, align=right] at (9.7, 1.8) {Doble Grado en Ingeniería Informática y Matemáticas\\Universidad de Granada};
        \end{tikzpicture}
    \end{figure}
    
    
    \restoregeometry        % Restaurar márgenes normales para las páginas subsiguientes
    \pagecolor{white}       % Restaurar el color de página
    
    
    \newpage
    \thispagestyle{empty}               % Sin encabezado ni pie de página
    \begin{tikzpicture}[remember picture, overlay]
        \node[anchor=south west, inner sep=3cm] at (current page.south west) {
            \begin{minipage}{0.5\paperwidth}
                \href{https://creativecommons.org/licenses/by-nc-nd/4.0/}{
                    \includegraphics[height=2cm]{../../_assets/Licencia.png}
                }\vspace{1cm}\\
                Esta obra está bajo una
                \href{https://creativecommons.org/licenses/by-nc-nd/4.0/}{
                    Licencia Creative Commons Atribución-NoComercial-SinDerivadas 4.0 Internacional (CC BY-NC-ND 4.0).
                }\\
    
                Eres libre de compartir y redistribuir el contenido de esta obra en cualquier medio o formato, siempre y cuando des el crédito adecuado a los autores originales y no persigas fines comerciales. 
            \end{minipage}
        };
    \end{tikzpicture}
    
    
    
    % 1. Foto de fondo
    % 2. Título
    % 3. Encabezado Izquierdo
    % 4. Color de fondo
    % 5. Coord x del titulo
    % 6. Coord y del titulo
    % 7. Fecha
    % 8. Autor


}


\newcommand{\portadaBook}[8]{

    % 1. Foto de fondo
    % 2. Título
    % 3. Encabezado Izquierdo
    % 4. Color de fondo
    % 5. Coord x del titulo
    % 6. Coord y del titulo
    % 7. Fecha
    % 8. Autor

    % Personaliza el formato del título
    \pretitle{\begin{center}\bfseries\fontsize{42}{56}\selectfont}
    \posttitle{\par\end{center}\vspace{2em}}
    
    % Personaliza el formato del autor
    \preauthor{\begin{center}\Large}
    \postauthor{\par\end{center}\vfill}
    
    % Personaliza el formato de la fecha
    \predate{\begin{center}\huge}
    \postdate{\par\end{center}\vspace{2em}}
    
    \title{#2}
    \author{\href{https://losdeldgiim.github.io/}{Los Del DGIIM, \texttt{\large losdeldgiim.github.io}}
    \\ \vspace{0.5cm}#8}
    \date{Granada, #7}
    \maketitle
    
    \tableofcontents
}




\newcommand{\portadaArticle}[8]{

    % 1. Foto de fondo
    % 2. Título
    % 3. Encabezado Izquierdo
    % 4. Color de fondo
    % 5. Coord x del titulo
    % 6. Coord y del titulo
    % 7. Fecha
    % 8. Autor

    % Personaliza el formato del título
    \pretitle{\begin{center}\bfseries\fontsize{42}{56}\selectfont}
    \posttitle{\par\end{center}\vspace{2em}}
    
    % Personaliza el formato del autor
    \preauthor{\begin{center}\Large}
    \postauthor{\par\end{center}\vspace{3em}}
    
    % Personaliza el formato de la fecha
    \predate{\begin{center}\huge}
    \postdate{\par\end{center}\vspace{5em}}
    
    \title{#2}
    \author{\href{https://losdeldgiim.github.io/}{Los Del DGIIM, \texttt{\large losdeldgiim.github.io}}
    \\ \vspace{0.5cm}#8}
    \date{Granada, #7}
    \thispagestyle{empty}               % Sin encabezado ni pie de página
    \maketitle
    \vfill
}
    \portada{ffccA4.jpg}{Ecuaciones\\Diferenciales\\ en Mecánica \\ y Biología}{Ecuaciones Diferenciales en Mecánica y Biología}{MidnightBlue}{-8}{28}{2026}{José Manuel Sánchez Varbas}

    \noindent
El temario se dividirá principalmente en tres bloques. El primer capítulo se corresponde con una breve introducción de ciertas herramientas necesarias para los otros tres. A partir del segundo comienza la susodicha división en bloques, cada uno de ellos correspondiente con un parcial de la asignatura.
El segundo capítulo se corresponde con el primer bloque, de transporte. El tercero con el segundo bloque, de mecánica de fluidos. Por último, el cuarto con el tercer bloque, relacionado con modelos en el ámbito de la biología.\newpage

\chapter{Introducción. Recordatorios}
\noindent

\section{Gradiente}

Consideramos $N \in \N$ fijo y $f : \R^N \to \R$ (campo escalar), así como $x = (x_1, \ldots, x_N) \in \R^N$. El gradiente 
recordemos que se define como $$\nabla f (x) = \left( \dfrac{\partial f}{\partial x_1}, \ldots, 
\dfrac{\partial f}{\partial x_N} \right)(x) \quad \text{\underline{VECTOR}}$$
Este vector nos indica la dirección de máximo cambio.

\begin{proof}
    Basta considerar $v \in \R^N, |v| = 1$. Recordando que $$D_vf = \langle \nabla f(x), v \rangle = |\nabla f(x)| 
    \overbrace{|v|}^{1} \cos(\widehat{\nabla f(x)}{v})$$
    Para que sea máximo (escalarmente), el coseno debe ser máximo (1), es decir, el ángulo entre el gradiente y 
    $v$ debe ser 0, por lo que el gradiente y $v$ son vectores paralelos, o sea, tienen la misma dirección.
    Por lo tanto, necesariamente\footnote{Suponiendo que $\nabla f(x) \neq 0$.}
    $$v = \dfrac{\nabla f(x)}{|\nabla f(x)|}$$
    y por definición de $v$, este tiene la dirección de máximo cambio.
\end{proof}

    Si consideramos un campo vectorial, $f: \R^N \to \R^N$, podemos considerar el jacobiano componente a componente. Así, si
    $f = (f_1, \ldots, f_N)$, tenemos que 
    $$\nabla f = \begin{pmatrix} 
        \dfrac{\partial f_1}{\partial x_1} & \cdots & \dfrac{\partial f_1}{\partial x_N} \\ 
        \vdots & \ddots & \vdots \\ 
        \dfrac{\partial f_N}{\partial x_1} & \cdots & \dfrac{\partial f_N}{\partial x_N} 
    \end{pmatrix} = J[f] \quad \text{\underline{MATRIZ JACOBIANA}}$$ 

\newpage

\section{Divergencia}\label{section:div}

\subsection{Motivación del concepto}

Supongamos que tenemos un fluido (líquido, por ejemplo, agua). Nos queremos fijar en una partícula (concreta pero arbitraria) 
del fluido, con la hipótesis de que dicho fluido se mueve en el tiempo. \\

Consideramos $u = (u_1, u_2, u_3)$, que será el campo de velocidades del fluido, a su vez supuesto conocido. Tenemos
$u = u(t,x)$, donde $t$ es la variable temporal y $q(t)$ es la posición de la partícula en el instante $t$. 
La situación queda representada en la Figura \ref{fig:f11}:

\begin{figure}[H]
    \begin{center}
        \tikzset{every picture/.style={line width=0.75pt}} %set default line width to 0.75pt        
        \begin{tikzpicture}[x=0.75pt,y=0.75pt,yscale=-1,xscale=1]
        %uncomment if require: \path (0,300); %set diagram left start at 0, and has height of 300

        %Shape: Circle [id:dp9868303673861135] 
        \draw  [fill={rgb, 255:red, 0; green, 0; blue, 0 }  ,fill opacity=1 ] (97.67,109.83) .. controls (97.67,108.45) and (98.79,107.33) .. (100.17,107.33) .. controls (101.55,107.33) and (102.67,108.45) .. (102.67,109.83) .. controls (102.67,111.21) and (101.55,112.33) .. (100.17,112.33) .. controls (98.79,112.33) and (97.67,111.21) .. (97.67,109.83) -- cycle ;
        %Curve Lines [id:da7650198021214434] 
        \draw    (100.17,109.83) .. controls (113,100.21) and (152.47,115.47) .. (144.33,87) .. controls (136.32,58.96) and (160.26,79.04) .. (187.74,61.18) ;
        \draw [shift={(189,60.33)}, rotate = 145.38] [color={rgb, 255:red, 0; green, 0; blue, 0 }  ][line width=0.75]    (10.93,-3.29) .. controls (6.95,-1.4) and (3.31,-0.3) .. (0,0) .. controls (3.31,0.3) and (6.95,1.4) .. (10.93,3.29)   ;
        %Shape: Circle [id:dp5572905211482128] 
        \draw  [fill={rgb, 255:red, 0; green, 0; blue, 0 }  ,fill opacity=1 ] (187.67,59.83) .. controls (187.67,58.45) and (188.79,57.33) .. (190.17,57.33) .. controls (191.55,57.33) and (192.67,58.45) .. (192.67,59.83) .. controls (192.67,61.21) and (191.55,62.33) .. (190.17,62.33) .. controls (188.79,62.33) and (187.67,61.21) .. (187.67,59.83) -- cycle ;

        % Text Node
        \draw (142.67,131.67) node   [align=left] {\begin{minipage}[lt]{68pt}\setlength\topsep{0pt}
        $\displaystyle q$
        \end{minipage}};
        % Text Node
        \draw (246,65.67) node   [align=left] {\begin{minipage}[lt]{68pt}\setlength\topsep{0pt}
        $\displaystyle q( t)$
        \end{minipage}};
        \end{tikzpicture}
    \end{center}
    \caption{Posición inicial e instantánea en el instante $t$}
    \label{fig:f11}
\end{figure}

La ley que nos describe este movimiento es el siguiente PVI\footnote{Para simplificar la notación, en lo que sigue escribiremos $u(q(t))\equiv u(t,q(t))$, 
entendiendo que la primera variable de $u$ es siempre el tiempo $t$.}

\begin{equation}\label{eq:ec1}
    \begin{cases}
        q'(t) = u(q(t)) \\
        q(0) = q
    \end{cases}
\end{equation}

Es decir, si conocemos el campo de velocidades, podemos escribir la ecuación, resolverla, y obtener la función $q(t)$
que nos describe el susodicho movimiento de la partícula. \\

Ahora consideramos la misma situación inicial, y queremos estudiar cómo cambia un volumen infinitesimal.
El tamaño de cada arista será $\varepsilon$, y queremos ver cómo se deforma el volumen del cubo con el tiempo.
Cuando podamos escribir dicha ecuación aparecerá de forma natural el concepto de divergencia. Las longitudes $a_q^i$, con
$i \in \{1,2,3\}$, son funciones que dependen del tiempo, pero que coinciden en el instante inicial. La situación en que nos 
encontramos viene ilustrada en la Figura \ref{fig:f12}:

\begin{figure}[H]
    \begin{center}

    \tikzset{every picture/.style={line width=0.75pt}} %set default line width to 0.75pt        

    \begin{tikzpicture}[x=0.75pt,y=0.75pt,yscale=-1,xscale=1]
    %uncomment if require: \path (0,300); %set diagram left start at 0, and has height of 300

    %Shape: Circle [id:dp9868303673861135] 
    \draw  [fill={rgb, 255:red, 0; green, 0; blue, 0 }  ,fill opacity=1 ] (207.67,99.83) .. controls (207.67,98.45) and (208.79,97.33) .. (210.17,97.33) .. controls (211.55,97.33) and (212.67,98.45) .. (212.67,99.83) .. controls (212.67,101.21) and (211.55,102.33) .. (210.17,102.33) .. controls (208.79,102.33) and (207.67,101.21) .. (207.67,99.83) -- cycle ;
    %Shape: Rectangle [id:dp8440612850022642] 
    \draw   (210,60) -- (250,60) -- (250,100) -- (210,100) -- cycle ;
    %Straight Lines [id:da8160136384076409] 
    \draw    (210,60) -- (230,40) ;
    %Straight Lines [id:da546497530395401] 
    \draw    (250,60) -- (270,40) ;
    %Straight Lines [id:da011934677438056873] 
    \draw    (250,100) -- (270,80) ;
    %Straight Lines [id:da28246078787903206] 
    \draw  [dash pattern={on 0.84pt off 2.51pt}]  (230,40) -- (230,80) ;
    %Straight Lines [id:da7718883113576894] 
    \draw  [dash pattern={on 0.84pt off 2.51pt}]  (230,80) -- (250,80) ;
    %Straight Lines [id:da7573142764220282] 
    \draw    (250,80) -- (270,80) ;
    %Straight Lines [id:da3645358326064383] 
    \draw    (270,40) -- (270,80) ;
    %Straight Lines [id:da04533787187694238] 
    \draw    (230,40) -- (270,40) ;
    %Straight Lines [id:da31951741461037897] 
    \draw  [dash pattern={on 0.84pt off 2.51pt}]  (210,100) -- (230,80) ;
    %Shape: Brace [id:dp616787207683195] 
    \draw   (210,100) .. controls (210,104.67) and (212.33,107) .. (217,107) -- (220,107) .. controls (226.67,107) and (230,109.33) .. (230,114) .. controls (230,109.33) and (233.33,107) .. (240,107)(237,107) -- (243,107) .. controls (247.67,107) and (250,104.67) .. (250,100) ;
    %Shape: Brace [id:dp44451260981777607] 
    \draw   (231,80) .. controls (228.25,77.12) and (225.44,77.05) .. (222.56,79.79) -- (222.56,79.79) .. controls (218.44,83.72) and (215.01,84.24) .. (212.26,81.35) .. controls (215.01,84.24) and (214.32,87.64) .. (210.21,91.56)(212.06,89.79) -- (210.21,91.56) .. controls (207.32,94.31) and (207.25,97.12) .. (210,100) ;
    %Curve Lines [id:da7568531294421236] 
    \draw  [dash pattern={on 0.84pt off 2.51pt}]  (212,80) .. controls (220.87,66.21) and (231.67,51.45) .. (201.41,36.67) ;
    \draw [shift={(200,36)}, rotate = 25.11] [color={rgb, 255:red, 0; green, 0; blue, 0 }  ][line width=0.75]    (10.93,-3.29) .. controls (6.95,-1.4) and (3.31,-0.3) .. (0,0) .. controls (3.31,0.3) and (6.95,1.4) .. (10.93,3.29)   ;
    %Shape: Brace [id:dp9587884219726038] 
    \draw   (210,60) .. controls (205.33,60.11) and (203.06,62.5) .. (203.17,67.17) -- (203.26,70.67) .. controls (203.42,77.34) and (201.17,80.73) .. (196.51,80.84) .. controls (201.17,80.73) and (203.58,84) .. (203.75,90.67)(203.67,87.67) -- (203.83,94.17) .. controls (203.94,98.84) and (206.33,101.11) .. (211,101) ;
    %Shape: Circle [id:dp8196257583758775] 
    \draw  [fill={rgb, 255:red, 0; green, 0; blue, 0 }  ,fill opacity=1 ] (247.67,98.83) .. controls (247.67,97.45) and (248.79,96.33) .. (250.17,96.33) .. controls (251.55,96.33) and (252.67,97.45) .. (252.67,98.83) .. controls (252.67,100.21) and (251.55,101.33) .. (250.17,101.33) .. controls (248.79,101.33) and (247.67,100.21) .. (247.67,98.83) -- cycle ;
    %Shape: Circle [id:dp9899510060278905] 
    \draw  [fill={rgb, 255:red, 0; green, 0; blue, 0 }  ,fill opacity=1 ] (207.67,60.83) .. controls (207.67,59.45) and (208.79,58.33) .. (210.17,58.33) .. controls (211.55,58.33) and (212.67,59.45) .. (212.67,60.83) .. controls (212.67,62.21) and (211.55,63.33) .. (210.17,63.33) .. controls (208.79,63.33) and (207.67,62.21) .. (207.67,60.83) -- cycle ;
    %Shape: Circle [id:dp14202122866986466] 
    \draw  [fill={rgb, 255:red, 0; green, 0; blue, 0 }  ,fill opacity=1 ] (227.67,79.83) .. controls (227.67,78.45) and (228.79,77.33) .. (230.17,77.33) .. controls (231.55,77.33) and (232.67,78.45) .. (232.67,79.83) .. controls (232.67,81.21) and (231.55,82.33) .. (230.17,82.33) .. controls (228.79,82.33) and (227.67,81.21) .. (227.67,79.83) -- cycle ;

    % Text Node
    \draw (201.33,114.83) node   [align=left] {\begin{minipage}[lt]{11.79pt}\setlength\topsep{0pt}
    $\displaystyle q$
    \end{minipage}};
    % Text Node
    \draw (217,116) node [anchor=north west][inner sep=0.75pt]   [align=left] {$\displaystyle a_{q}^{1}( t)$};
    % Text Node
    \draw (162,21) node [anchor=north west][inner sep=0.75pt]   [align=left] {$\displaystyle a_{q}^{2}( t)$};
    % Text Node
    \draw (161,67) node [anchor=north west][inner sep=0.75pt]   [align=left] {$\displaystyle a_{q}^{3}( t)$};

    \end{tikzpicture}
    \end{center}
    \caption{Cubo en el instante $t$}
    \label{fig:f12}
\end{figure}

\newpage

Así, planteamos las siguientes condiciones iniciales:

$$\begin{cases}
    a_q^1(0) = (\varepsilon, 0, 0) \\
    a_q^2(0) = (0, \varepsilon, 0) \\
    a_q^3(0) = (0, 0, \varepsilon)
\end{cases}$$

Buscamos entonces describir $a_q^i(t)$, para $i \in \{1,2,3\}$. Vemos que $q(t) + a_q^i(t)$ nos da la posición
de cada uno de los vértices (para cada $i$), y dicha evolución en el tiempo la medimos con su derivada
$$(q(t) + a_q^i(t))'.$$
De aquí deducimos una ley análoga a la que teníamos antes
$$(q(t) + a_q^i(t))' = u(q(t) + a_q^i(t)).$$
Aplicando la linealidad de la derivada, obtenemos
$$q'(t) + (a_q^i)'(t) = u(q(t) + a_q^i(t))$$
y usando \eqref{eq:ec1}~llegamos a 
$$u(q(t)) + (a_q^i)'(t) = u(q(t) + a_q^i(t)).$$

Nos disponemos en este punto a realizar un desarrollo de Taylor de la función $u$ en torno a $q(t)$, es decir\footnote{Suponemos que, al menos, $u(t, \cdot) \in C^2(\R^3, \R^3)$,
es decir, para cada $t$ instante de tiempo fijo, pongamos en $\left[0, T\right[$, el campo es $C^2$ respecto de $x$.)}

$$u(q(t) + a_q^i(t)) = u(q(t)) + \nabla u (q(t)) a_q^i(t) + O(a_q^i(t)^2)$$

La matriz jacobiana de $u$ ($\nabla u$) aparece debido a que $u$ es un campo vectorial. El segundo sumando se debe a que estamos haciendo la derivada de un campo vectorial, por lo que tenemos una matriz y 
el segundo factor es un vector\footnote{Será fundamental saber con lo que estamos trabajando (matriz, vector, etc...) 
durante el curso.}, y el término de error suponemos que es lo suficientemente pequeño como para poder 
despreciarlo. \\

Juntando todo, llegamos a

$$u(q(t)) + (a_q^i)'(t) = u(q(t)) + \nabla u (q(t)) a_q^i(t) + \cancel{O(a_q^i(t)^2)},$$

por lo que simplificando el primer sumando de cada miembro de la igualdad anterior deducimos que 

$$(a_q^i)'(t) = \nabla u (q(t)) a_q^i(t)$$

Esta ecuación será fundamental para el último paso, que es estudiar cómo cambia el volumen del cubo. Llamando $Q$ al cubo,
y denotando a su volumen por $V_Q$, obtenemos\footnote{Recordando que el volumen del cubo se halla con el producto mixto.} que
$$V_Q(t) = \det(a_q^1(t), a_q^2(t), a_q^3(t)) = \begin{vmatrix}
    a_q^1(t) | a_q^2(t) | a_q^3(t) 
\end{vmatrix}.$$

Y la variación del volumen viene determinada por $V'_Q(t)$\footnote{Recordando que la derivada de un determinante se hace
por columnas (o filas dado que coincide con su traspuesta).}, que es
$$V_Q'(t) = \det((a_q^1)'(t), a_q^2(t), a_q^3(t)) + \det(a_q^1(t), (a_q^2)'(t), a_q^3(t)) + 
\det(a_q^1(t), a_q^2(t), (a_q^3)'(t))$$ definido para todo $t$, y lo vamos a evaluar en $t=0$. \\

Estudiando uno de los sumandos (por ejemplo el primero), nos percatamos de que 
$$\det((a_q^1)'(t), a_q^2(t), a_q^3(t)) = \begin{vmatrix}
    \nabla u (q) a_q^1(t) \\
    a_q^2(t) \\
    a_q^3(t)
\end{vmatrix} \stackrel{t=0}{=} \begin{vmatrix}
    \varepsilon \dfrac{\partial u_1}{\partial x_1} & \varepsilon \dfrac{\partial u_2}{\partial x_2} & \varepsilon \dfrac{\partial u_3}{\partial x_3} \\
    0 & \varepsilon & 0 \\
    0 & 0 & \varepsilon
\end{vmatrix} = \varepsilon^3 \dfrac{\partial u_1}{\partial x_1}$$
Siguiendo este esquema con los otros dos sumandos, es fácil ver que 
$$V_Q'(0) = \varepsilon^3 \dfrac{\partial u_1}{\partial x_1} + \varepsilon^3 \dfrac{\partial u_2}{\partial x_2} + \varepsilon^3 \dfrac{\partial u_3}{\partial x_3} = 
\varepsilon^3 \overbrace{\left( \dfrac{\partial u_1}{\partial x_1} + \dfrac{\partial u_2}{\partial x_2} + \dfrac{\partial u_3}{\partial x_3} \right)}^{\tr(\nabla u) = \di(u)} = 
V_Q(0) \di(u)$$

Así, vemos que la divergencia se define como sigue:
\begin{definicion}[Divergencia de un Campo Vectorial]
    Sea $u : \R^N \to \R^N$ un campo vectorial de clase $C^1(\R^N, \R^N)$, con $u = (u_1, \ldots, u_N)$. Se define la \underline{divergencia} de $u$ 
    como la aplicación dada por:
    \Func{\di(u)}{\R^N}{\R}{x}{\di(u)(x) = (\nabla \cdot u)(x) = \sum_{i=1}^{N} \dfrac{\partial u_i(x)}{\partial x_i}}
    Equivalentemente, la divergencia de $u$ coincide con la traza de su matriz jacobiana.
\end{definicion}

\subsection{Interpretación}
\label{t1:interpretación}

En pocas palabras, la variación del volumen en el instante inicial depende del propio volumen en el instante inicial, multiplicado por la
divergencia del campo de velocidades. Lo importante a recordar es que \\\\
\boxed{\text{la divergencia mide la capacidad de contracción o dilatación de un fluido con el tiempo.}}

\vspace{0.5cm}

\noindent Por lo tanto, si $\di(u) = 0$, el volumen no se comprime (ni se dilata), y se dice entonces que el fluido es \textit{incompresible}. 

\section{Rotacional}\label{section:rot}

\subsection{Motivación del concepto}

Supongamos que $\nabla u$ es una matriz antisimétrica. En tal caso se verifica que 
$$2 \nabla u = \nabla u - (\nabla u)^T = \begin{pmatrix}
    0 & - \eta_3 & \eta_2 \\
    \eta_3 & 0 & -\eta_1 \\
    -\eta_2 & \eta_1 & 0
\end{pmatrix} \Longrightarrow \begin{cases}
    \eta_1 = \dfrac{\partial u_3}{\partial x_2} - \dfrac{\partial u_2}{\partial x_3} \\\\
    \eta_2 = \dfrac{\partial u_1}{\partial x_3} - \dfrac{\partial u_3}{\partial x_1} \\\\
    \eta_3 = \dfrac{\partial u_2}{\partial x_1} - \dfrac{\partial u_1}{\partial x_2}
    \end{cases} \Longrightarrow \eta = (\eta_1, \eta_2, \eta_3) = \rot(u)$$

Podemos expresar el rotacional también como un determinante:
$$\rot(u) = \begin{vmatrix}
    i & j & k \\
    \dfrac{\partial}{\partial x_1} & \dfrac{\partial}{\partial x_2} & \dfrac{\partial}{\partial x_3} \\
    u_1 & u_2 & u_3
\end{vmatrix} = (\eta_1, \eta_2, \eta_3) = \eta$$

Ahora, realizaremos, igual que en el caso de la divergencia, un desarrollo de Taylor en torno al instante $t=0$, y asumiremos que la partícula
parte del reposo. Entonces
$$q(t) = q(0) + q'(0)t + O(t^2) = q + O(t^2)$$

Por lo tanto, las ecuaciones de las aristas vienen definidas por el siguiente sistema diferencial:
$$(a_q^i)'(t) = \nabla u (q) a_q^i(t) \quad i \in \{1,2,3\}$$
donde la matriz de coeficientes es constante. Su solución\footnote{Recordemos que si tenemos el sistema $x' = Ax$, con $A$ una
matriz de coeficientes constante, entonces la solución es $x(t) = e^{At}x_0$, donde la condición inicial es $x(t_0) = x_0$.}, utilizando la condición inicial, es:
$$a_q^i(t) = e^{\nabla u (q) t} a_q^i(0) \quad i \in \{1,2,3\}$$
En dimensión 2, el rotacional es más sencillo de calcular: para $u = (u_1, u_2)$, se tiene que
$$\rot(u) = \dfrac{\partial u_2}{\partial x_1} - \dfrac{\partial u_1}{\partial x_2}$$

% TODO: Hacerlo

\begin{ejercicio}~
    \begin{itemize}
        \item[i)] Comprobar que $e^{\nabla u(q) t}$ es una matriz de giro (propiedad de la matriz exponencial). ¿Es una matriz ortogonal? ¿Tiene determinante 1? \\
        \item[ii)] Verificar que el ángulo de giro es $\theta(t) = \frac{|\eta| \cdot t}{2}$, y el eje de giro $\eta$. 
    \end{itemize}
\end{ejercicio}

\section{Recapitulación}

Tenemos entonces $u = (u_1, u_2, u_3)$. Si $u = (u_1, u_2, u_3)$, entonces, para dimensión $N = 3$,\footnote{Donde $\times$ denota al producto vectorial.} $\rot(u) = \nabla \times u$, y para 
$N = 2$,\footnote{Donde $\cdot$ denota al producto escalar.} $u = (u_1, u_2)$, y $\rot(u) = \nabla^{\perp} \cdot u$, donde $$\nabla^{\perp} = \left( - \dfrac{\partial}{\partial x_2}, 
\dfrac{\partial}{\partial x_1} \right)$$

Necesitamos la hipótesis de antisimetría para que nos aparezca el rotacional. \\

Para el siguiente ejercicio, definimos primero el concepto de laplaciano. \\

\begin{definicion}[Laplaciano]
    Sea $u : \R^N \to \R$. Se define el laplaciano como la traza de la matriz hessiana de $u$, es decir:
    $$\Delta u = \sum_{i=1}^{N} \dfrac{\partial^2 u}{\partial x_i^2}$$
    Para un campo vectorial, $u : \R^N \to \R^N$, con $u = (u_1, \ldots, u_N)$, se define como el siguiente vector:
    $$\Delta u = \left( \sum_{i=1}^{N} \dfrac{\partial^2 u_1}{\partial x_i^2}, \ldots, \sum_{i=1}^{N} \dfrac{\partial^2 u_N}{\partial x_i^2} \right)$$
\end{definicion}

% TODO: Hacer lo que falte

\begin{ejercicio}~\label{ejercicio:ej12}
    \begin{itemize}
        \item[(i)] Comprobar que la divergencia del gradiente es el laplaciano (la traza de la matriz hessiana). Es decir, $$\di(\nabla) = \Delta$$
        Esto tiene sentido para todo $N \geq 1$. Sin embargo, comprobar que $\di(\rot) = 0$ solo tiene sentido en dimensión $N \in \{2,3\}$, porque
        el rotacional solo tiene sentido en estas dimensiones. Además, $\rot(\rot) = \nabla(\di) - \Delta$ y $\rot(\nabla^{\perp})\footnote{Rotacional del conjugado del 
        gradiente.} = \Delta$. Comprobar además las siguientes propiedades:
        $$\di(\nabla) = \Delta$$
        $$\di(\rot) = 0$$
        $$\rot(\nabla) = 0$$
        $$\rot(\rot) = \nabla(\di) - \Delta$$
        $$\rot(\nabla^{\perp}) = \Delta$$
        \item[(ii)] Dadas $f,g : \R^3 \to \R$, y $F : \R^3 \to \R^3$, se verifica lo siguiente:
        \begin{itemize}
            \item $\di(fF) = f \di(F) + F \cdot \nabla f$
            \item $\di(\nabla f \times \nabla g) = 0$
            \item $\rot(fF) = f \rot(F) + \nabla f \times F$
            \item $\Delta(fg) = f \Delta g + g \Delta f + 2 \nabla f \cdot \nabla g$ 
        \end{itemize}
        \item[(iii)] Si $G : \R^2 \to \R^2$ y $\rot(G) = 0$, entonces existe\footnote{Denominada \textit{función potencial escalar}.} $\psi : \R^2 \to \R$ tal que $\nabla \psi = G$. \\
        
        La hipótesis de que $\rot(G) = 0$ se traduce en la siguiente igualdad: 
        
        \begin{equation}\label{eq:ej21}
            \dfrac{\partial G_2}{\partial x_1} = \dfrac{\partial G_1}{\partial x_2}
        \end{equation}

        Veamos cómo se construye la función fuente $\psi$, y luego comprobaremos que cumple con lo pedido en el enunciado. \\

        Consideramos \Func{\lm}{[0,1]}{\R^2}{t}{(tx_1, tx_2)}
        Definimos $$\psi(x_1, x_2) = \int_\lm G d \lm$$
        y falta ver que $\nabla \psi = G$. En efecto:  
        $$\psi(x_1, x_2) = \int_\lm G d \lm = \int_0^1 G(\lm(t)) \cdot \lm'(t) dt = 
        \int_0^1 (G_1(\lm(t)), G_2(\lm(t))) \cdot (x_1, x_2) dt = $$
        $$\int_0^1 (G_1(tx_1, tx_2), G_2(tx_1, tx_2)) \cdot (x_1, x_2) dt = \int_0^1 \left[ x_1 G_1(tx_1, tx_2) + x_2 G_2(tx_1, tx_2) \right] dt$$

        Comprobamos que esto es correcto:
        \begin{itemize}
            \item $\dfrac{\partial \psi}{\partial x_1} =\footnote{Usamos linealidad de la derivada para meterla dentro de la integral.} \displaystyle\int_0^1 
            \left[G_1(tx_1, tx_2) + x_1 t \dfrac{\partial G_1}{\partial x_1}(tx_1, tx_2) + x_2 t \dfrac{\partial G_2}{\partial x_1}(tx_1, tx_2) \right] dt$. \\

            Ahora, usando la hipótesis \eqref{eq:ej21}~deducimos que 
            $$\dfrac{\partial \psi}{\partial x_1} = \int_{0}^{1} \dfrac{d}{dt} \left[ t G_1(tx_1, tx_2) \right] dt \stackrel{Barrow}{=} G_1(x_1, x_2)$$
            \item De forma completamente análoga deducimos que $$\dfrac{\partial \psi}{\partial x_2} = G_2$$
        \end{itemize}
        Concluimos entonces que $$\nabla \psi = \left( \dfrac{\partial \psi}{\partial x_1}, \dfrac{\partial \psi}{\partial x_2} \right) = (G_1, G_2) = G$$
        
        \item[(iv)] Si $G: \R^2 \to \R^2$ y $\di(G) = 0$,\footnote{Estaríamos tratando con el caso de un fluido incompresible.} entonces 
        existe\footnote{Denominada \textit{función fuente}.} $\psi: \R^2 \to \R$ tal que $\nabla^{\perp} \psi = G$. \\

        Debemos repetir la misma técnica que en el apartado anterior, pero definiendo 
        $$\psi = - \int_\lm G^{\perp} d \lm$$
        donde $G^\perp = (-G_2, G_1)$.
    \end{itemize}

    % Próximas semanas con modelos lineales, y al final con modelos no lineales. Evidentemente los no lineales se complican mucho.
\end{ejercicio}

    \chapter{Modelos de Transporte}
\noindent

\section{Modelo de transporte lineal}

Sean $u,\omega,f : \left[0,T \right[ \times \R^N \to \R^N$ y $\omega_0: \left[0,T \right[ \times \R^N \to \R$. Los modelos que vamos a estudiar
se corresponden con la ecuación siguiente\footnote{Puede leerse $\di_x$ como divergencia respecto de $x$.}:
\begin{equation}\label{eq:ec10}
    \dfrac{\partial u}{\partial t} + \di_x(\overbrace{\omega:u}^{\text{flujo}}) + \omega_0 u = f
\end{equation}
Al final solo nos quedará un modelo de dos términos. El primer sumando es un vector porque $u$ es un campo vectorial. El tercer sumando también,
y $f$ es vectorial. Para que la ecuación tenga sentido, el sumando de en medio debe ser un vector\footnote{En otro caso estaría mal escrita la ecuación.}. 
Así, el argumento de dicho sumando debe ser una matriz\footnote{Repasar la Sección \ref{section:div}.}. Como la ecuación es vectorial, no hay una sola ecuación, hay $N$.
Componente a componente sería
$$\dfrac{\partial u_j}{\partial t} + (\di_x(\overbrace{\omega:u}^{\text{flujo}}))_j + \omega_0 u_j = f_j \quad 1 \leq j \leq N$$
Nos interesa ahora describir el sumando $\di_x(\omega:u)$. El argumento $\omega:u$ es el producto producto tensorial o producto exterior de vectores, habitualmente denotado por 
$\omega\otimes u$. Se corresponde con la siguiente matriz $(\omega_iu_j)_{1 \leq i,j \leq N}$. Ahora nos hacemos la siguiente pregunta: ¿qué sentido tiene, si lo tiene, la divergencia de una matriz\footnote{La de una función
escalar ya la conocemos.}? La idea será aplicar la definición columna a columna de la divergencia de un vector. \\

Partimos entonces de $(\di_x(\omega : u))_j$, para cierta columna $1 \leq j \leq N$, y entonces aplicamos el concepto de divergencia ya conocido:
$$(\di_x(\omega : u))_j = \sum_{i=1}^{N} \dfrac{\partial}{\partial x_i} (w_i u_j) = \sum_{i=1}^{N} \left[\dfrac{\partial w_i}{\partial x_i} u_j + w_i \dfrac{\partial u_j}{\partial x_i} \right] =$$
$$u_j\sum_{i=1}^{N} \dfrac{\partial w_i}{\partial x_i} + \sum_{i=1}^{N} w_i \dfrac{\partial u_j}{\partial x_i} = \di_x(\omega) u_j + (\omega \cdot \nabla_x) u_j$$
Ahora vienen las simplificaciones, porque el modelo genérico es muy complicado.

\subsection{Simplificaciones del modelo lineal}

Consideraremos 

$$
\boxed{
\begin{gathered}
f = 0 \\
\omega_0 = -\di_x(\omega)
\end{gathered}
}
$$

La primera simplificación significa que no vamos a considerar fuerza. Será el modelo homogéneo. Además, es importante
tener en cuenta que $\omega$ es conocida desde el principio. Es el campo (vectorial fijo) de velocidades asociado a nuestro sistema de partículas.
Sustituyendo estas simplificaciones en el modelo general, obtenemos
$$\dfrac{\partial u}{\partial t} + \di_x(\omega) u + (\omega \cdot \nabla_x) u + \omega_0 u = f \Longrightarrow 
\dfrac{\partial u}{\partial t} + (\omega \cdot \nabla_x) u = 0$$
Asociaremos a esta ecuación la siguiente condición inicial, para cierta $u_0 : \R^N \to \R^N$ conocida:
$$u(0,x) = u_0(x)$$
de tal manera que ya podemos definir el modelo de transporte lineal (simplificado):

\begin{definicion}[Modelo de Transporte Lineal Simplificado (MTLS)]
    \begin{equation}\label{eq:ec3}
        \begin{cases}
            \dfrac{\partial u}{\partial t} + (\omega \cdot \nabla_x) u = 0 \\
            u(0,x) = u_0(x)
        \end{cases}
    \end{equation}
\end{definicion}

\begin{observacion}
    En el caso en que $N>1$, se tiene:
    $$\omega = (\omega_1, \ldots, \omega_N), \quad \nabla_x = \left( \dfrac{\partial}{\partial x_1}, \ldots, \dfrac{\partial}{\partial x_N} \right)$$
    y el producto escalar queda:
    $$(\omega \cdot \nabla_x)u = \sum_{i=1}^{N} \omega_i \dfrac{\partial u}{\partial x_i}$$
    Si $N=1$, entonces:
    \begin{equation*}
        \begin{cases}
            \dfrac{\partial u}{\partial t} + \omega(t,x) \dfrac{\partial u}{\partial x} = 0 \\
            u(0,x) = u_0(x)
        \end{cases}
    \end{equation*}
    Y se usará la notación $u_t = \frac{\partial u}{\partial t}$ y $u_x = \frac{\partial u}{\partial x}$ indistintamente.
\end{observacion}

Para resolver el PVI \eqref{eq:ec3}, debemos resolver primeramente un problema auxiliar, conocido como \textit{sistema característico}, que consiste en
encontrar ciertas trayectorias que determinarán cómo será el transporte.

\newpage

\begin{definicion}[Sistema Característico]
    Viene determinado por la ecuación $$X'(t) = \omega(t,X(t))$$
    Las trayectorias soluciones de esta ecuación nos explican cómo son las curvas que siguen nuestras partículas en $\R^N$ si
    la velocidad a la que están sujetas dichas partículas es $\omega$. Para obtener una única solución, impondremos una condición inicial de la forma
    $X(s) = x \in \R^N$. Típicamente $s=0$. Así, consideramos el siguiente PVI:
    \begin{equation}\label{eq:ec2}
        \begin{cases}
            X'(t) = \omega(t,X(t)) \\
            X(s) = x \in \R^N \quad s \in \left[0, T \right[
        \end{cases}
    \end{equation}
\end{definicion}

Debemos resolver el problema \eqref{eq:ec2}~forzosamente para resolver \eqref{eq:ec3}. Nos situaremos entonces en las siguientes condiciones favorables:
\begin{itemize}
    \item[(H1)] $\omega$ es continua en los dos argumentos. Con esto nos garantizamos la existencia.
    \item[(H2)] $\omega$ es localmente Lipschitz en el segundo argumento (l.l. resp. de $x$). Juntando ambas cosas tenemos Cauchy-Peano y Picard-Lindelöf, es decir,
    existencia y unicidad local. Habitualmente esta segunda hipótesis se deducirá de ser la función $C_x^1$.
    \item[(H3)] El crecimiento de $\omega$ es sublineal. Con esto nos garantizaremos que la solución esté definida en $\left[0, T\right[$.
    $$\norma{\omega(t,x)} \leq A|x| + B; \quad A,B \in \R$$
    \item[(H4)] La aplicación \Func{\dfrac{\partial \omega_i}{\partial x_j}}{\left[0, T \right[ \times \R^N}{\R}{(t,x)}{\dfrac{\partial \omega_i}{\partial x_j}(t,x)} es continua para todo $i,j \in \{1, \ldots, N\}$.
\end{itemize} 

Con estas 4 hipótesis, \eqref{eq:ec2} tiene una única solución. Dicha solución se denotará en lo que sigue como $X(t;s,x)$, es decir,
siempre seguiremos la pista de la condición inicial, porque será esencial. Reescribiendo \eqref{eq:ec2} usando esta última notación nos queda:
\begin{equation}\label{eq:ec32}
    \left\{
    \begin{array}{l}
    X'(t)=\omega(t,X(t))\\[6pt]
    X(s)=x
    \end{array}
    \right.
    \quad
    \overset{X(t)\equiv X(t;s,x)}{\longleftrightarrow}
    \quad
    \left\{
    \begin{array}{l}
    \dfrac{\partial X}{\partial t}(t;s,x)
    =\omega(t,X(t;s,x))\\[6pt]
    X(s;s,x)=x
    \end{array}
    \right.
\end{equation}

\begin{teo}\label{teo:t3}
    Bajo las hipótesis (H1)$-$(H4), se satisface:
    \begin{itemize}
        \item[(i)] Fórmula de traslación característica:\label{ftc}
        $$X(t;\tau, X(\tau;s,x)) = X(t;s,x) \quad \forall t, \tau, s \in \left[0, T\right[, \quad \forall x \in \R^N$$

        \begin{figure}[H]
            \begin{center}
                \tikzset{every picture/.style={line width=0.75pt}} %set default line width to 0.75pt        

                \begin{tikzpicture}[x=0.75pt,y=0.75pt,yscale=-1,xscale=1]
                %uncomment if require: \path (0,300); %set diagram left start at 0, and has height of 300

                %Shape: Axis 2D [id:dp3242275334401292] 
                \draw  (220,150) -- (360,150)(240,69) -- (240,181) (353,145) -- (360,150) -- (353,155) (235,76) -- (240,69) -- (245,76)  ;
                %Straight Lines [id:da6800909913200827] 
                \draw    (270,144) -- (270,156) ;
                %Straight Lines [id:da504363558290425] 
                \draw    (290,144) -- (290,156) ;
                %Straight Lines [id:da4011575225736249] 
                \draw    (310,144) -- (310,156) ;
                %Straight Lines [id:da4288458752818559] 
                \draw  [dash pattern={on 0.84pt off 2.51pt}]  (270,150) -- (270,119) ;
                %Straight Lines [id:da5200175588918502] 
                \draw  [dash pattern={on 0.84pt off 2.51pt}]  (269.98,119.83) -- (239.98,119.83) ;
                %Shape: Circle [id:dp34120311380094603] 
                \draw  [fill={rgb, 255:red, 0; green, 0; blue, 0 }  ,fill opacity=1 ] (268,119.5) .. controls (268,118.67) and (268.67,118) .. (269.5,118) .. controls (270.33,118) and (271,118.67) .. (271,119.5) .. controls (271,120.33) and (270.33,121) .. (269.5,121) .. controls (268.67,121) and (268,120.33) .. (268,119.5) -- cycle ;
                %Straight Lines [id:da41768238504122124] 
                \draw  [dash pattern={on 0.84pt off 2.51pt}]  (290,150) -- (290,99) ;
                %Straight Lines [id:da8451658517460895] 
                \draw    (269.5,121) -- (308,79.5) ;
                %Shape: Circle [id:dp23983039955771324] 
                \draw  [fill={rgb, 255:red, 47; green, 0; blue, 0 }  ,fill opacity=1 ] (288,100.5) .. controls (288,99.67) and (288.67,99) .. (289.5,99) .. controls (290.33,99) and (291,99.67) .. (291,100.5) .. controls (291,101.33) and (290.33,102) .. (289.5,102) .. controls (288.67,102) and (288,101.33) .. (288,100.5) -- cycle ;
                %Shape: Arc [id:dp02494682923156455] 
                \draw  [draw opacity=0] (269.98,118.83) .. controls (269.02,116.72) and (268.37,114.32) .. (268.13,111.74) .. controls (267.32,102.69) and (271.94,94.88) .. (278.46,94.3) .. controls (282.33,93.95) and (285.99,96.21) .. (288.5,100) -- (279.94,110.67) -- cycle ; \draw   (269.98,118.83) .. controls (269.02,116.72) and (268.37,114.32) .. (268.13,111.74) .. controls (267.32,102.69) and (271.94,94.88) .. (278.46,94.3) .. controls (282.33,93.95) and (285.99,96.21) .. (288.5,100) ;  
                %Curve Lines [id:da623468907736261] 
                \draw    (263,114.5) .. controls (257.09,110.56) and (246.33,96.92) .. (257.48,57.33) ;
                \draw [shift={(258,55.5)}, rotate = 106.31] [color={rgb, 255:red, 0; green, 0; blue, 0 }  ][line width=0.75]    (10.93,-3.29) .. controls (6.95,-1.4) and (3.31,-0.3) .. (0,0) .. controls (3.31,0.3) and (6.95,1.4) .. (10.93,3.29)   ;
                %Shape: Arc [id:dp26745949184759055] 
                \draw  [draw opacity=0] (289.71,98.58) .. controls (288.84,96.41) and (288.25,93.97) .. (288.01,91.37) .. controls (287.13,81.57) and (291.69,73.15) .. (298.21,72.56) .. controls (302.23,72.2) and (306.05,74.91) .. (308.6,79.34) -- (299.81,90.31) -- cycle ; \draw   (289.71,98.58) .. controls (288.84,96.41) and (288.25,93.97) .. (288.01,91.37) .. controls (287.13,81.57) and (291.69,73.15) .. (298.21,72.56) .. controls (302.23,72.2) and (306.05,74.91) .. (308.6,79.34) ;  
                %Straight Lines [id:da6666191987751137] 
                \draw  [dash pattern={on 0.84pt off 2.51pt}]  (310,150) -- (310,78) ;
                %Shape: Circle [id:dp19053748007169846] 
                \draw  [color={rgb, 255:red, 0; green, 0; blue, 0 }  ,draw opacity=1 ][fill={rgb, 255:red, 0; green, 0; blue, 0 }  ,fill opacity=1 ] (308,79.5) .. controls (308,78.67) and (308.67,78) .. (309.5,78) .. controls (310.33,78) and (311,78.67) .. (311,79.5) .. controls (311,80.33) and (310.33,81) .. (309.5,81) .. controls (308.67,81) and (308,80.33) .. (308,79.5) -- cycle ;
                %Curve Lines [id:da7070550816349991] 
                \draw    (318,76.5) .. controls (323.76,71.7) and (334.13,70.59) .. (348.22,71.39) ;
                \draw [shift={(350,71.5)}, rotate = 183.81] [color={rgb, 255:red, 0; green, 0; blue, 0 }  ][line width=0.75]    (10.93,-3.29) .. controls (6.95,-1.4) and (3.31,-0.3) .. (0,0) .. controls (3.31,0.3) and (6.95,1.4) .. (10.93,3.29)   ;

                % Text Node
                \draw (264,160) node [anchor=north west][inner sep=0.75pt]   [align=left] {$\displaystyle s$};
                % Text Node
                \draw (285,160) node [anchor=north west][inner sep=0.75pt]   [align=left] {$\displaystyle \tau $};
                % Text Node
                \draw (305,158) node [anchor=north west][inner sep=0.75pt]   [align=left] {$\displaystyle t$};
                % Text Node
                \draw (226,114) node [anchor=north west][inner sep=0.75pt]   [align=left] {$\displaystyle x$};
                % Text Node
                \draw (210,28) node [anchor=north west][inner sep=0.75pt]   [align=left] {$\displaystyle X( t;\tau ,X( \tau ;s,x))$};
                % Text Node
                \draw (354,59.5) node [anchor=north west][inner sep=0.75pt]   [align=left] {$\displaystyle X( t;s,x)$};
                % Text Node
                \draw (214,60) node [anchor=north west][inner sep=0.75pt]   [align=left] {$\displaystyle X$};


                \end{tikzpicture}
                \caption{Fórmula de traslación característica. Primero se llega hasta $\tau$, y luego hasta $t$. La fórmula nos dice que llegar de un tirón es como llegar haciendo paradas.}
                \label{fig:f21}
            \end{center}
        \end{figure}
        \item[(ii)] Fijados $s,t \in \left[0, T\right[$, la aplicación 
        \Func{T_s^t}{\R^N}{\R^N}{x}{T_s^t(x) = X(t;s,x)}
        es un difeomorfismo de clase $C^1(\R^N, \R^N)$ con operador inverso\footnote{También de clase $C^1(\R^N, \R^N)$.} 
        \Func{T_t^s}{\R^N}{\R^N}{x}{T_t^s(x) = X(s;t,x)}
        En palabras más sencillas, la inversa de ir de $s$ a $t$, es ir de $t$ a $s$ siguiendo la misma trayectoria\footnote{Podemos utilizar como ejemplo la situación de la Figura \ref{fig:f21}.}.
        \item[(iii)] El determinante jacobiano de las trayectorias características, es decir, 
        la aplicación dada por 
        \Func{J}{\left[0, T\right[^2 \times \R^N}{\R}{(t;s,x)}{J(t;s,x) = \begin{vmatrix}
            \dfrac{\partial X_i}{\partial x_j} (t;s,x) 
        \end{vmatrix}_{1 \leq i,j \leq N}}
        satisface\footnote{Podemos establecer un paralelismo entre esta expresión y la que se obtuvo de la variación del volumen del cubo en la motivación del concepto de divergencia (\ref{t1:interpretación}).} 
        $$\dfrac{\partial J}{\partial t}(t;s,x) = \di_x(\omega(t, X(t;s,x)))J(t;s,x)$$
        y además 
        $$J(t;t,x) = 1$$
        \item[(iv)] Si $\di_x(\omega) = 0$ (caso incompresible), el operador del apartado (ii), $T_s^t$, preserva la medida de Lebesgue.
    \end{itemize}
\end{teo}

\begin{proof}~
    \begin{itemize}
        \item[(i)] Por reducción al absurdo, si no fuera cierta la igualdad, habríamos encontrado dos soluciones distintas del mismo PVI \eqref{eq:ec32} 
        que coinciden en $X(\tau; s, x)$. Sustituyendo en la fórmula de traslación característica en $t = \tau$:
        $$X(\tau; \tau, X(\tau; s,x)) = X(\tau; s,x)$$
        Y tendríamos dos soluciones distintas del mismo PVI que coinciden en un punto, lo cual es una contradicción porque sabemos que este PVI
        tiene unicidad.
        \item[(ii)] Tenemos que probar que $T_s^t$ es una biyección, y luego que dicha biyección es bi-diferenciable (solo con la bicontinuidad tendríamos
        que es un homeomorfismo), y ver que ambos operadores son de clase $C^1$. \\

        Es trivial que son de clase $C^1$ con respecto al primer argumento por ser soluciones del PVI \eqref{eq:ec32}.
        Además, son de clase $C^1$ respecto a los otros dos argumentos por la hipótesis (H4), ya que dependía continua y diferenciablemente de los datos 
        iniciales. \\

        Para la inyectividad, basta con tener en cuenta que las imágenes de $T_s^t$ y $T_t^s$ son trayectorias características evolucionadas, y si 
        coinciden en algún punto, deben partir del mismo para evitar contradicción con la fórmula de traslación característica\footnote{Gráficamente lo vemos en la Figura \ref{fig:f22}.}. \\

        \begin{figure}
            \begin{center}
                \tikzset{every picture/.style={line width=0.75pt}} %set default line width to 0.75pt        

                \begin{tikzpicture}[x=0.75pt,y=0.75pt,yscale=-1,xscale=1]
                %uncomment if require: \path (0,300); %set diagram left start at 0, and has height of 300

                %Shape: Axis 2D [id:dp3242275334401292] 
                \draw  (220,150) -- (360,150)(240,69) -- (240,181) (353,145) -- (360,150) -- (353,155) (235,76) -- (240,69) -- (245,76)  ;
                %Straight Lines [id:da6800909913200827] 
                \draw    (270,144) -- (270,156) ;
                %Straight Lines [id:da504363558290425] 
                \draw    (290,144) -- (290,156) ;
                %Straight Lines [id:da4011575225736249] 
                \draw    (310,144) -- (310,156) ;
                %Straight Lines [id:da4288458752818559] 
                \draw  [dash pattern={on 0.84pt off 2.51pt}]  (270,150) -- (270,109) ;
                %Straight Lines [id:da5200175588918502] 
                \draw  [dash pattern={on 0.84pt off 2.51pt}]  (269.98,110.83) -- (239.98,110.83) ;
                %Straight Lines [id:da12444846828280065] 
                \draw  [dash pattern={on 0.84pt off 2.51pt}]  (290,150) -- (290,121) ;
                %Straight Lines [id:da2984178356461328] 
                \draw  [dash pattern={on 0.84pt off 2.51pt}]  (289.98,119.83) -- (239.98,119.83) ;
                %Shape: Circle [id:dp3001811327779722] 
                \draw  [fill={rgb, 255:red, 0; green, 0; blue, 0 }  ,fill opacity=1 ] (309,79.5) .. controls (309,78.67) and (309.67,78) .. (310.5,78) .. controls (311.33,78) and (312,78.67) .. (312,79.5) .. controls (312,80.33) and (311.33,81) .. (310.5,81) .. controls (309.67,81) and (309,80.33) .. (309,79.5) -- cycle ;
                %Straight Lines [id:da6340522878310485] 
                \draw  [dash pattern={on 0.84pt off 2.51pt}]  (310.5,150.5) -- (310.5,78.5) ;
                %Straight Lines [id:da7786265993123673] 
                \draw  [dash pattern={on 0.84pt off 2.51pt}]  (309.98,79.83) -- (239.98,79.83) ;
                %Curve Lines [id:da9873085816340168] 
                \draw    (269.5,109.5) .. controls (267,101) and (282,81) .. (309,79.5) ;
                %Curve Lines [id:da7083634369024402] 
                \draw    (290.5,119) .. controls (300,109) and (309,99) .. (310.5,78) ;

                % Text Node
                \draw (264,105) node [anchor=north west][inner sep=0.75pt]   [align=left] {x};
                % Text Node
                \draw (285,114) node [anchor=north west][inner sep=0.75pt]   [align=left] {x};


                \end{tikzpicture}
                \caption{Inyectividad. Si las dos características llegan al mismo punto, necesariamente deben partir del mismo punto. Las cruces denotan que esas características no son compatibles.}
                \label{fig:f22}
            \end{center}
        \end{figure}

        Respecto a la sobreyectividad, nos preguntamos si $\exists x \in \R^N$ tal que $y = X(t;s,x)$. Es fácil ver\footnote{Gráficamente lo vemos en la Figura \ref{fig:f23}.} que 
        dicho $x$ será $x = X(s;t, y)$.  \\

        \begin{figure}
            \begin{center}
                \tikzset{every picture/.style={line width=0.75pt}} %set default line width to 0.75pt        

                \begin{tikzpicture}[x=0.75pt,y=0.75pt,yscale=-1,xscale=1]
                %uncomment if require: \path (0,300); %set diagram left start at 0, and has height of 300

                %Shape: Axis 2D [id:dp3242275334401292] 
                \draw  (220,150) -- (360,150)(240,49) -- (240,181) (353,145) -- (360,150) -- (353,155) (235,56) -- (240,49) -- (245,56)  ;
                %Straight Lines [id:da6800909913200827] 
                \draw    (270,144) -- (270,156) ;
                %Straight Lines [id:da4011575225736249] 
                \draw    (310,144) -- (310,156) ;
                %Straight Lines [id:da4288458752818559] 
                \draw  [dash pattern={on 0.84pt off 2.51pt}]  (270,150) -- (270,119) ;
                %Straight Lines [id:da5200175588918502] 
                \draw  [dash pattern={on 0.84pt off 2.51pt}]  (269.98,119.83) -- (239.98,119.83) ;
                %Shape: Circle [id:dp34120311380094603] 
                \draw  [fill={rgb, 255:red, 0; green, 0; blue, 0 }  ,fill opacity=1 ] (268,119.5) .. controls (268,118.67) and (268.67,118) .. (269.5,118) .. controls (270.33,118) and (271,118.67) .. (271,119.5) .. controls (271,120.33) and (270.33,121) .. (269.5,121) .. controls (268.67,121) and (268,120.33) .. (268,119.5) -- cycle ;
                %Straight Lines [id:da6666191987751137] 
                \draw  [dash pattern={on 0.84pt off 2.51pt}]  (310,150) -- (310,78) ;
                %Shape: Circle [id:dp19053748007169846] 
                \draw  [color={rgb, 255:red, 0; green, 0; blue, 0 }  ,draw opacity=1 ][fill={rgb, 255:red, 0; green, 0; blue, 0 }  ,fill opacity=1 ] (308,79.5) .. controls (308,78.67) and (308.67,78) .. (309.5,78) .. controls (310.33,78) and (311,78.67) .. (311,79.5) .. controls (311,80.33) and (310.33,81) .. (309.5,81) .. controls (308.67,81) and (308,80.33) .. (308,79.5) -- cycle ;
                %Curve Lines [id:da7070550816349991] 
                \draw    (318,76.5) .. controls (323.76,71.7) and (334.13,70.59) .. (348.22,71.39) ;
                \draw [shift={(350,71.5)}, rotate = 183.81] [color={rgb, 255:red, 0; green, 0; blue, 0 }  ][line width=0.75]    (10.93,-3.29) .. controls (6.95,-1.4) and (3.31,-0.3) .. (0,0) .. controls (3.31,0.3) and (6.95,1.4) .. (10.93,3.29)   ;
                %Curve Lines [id:da691035537947874] 
                \draw    (311,79.5) .. controls (288.58,86.81) and (277.56,95.55) .. (270.07,119.15) ;
                \draw [shift={(269.5,121)}, rotate = 286.7] [color={rgb, 255:red, 0; green, 0; blue, 0 }  ][line width=0.75]    (10.93,-3.29) .. controls (6.95,-1.4) and (3.31,-0.3) .. (0,0) .. controls (3.31,0.3) and (6.95,1.4) .. (10.93,3.29)   ;
                %Straight Lines [id:da31623284526739026] 
                \draw  [dash pattern={on 0.84pt off 2.51pt}]  (308,79.5) -- (240,79.5) ;

                % Text Node
                \draw (264,160) node [anchor=north west][inner sep=0.75pt]   [align=left] {$\displaystyle s$};
                % Text Node
                \draw (305,158) node [anchor=north west][inner sep=0.75pt]   [align=left] {$\displaystyle t$};
                % Text Node
                \draw (226,114) node [anchor=north west][inner sep=0.75pt]   [align=left] {$\displaystyle x$};
                % Text Node
                \draw (354,59.5) node [anchor=north west][inner sep=0.75pt]   [align=left] {$\displaystyle X( t;s,x)$};
                % Text Node
                \draw (217,35) node [anchor=north west][inner sep=0.75pt]   [align=left] {$\displaystyle X$};
                % Text Node
                \draw (225,75) node [anchor=north west][inner sep=0.75pt]   [align=left] {$\displaystyle y$};


                \end{tikzpicture}
                \caption{Sobreyectividad. Recorremos hacia delante partiendo de $x$ en el instante $s$ y llegamos a $y$ en el instante $t$, luego $y=X(t;s,x)=T_s^t(x)$. Para volver, hacemos el recorrido hacia atrás: partimos de $y$ en el instante $t$ y retrocedemos hasta $x$ en el instante $s$,
                o sea $x=X(s;t,y)=T_t^s(y)$}
                \label{fig:f23}
            \end{center}
        \end{figure}

        Así, falta comprobar que son inversas la una de la otra, lo cual se deduce de que 
        $$T_t^s(T_s^t(x)) = T_t^s(X(t;s,x)) = X(s;t,X(t;s,x)) \stackrel{(i)}{=} X(s;s,x) = x$$
        $$T_s^t(T_t^s(x)) = T_s^t(X(s;t,x)) = X(t;s,X(s;t,x)) \stackrel{(i)}{=} X(t;t,x) = x$$

        \item[(iii)] Se deja como ejercicio.
        \item[(iv)] Por hipótesis, sustituyendo en (iii)
        $$\dfrac{\partial J}{\partial t}(t;s,x) = \di_x(\omega)J(t;s,x) = 0 \stackrel{J(t;t,x)}{\Longrightarrow} J \equiv 1$$
        Para ver que $T_s^t$ preserva la medida de Lebesgue, tomamos una función $f$ integrable y vemos que 
        $$\int_{\R^N} f(x) dx \stackrel{x = T_s^t(y)}{=} \int_{\R^N} f(T_s^t(y)) \overbrace{J}^{1} dy = \int_{\R^N} f(T_s^t(y)) dy$$
    \end{itemize}
\end{proof}

A partir de ahora, recordamos dos teorías para tratar el modelo de transporte simplificado (mtls). Recordemos que dicho modelo se corresponde con el PVI \eqref{eq:ec3}. Las dos 
variantes que tenemos son:
\begin{enumerate}
    \item Teoría fuerte / clásica, considerando $u_0 \in C^1(\R^N)$.
    \item Teoría débil, considerando $u_0 \in L_{loc}^1(\R^N)$.
\end{enumerate}

\subsection{Teoría clásica para el MTLS}

\begin{teo}[Existencia y Unicidad de Soluciones Clásicas de \eqref{eq:ec3}]\label{teo:t1}
    Bajo las hipótesis (H1)-(H4) sobre $\omega$ y dado un dato inicial $u_0 \in C^1(\R^N)$, el PVI \eqref{eq:ec3} \underline{admite una única solución}
    $u \in C^1(\left[0, T \right[ \times \R^N)$ que viene dada por 
    $$\boxed{u(t,x) = u_0(X(0;t,x))}$$
\end{teo}

Teniendo esto en cuenta, escribimos el sistema característico, resolvemos el sistema, obtendremos las trayectorias características, las recorremos a la inversa\footnote{Es decir, la evaluamos en $0$.}, y una vez tengamos ese valor, metemos esa expresión en el dato inicial.
Realmente entonces lo único que hay que hacer es encontrar las trayectorias características, es decir, resolver \eqref{eq:ec2}. \\

\begin{ejercicio}
    Resuélvanse los siguientes problemas de transporte lineal:
    \begin{enumerate}
        \item[(i)] $N=1$. $\begin{cases}
            \dfrac{\partial u}{\partial t} + x \dfrac{\partial u}{\partial x} = 0 \\
            u(0,x) = x
        \end{cases} \quad (x \in \R)$

        Tenemos que comprobar que se satisfacen las hipótesis (H1)-(H4) (sencillo en este caso), y luego resolver el sistema característico asociado al 
        campo de velocidades, es decir, \eqref{eq:ec2}. En este caso, $\omega(t,x) = x$, luego debemos resolver\footnote{Usaremos para resolver los sistemas característicos la notación $X(t) \equiv X(t;s,x)$ cuando la condición inicial de \eqref{eq:ec2} sea $X(s) = x$.}
        $$\begin{cases}
            \dfrac{\partial X}{\partial t}(t;s,x)
            =\omega(t,X(t;s,x)) = X(t;s,x)\\[6pt]
            X(s;s,x)=x
        \end{cases}$$
        y nos queda un problema unidimensional con una ecuación diferencial de variables separadas cuya solución es conocida: $X(t;s,x) = xe^{t-s}$. Como $u_0(x) = x \in C^1(\R)$, 
        por el Teorema \ref{teo:t1}, la única solución del PVI es, teniendo en cuenta que $X(0;t,x) = xe^{-t}$, la siguiente:
        $$u(t,x) = u_0(X(0;t,x)) = u_0(xe^{-t}) = xe^{-t} \in C^1(\left[0, T \right[ \times \R)$$
        \item[(ii)] $N=2$.
        $\begin{cases}
            \dfrac{\partial u}{\partial t} \overbrace{- 2y \dfrac{\partial u}{\partial x} + 2x \dfrac{\partial u}{\partial y}}^{(\omega \cdot \nabla_x) u} = 0 \\
            u(0,x,y) = x + y
        \end{cases}$

        En este caso, $u : \left[0, T \right[ \times \R^2 \to \R$ y $\omega : \left[0, T \right[ \times \R^2 \to \R^2$.

        % TODO: Terminar
    \end{enumerate}
\end{ejercicio}

Probamos ahora el Teorema \ref{teo:t1}.

\begin{proof}
    Primero demostraremos la existencia, y luego la unicidad.
    \begin{itemize}
        \item[$(i)$] Existencia. \\
        
        El dato inicial se satisface trivialmente. Ahora, veamos que se cumple \eqref{eq:ec3}. Basta con derivar y aplicar la regla de la cadena:
        $$\dfrac{\partial}{\partial t}(u_0(X(0;t,x))) = \nabla_x u_0(X(0;t,x)) \cdot \dfrac{\partial}{\partial t}(X(0;t,x))$$
        Además
        $$\left[(\omega \cdot \nabla_x)(u_0(X(0;t,x))) \right]_j = \sum_{i=1}^{N} \omega_i \dfrac{\partial}{\partial x_i} \left( u_0^j (X(0;t,x)) \right) = $$
        $$\sum_{i=1}^{N} \omega_i \sum_{k=1}^{N} \dfrac{\partial u_0^j}{\partial x_k} \left((X(0;t,x)) \right) \dfrac{\partial X_k}{\partial x_i}(0;t,x) =$$
        $$\nabla_x u_0(X(0;t,x)) \cdot ((\omega \cdot \nabla_x) X(0;t,x)) \quad 1 \leq j \leq N$$

        Falta entonces sumar ambos y ver que dicha suma es $0$. Así pues
        $$\left( \dfrac{\partial}{\partial t} + (\omega \cdot \nabla_x) \right)(u_0(X(0;t,x))) = \nabla_x u_0(X(0;t,x)) \cdot \left( \dfrac{\partial}{\partial t} + (\omega \cdot \nabla_x) \right) X(0;t,x)$$
        Para ver que esto es $0$, recurrimos a la fórmula de traslación característica (\ref{ftc}). Derivamos con respecto a $\tau$:
        $$0 = \dfrac{\partial X}{\partial \tau} (t;\tau;X(\tau;s,x)) + \overbrace{\nabla_x X}^{J \text{ (jacobiana)}} (t;\tau,X(\tau;s,x)) \cdot \dfrac{\partial X}{\partial t} (\tau; s,x)$$
        y viendo que\footnote{Por el sistema característico.} 
        $$\dfrac{\partial X}{\partial t}(\tau; s,x) = \omega(\tau; X(\tau;s,x))$$
        y evaluando en $\tau = s$:
        $$0 = \dfrac{\partial X}{\partial s} (t;s;\overbrace{X(s;s,x)}^{x}) + \overbrace{\nabla_x X}^{J \text{ (jacobiana)}} (t;s,\overbrace{X(s;s,x)}^{x}) \cdot \overbrace{\dfrac{\partial X}{\partial t} (s; s,x)}^{\omega(s,x)}$$
        Por lo tanto, evaluamos en $t=0$, inferimos que 
        $$\nabla_x u_0(X(0;t,x)) \cdot \overbrace{\left( \dfrac{\partial}{\partial t} + (\omega \cdot \nabla_x) \right) X(0;t,x)}^{0} = 0$$
        por lo que en efecto $u(t,x) = u_0(X(0;t,x))$ es solución. 

        \item[$(ii)$] Unicidad. \\
        
        Sea $u$ una solución cualquiera de \eqref{eq:ec3}. Nuevamente aplicando la regla de la cadena:
        $$\dfrac{d}{dt}(u(t,X(t;0,x))) = \dfrac{\partial u}{\partial t}(t,X(t;0,x)) + \nabla_x u (t,X(t;0,x)) \cdot \dfrac{\partial X}{\partial t}(t;0,x) = $$
        $$\dfrac{\partial u}{\partial t} (t,X(t;0,x)) + \nabla_x u(t,X(t;0,x)) \cdot \omega(t,X(t;0,x)) \stackrel{(*)}{=} 0$$
        donde en $(*)$ se ha usado que $u$ es solución de \eqref{eq:ec3}. \\

        De la última igualdad se deduce que la solución se transporta a lo largo de las trayectorias características\footnote{De ahí su nombre: se llama modelo de transporte porque el campo de velocidades $\omega$ mueve las partículas a lo largo de las trayectorias características, 
        transportando los valores iniciales de $u$ sin modificarlos.}. 
        En particular
        $$u(t;X(t;0,x)) = \overbrace{u(0,X(0;0,x))}^{u_0(x)} \quad \forall t$$
        
        Dado un dato inicial, lo que hace el operador de transporte es llevarlo por las trayectorias características. \\

        Para concluir la unicidad, hacemos el siguiente cambio de variable:
        $$u_0(X(0;t,y)) = u(t,y)\Longleftarrow y = \overbrace{X(t;0,x)}^{T_0^t(x)} \Longrightarrow x = T_t^{0}(y) = X(0;t,y)$$
        O sea, cualquier solución de \eqref{eq:ec3} tiene la forma anterior, por lo que se tiene la unicidad. \qedhere
    \end{itemize}
\end{proof}

En resumen, tomamos el dato inicial, y lo evaluamos en las trayectorias características hacia atrás. Así se soluciona el problema de transporte.

% Teoría débil mañana. Mucho más débil. Ni siquiera son continuas las soluciones, solo L^1_loc.

% Miércoles mínimo 2 horas, lo más complicado del Bloque a nivel teórico. Estudiaremos X' = \omega(t,X) porque hará falta para el bloque de fluidos
% sin las hipótesis (H1)-(H4)
% En el segundo argumento será logarítmica Lipschitz, más débil que localmente Lipschitz, y en el primero no va a ser ni continua.

\subsection{Teoría débil para el MTLS}

Consideraremos funciones $\varphi = \varphi(t,x)$ con $\varphi \in C^1_c(\left[0,T \right[ \times \R^N)$\footnote{Recordamos estos conceptos. Sea $\Omega\in\mathcal{P}(\R^N)$ y $f:\Omega\to\R$. Se define el soporte de $f$ como el cierre, relativo a $\Omega$, del conjunto de puntos en los que $f$ no se anula, es decir:
$$\sop f = \overline{\{x \in \Omega : f(x) \neq 0\}} \subset \Omega$$ Se dice que $f$ tiene soporte compacto en $\Omega$ si $\sop f$ es un subconjunto compacto de $\Omega$. En particular definimos 
$$C_c^1(\Omega) = \{f : \Omega \to \R : f \in C^1(\Omega) \quad \land \quad \sop f \text{ es compacto}\}$$ En nuestro caso consideramos una componente (temporal) más ($N+1$ en total), y $\Omega = \left[0,T\right[ \times \R^N \subset \R^{N+1}$. 
Como las funciones test $\varphi$ tienen soporte compacto en $\Omega=\left[0,T \right[\times\R^N$, existe $t_0 < T$ tal que $\varphi$ se anula en $\left]t_0, T \right[$. Por tanto, extendiéndola hacia la derecha de manera continua como $\varphi(t,x) = 0$, tenemos que $\varphi(T,x)=0$ para todo $x\in\R^N$.}. \\ 

La formulación débil será:

$$\int_0^T \int_{\R^N} \left[ \dfrac{\partial u}{\partial t} + (\omega \cdot \nabla_x) u \right] \varphi\,dx\,dt = 0$$

Entonces, hacemos lo siguiente. Por un lado
$$\int_0^T \int_{\R^N} \left[ \dfrac{\partial u}{\partial t} \right] \varphi\,dx\,dt \stackrel{Fubini | Partes}{=} 
\int_{\R^N} (u \cancelto{0}{\varphi(T,x)} - \overbrace{u \varphi(0,x))}^{u(0,x)\varphi(0,x)} \,dx - \int_{\R^N} \int_{0}^{T} u \dfrac{\partial \varphi}{\partial t} \,dt\,dx = $$
$$\boxed{- \int_{\R^N} \int_{0}^{T} u \dfrac{\partial \varphi}{\partial t} \,dt\,dx - \int_{\R^N}u_0(x) \varphi(0,x)\,dx}$$

y por el otro
$$\sum_i \int_{0}^{T} \int_{R^N} \omega_i \dfrac{\partial u_j}{\partial x_i} \varphi \,dx\,dt \stackrel{F|P}{=} 
\sum_i \int_{0}^{T} \left[ \cancelto{0}{\omega_i \varphi u_j |_{|x| \nearrow \infty}} - \int_{\R^N} u_j \dfrac{\partial}{\partial x_i} (\omega_i \varphi) \,dx\right] = $$
$$- \sum_i \int_{0}^{T} \int_{\R^N} u_j \left[ \dfrac{\partial \omega_i}{\partial x_i} \varphi + \omega_i \dfrac{\partial \varphi}{\partial x_i} \right] \,dx\,dt = $$
$$\boxed{- \int_{0}^{T} \int_{\R^N} \overbrace{\left[ \di_x(\omega) \varphi + (\omega \cdot \nabla_x) \varphi \right]}^{\di_x(\omega \varphi)} u \,dx\,dt}$$

Juntando ambas cosas, la formulación débil asociada a \eqref{eq:ec3} es \eqref{eq:ec4}:

\begin{definicion}[Formulación débil del MTLS]
    \begin{equation}\label{eq:ec4}
        - \int_{0}^{T} \int_{\R^N} \left[ \dfrac{\partial \varphi}{\partial t} + \di_x(\omega \varphi) \right] u \,dx\,dt = \int_{\R^N} u_0(x) \varphi(0,x)\,dx
    \end{equation}
\end{definicion}

Para facilitar las cosas, supondremos el caso de un fluido incompresible, es decir, $\di_x(\omega) = 0$. De esta manera\footnote{Consultar el primer punto del apartado ii) del Ejercicio \ref{ejercicio:ej12}. Aquí se ha aplicado con $f = \varphi(t, \cdot)$ y $F = \omega(t, \cdot)$
fijando el tiempo $t$, y teniendo en cuenta que la propiedad es válida para cualquier dimensión $N \geq 1$, no solo para $N = 3$.}:
$$\di_x(\omega \varphi) = \cancelto{0}{\di_x(\omega) \varphi} + (\omega \cdot \nabla_x) \varphi = (\omega \cdot \nabla_x) \varphi$$
y llegamos a la formulación con que trabajaremos:
\begin{definicion}[Formulación débil simplificada del MTLS]
    \begin{equation}\label{eq:ec5}
        - \int_{0}^{T} \int_{\R^N} \left[ \dfrac{\partial \varphi}{\partial t} + (\omega \cdot \nabla_x) \varphi \right] u \,dx\,dt = \int_{\R^N} u_0(x) \varphi(0,x)\,dx
    \end{equation}
\end{definicion}

Podemos definir entonces qué es una solución débil de \eqref{eq:ec3}:

\begin{definicion}[Solución débil de \eqref{eq:ec3}]\label{definicion:d1}
    Sean $\omega_i \in L_{loc}^{\infty}(\left[0,T \right[ \times \R^N)$, con $1 \leq i \leq N$, y $u_0 \in L_{loc}^{1}(\R^N)$.
    Diremos que $u \in L^1_{loc}(\left[0,T \right[ \times \R^N)$ es una \underline{solución débil} de \eqref{eq:ec3} si satisface \eqref{eq:ec5} para todo 
    $\varphi \in C_c^1(\left[0,T \right[ \times \R^N)$.
\end{definicion}

Veamos entonces un teorema (que asume hipótesis muy convenientes, del mismo modo que el Teorema \ref{teo:t1}) para encontrar estas soluciones débiles:

\begin{teo}[Existencia y Unicidad de Soluciones Débiles de \eqref{eq:ec3}]
    Sea $\omega : \left[0,T \right[ \times \R^N \to \R^N$ una función que satisface (H1)-(H4) y tal que $\di_x(\omega) = 0$. Sea también $u_0 \in L_{loc}^{1}(\R^N)$.
    Entonces \eqref{eq:ec3}~\underline{admite una única solución}\footnote{Por la Definición \ref{definicion:d1}, esta será de la forma $u \in L^1_{loc}(\left[0,T \right[ \times \R^N)$.} que 
    viene dada por 
    $$\boxed{u(t,x) = u_0(X(0;t,x))}$$ 
\end{teo}

\begin{observacion}
    Vemos que tanto con la formulación fuerte como con la débil, la solución de \eqref{eq:ec3}~tiene la misma forma.
\end{observacion}

\begin{proof}
    \begin{itemize}
        \item[$(i)$] Existencia. \\
        
        Nos preguntamos si 
        $$\hspace{-1cm} \overbrace{-\int_{0}^{T} \int_{\R^N} \left[ \dfrac{\partial \varphi}{\partial t} + (\omega \cdot \nabla_x) \varphi \right](t,x) u_0(X(0;t,x))\,dx\,dt}^{I}
        \stackrel{?}{=} \int_{\R^N} \varphi(0,x) u_0(x)\,dx \quad \forall \varphi \in C_c^1(\left[0,T \right[ \times \R^N)$$
        Para responder a la pregunta hacemos el siguiente cambio de variable: $y = X(0;t,x)$\footnote{Al aplicar el teorema de cambio de variable con $y=X(0;t,x)$, tenemos que 
        $dx=|J(t;0,y)|\,dy$. Por el Teorema \ref{teo:t3}, apartado (iii), como $\di_x(\omega)=0$, se deduce que $\frac{\partial J}{\partial t}(t;0,y)=0$. Además, por dicho apartado se tiene que $J(0;0,y)=1$. Por tanto, como $J$ es constante respecto de la primera
        variable $t$, $J(t;0,y) = J(0;0,y) =1$ y, en consecuencia, $dx=dy$.}, de tal manera que nos queda
        $$I = -\int_{0}^{T} \int_{\R^N} \left[ \dfrac{\partial \varphi}{\partial t} + (\omega \cdot \nabla_x) \varphi \right] (t,X(t;0,y)) u_0(y)\,dy\,dt$$
        Definimos $\psi(t,y) = \varphi(t,X(t;0,y))$, y nos queda
        $$I = -\int_{0}^{T} \int_{\R^N} \dfrac{d \psi}{dt} u_0(y) \,dy\,dt = -\int_{\R^N} (\cancelto{0}{\psi(T,y)} - \psi(0,y)) u_0(y) \,dy\,dt =$$
        $$\int_{\R^N} \psi(0,y) u_0(y) \,dy\,dt$$
        Respondemos afirmativamente a la pregunta que nos hemos hecho, por lo que queda probada la existencia.

        \item[$(ii)$] Unicidad. Por reducción al absurdo, supongamos que $u_1$ y $u_2$ son soluciones débiles. Definimos la diferencia $v = u_1 - u_2$. Entonces
        $$-\int_{0}^{T} \int_{\R^N} \left[ \dfrac{\partial \varphi}{\partial t} + (\omega \cdot \nabla_x) \varphi \right](t,x) v(t,x) \,dx\,dt = 0 \quad \forall \varphi \in C_c^1(\left[0, T \right[ \times \R^N)$$
        Podríamos pensar en aplicar el lema fundamental del cálculo de variaciones, pero $v$ no es continua, solo $L^1_{loc}$, aunque el lema fundamental también
        sirve en este caso. Para ubicarnos en esta situación, consideramos el siguiente problema:
        $$\begin{cases}
            \dfrac{\partial \varphi}{\partial t} + (\omega \cdot \nabla_x) \varphi = \psi \\
            \varphi(T,x) = 0
        \end{cases}$$
        donde $\psi$ es una función test arbitraria. El Ejercicio \ref{ejercicio:ej25} demuestra la unicidad.
        Si encontrásemos dos soluciones distintas, diferirían en un conjunto de medida nula, por lo que concluye la demostración. \qedhere
    \end{itemize}
\end{proof}

\begin{ejercicio}\label{ejercicio:ej25} % TODO: Hacer
    Sea $\varphi(t,x) = \int_{T}^{t} \psi(s, X(s;t,x)) \,ds$.
    La forma de encontrar la solución si tiene segundo miembro no negativo es integrar el segundo miembro a través de las trayectorias características.
\end{ejercicio}

\newpage

\begin{observacion}~
    \begin{itemize}
        \item Se puede integrar desde $t$ hasta $T$ y poner el signo $-$.
        \item Es obvio que se verifica la condición inicial.
    \end{itemize}
\end{observacion}

Terminada la teoría débil, tenemos el siguiente ejercicio también:

\begin{ejercicio} % TODO: Hacer
    Consideramos un fluido compresible, es decir, $\di_x(\omega) \neq 0$. Entonces se cumple que 
    $$-\int_{0}^{T} \int_{\R^N} \left[ \dfrac{\partial \varphi}{\partial t} + \di_x(\omega \varphi) \right](t,x) u(t,x) \,dx\,dt = \int_{\R^N} \varphi(0,x) u_0(x) \,dx \quad \forall \varphi \in C_c^1(\left[0, T \right[ \times \R^N)$$
\end{ejercicio}

\begin{observacion}
    La principal diferencia es que ahora, al aplicar el teorema de cambio de variable, el determinante del jacobiano $|J|$ ya no es uno. Ahora la relación
    entre la divergencia y el jacobiano es la siguiente:
    $$\dfrac{\partial J}{\partial t} = \di_x(\omega) J$$
\end{observacion}

El teorema anterior se puede demostrar para el caso de fluidos compresibles, siguiendo los mismos pasos, y considerando
$$\varphi(t,x) = \int_{T}^{t} \psi(s,X(s;t,x)) e^{\displaystyle\int_{t}^{s} \di_x(\omega (\tau; X(\tau; t,x))) \,d \tau}\,ds$$

\subsection{Preparación para el bloque de fluidos}

Hasta ahora habíamos usado que $\omega$ era continua, lo cual facilitaba bastante las cosas. Ahora $\omega$ será menos ``generoso''. \\

Consideramos el siguiente problema de trayectorias características:
\begin{equation}\label{eq:ec6}
    \begin{cases}
        X'(t) = \omega(t,X(t)) \\
        X(0) = x \in \R^N
    \end{cases}
\end{equation}

Necesitaremos una serie de conceptos, que definimos a continuación:
\begin{definicion}[Campo Vectorial Logarítmicamente Lipschitziano]
    Diremos que un campo vectorial $f : \R^N \to \R^N$ es logarítmicamente lipschitziano (abreviado en adelante por LL) en $\R^N$, denotado por $LL(\R^N)$ si 
    $\exists C > 0$ tal que 
    $$|f(x) - f(y)| \leq C |x-y|(1-\log^{*}|x-y|)$$
    donde 
    $$\log^{*}(z) = \begin{cases}
        \log(z) &\quad 0 < z \leq 1 \\
        0 &\quad z > 1
    \end{cases}$$
    Gráficamente esta última función $\log^{*}(z)$ es la siguiente:

    \begin{figure}
        \begin{center}
            \tikzset{every picture/.style={line width=0.75pt}} %set default line width to 0.75pt        

            \begin{tikzpicture}[x=0.75pt,y=0.75pt,yscale=-1,xscale=1]
            %uncomment if require: \path (0,300); %set diagram left start at 0, and has height of 300

            %Shape: Axis 2D [id:dp7966055758022831] 
            \draw  (221,150) -- (420.67,150)(239.67,60) -- (239.67,251) (413.67,145) -- (420.67,150) -- (413.67,155) (234.67,67) -- (239.67,60) -- (244.67,67)  ;
            %Straight Lines [id:da015899118604729257] 
            \draw    (289.67,139) -- (289.67,161) ;
            %Straight Lines [id:da7754339207403739] 
            \draw [color={rgb, 255:red, 243; green, 0; blue, 0 }  ,draw opacity=1 ]   (289.67,150) -- (397.67,150) ;
            \draw [shift={(399.67,150)}, rotate = 180] [color={rgb, 255:red, 243; green, 0; blue, 0 }  ,draw opacity=1 ][line width=0.75]    (10.93,-3.29) .. controls (6.95,-1.4) and (3.31,-0.3) .. (0,0) .. controls (3.31,0.3) and (6.95,1.4) .. (10.93,3.29)   ;
            %Curve Lines [id:da7496613728346398] 
            \draw [color={rgb, 255:red, 0; green, 62; blue, 255 }  ,draw opacity=1 ]   (289.67,150) .. controls (241.39,150.63) and (245.98,201.73) .. (243.22,244.51) .. controls (243.1,246.25) and (242.98,247.97) .. (242.84,249.68) ;
            \draw [shift={(242.67,251.67)}, rotate = 275.19] [color={rgb, 255:red, 0; green, 62; blue, 255 }  ,draw opacity=1 ][line width=0.75]    (10.93,-3.29) .. controls (6.95,-1.4) and (3.31,-0.3) .. (0,0) .. controls (3.31,0.3) and (6.95,1.4) .. (10.93,3.29)   ;

            % Text Node
            \draw (295,156) node [anchor=north west][inner sep=0.75pt]   [align=left] {$\displaystyle 1$};
            \end{tikzpicture}
        \end{center}
    \end{figure}

\end{definicion}

\begin{observacion}
    Para distancias pequeñas estamos corrigiendo la condición de lipschitzianidad. Para distancias grandes ($>1$), tenemos la condición de lipschitzianidad
    tal y como la conocemos hasta ahora. \\

    Ser lipschitziano es más fuerte que ser logarítmicamente lipschitziano. Es decir, se puede ser logarítmicamente lipschitziano sin ser lipschitziano.
    De esta manera Picard-Lindelöf ya no nos sirve. \\

    Definiendo \Func{\lm}{\left] 0, 1 \right]}{\R}{r}{\lm(r) = r(1-\log(r))}
    y \Func{\lm^{*}}{\R^{+}}{\R}{r}{\lm^{*}(r) = r(1-\log^{*}(r))} 
    para ver que se tiene también la continuidad, si $r \to 0 \Longrightarrow \lm(r) \to 0$. Se puede ver por L'Hôpital por ejemplo. \\
    
    Así, la condición LL implica la continuidad. \\

    Usaremos que $\lm^{*}$ es creciente en $\Rmas$ (basta derivar y ver que es creciente en dicho intervalo).
\end{observacion}

\begin{ejercicio}\label{ejercicio:ej22} % TODO: HACER
    Verificar que el crecimiento de $\lm^{*}$ es sublineal. Concretamente:
    $$\lm^{*}(r) \leq - \log(\varepsilon) r + \varepsilon, \quad 0 < \varepsilon < \dfrac{1}{e}$$
\end{ejercicio}

Vamos a ver que la única trayectoria característica que resuelve \eqref{eq:ec6} lo hace con formulación integral. Es decir, 
si tenemos una solución de \eqref{eq:ec6}, también lo será de la siguiente ecuación integral:

\begin{equation}
    X(t) = x + \int_{0}^{t} \omega(s,X(s))\,ds
\end{equation}

Bastaría cambiar $L^1_{loc}$ por continuidad para tener existencia y unicidad de \eqref{eq:ec6}, donde la solución sería $C^1$. Más adelante 
veremos que cuando lleguemos al bloque de fluidos y a la ecuación de Euler, la condición será $L^1_{loc}$. \\

Resumimos lo anterior en el siguiente teorema:

\begin{teo}
    Sea $\omega : \left[ 0, T \right[ \times \R^N \to \R^N$ un campo vectorial localmente integrable en la primera variable, y LL en la segunda variable.
    Entonces \underline{existe una única solución continua}, $X \in C(\left[0, T \right[)$ del PVI \eqref{eq:ec6}. 
\end{teo}


\begin{proof}
    La demostración es una especie de reconstrucción del teorema de Picard. Construiremos las iterantes de Picard, veremos que están bien definidas
    y generan una sucesión que pasa bien al límite, y el límite es el que buscamos. \\

    \begin{itemize}
        \item[(i)] Existencia. \\
        
        Definimos recursivamente
        \begin{equation}\label{eq:ec7}
            \begin{cases}
                X_{n+1}(t) = x + \displaystyle\int_{0}^{t} \omega(s,X_n(s))\,ds \quad n \geq 0 \\
                X_0 \equiv x
            \end{cases}
        \end{equation}
        Ahora veamos que está bien definido \eqref{eq:ec7}. El caso $n=0$ está trivialmente bien definido. Para $n=1$:
        $$X_1 = x + \displaystyle\int_{0}^{t} \omega(s,X_0(s))\,ds = x + \displaystyle\int_{0}^{t} \omega(s,x)\,ds$$
        Sea ahora $I=\left[0,T\right[$. Por hipótesis, para cada $x\in\mathbb{R}^N$
        fijo, se tiene que
        \[
        \omega(\cdot,x)\in L^1_{\mathrm{loc}}(I,\mathbb{R}^N).
        \]
        Fijado $0<T_0<T$ arbitrario, como $[0,T_0]$ es un
        compacto contenido en $I$, resulta que
        \[
        \int_0^{T_0}|\omega(s,x)|\,ds<\infty.
        \]

        Por tanto, para todo $t\in[0,T_0]$, aplicando la
        desigualdad triangular:
        \[
        |X_1(t)|\leq |x|+\int_0^t|\omega(s,x)|\,ds
        \leq |x|+\int_0^{T_0}|\omega(s,x)|\,ds<\infty.
        \]

        Además, $X_1$ es continua, pues la integral indefinida
        de una función localmente integrable es absolutamente
        continua en cada intervalo compacto\footnote{Visto en Análisis Matemático II.}. Como $T_0<T$ es
        arbitrario, y la composición de una función constante con una función absolutamente continua es absolutamente continua (luego continua), concluimos que $X_1\in C(I,\mathbb{R}^N)$. \\

        Ahora argumentamos inductivamente. Supuesto que $X_n(t)$ está bien definida y es continua, probamos esta misma propiedad para $X_{n+1}(t)$.
        Como
        $$\left|\int_{0}^{t} \omega(s,X_n(s))\,ds \right| \leq \int_{0}^{T_0} |\omega(s,X_n(s))|\,ds$$
        y ahora el ``truco'' está en restar y sumar dentro del argumento del valor absoluto la cantidad $\omega(s,x)$, es decir:
        $$\int_{0}^{T_0} |\omega(s,X_n(s))|\,ds \leq \int_{0}^{T_0} |\omega(s,X_n(s)) - \omega(s,x)|\,ds + \int_{0}^{T_0} |\omega(s,x)|\,ds
        \stackrel{LL}{\leq}$$
        $$C\int_{0}^{T_0} \lm^{*}(|X_n(s) - x|)\,ds + \overbrace{\int_{0}^{T_0} |\omega(s,x)|\,ds}^{<\infty \text{ pq } \omega \in L^1_{loc}}$$
        Ahora como $\lm^{*}(|X_n(s) - x|)$ es continua (composición de continuas) en un compacto ($[0,T_0]$), por el Teorema de Weierstrass 
        tiene máximo, pongamos $M = \max_{0 \leq s \leq T_0} |X_n(s) - x|$. Análogamente, como la segunda integral es finita, pongamos
        que su valor es $M_{T_0}$, y nos queda que 
        $$C\int_{0}^{T_0} \lm^{*}(|X_n(s) - x|)\,ds + \int_{0}^{T_0} |\omega(s,x)|\,ds \leq C \int_{0}^{T_0} \lm^{*}(M)\,ds + M_{T_0} \leq$$
        $$C \lm^{*}(M)T_0 + M_{T_0} < \infty$$
        Ahora usaremos el criterio de la mayorante de Weierstrass (para series convergentes). La idea consiste en estudiar:
        $$|X_{n+1}(t) - X_n(t)| = \left| \int_{0}^{t} \left[ \omega(s,X_n(s)) - \omega(s,X_{n-1}) \right] \,ds \right| \leq$$
        $$\int_{0}^{t} |\omega(s,X_n(s)) - \omega(s,X_{n-1}(s))|\,ds \stackrel{LL}{\leq} C \int_{0}^{t} \lm^{*}(|X_n(s) - X_{n-1}(s)|)\,ds$$
        Por el Ejercicio \ref{ejercicio:ej22}, se tiene que 
        $$\lm^{*}(|X_n(s) - X_{n-1}(s)|) \leq a_{\varepsilon} |X_n(s) - X_{n-1}(s)| + \varepsilon, \quad a_{\varepsilon} = -\log(\varepsilon)$$
        y juntándolo con la última acotación superior:
        $$C \int_{0}^{t} \lm^{*}(|X_n(s) - X_{n-1}(s)|)\,ds \leq C \varepsilon t + C a_{\varepsilon} \int_{0}^{t} |X_n(s) - X_{n-1}(s)|\,ds$$
        A su vez con el último integrando podemos hacer una acotación análoga:
        $$|X_n(s) - X_{n-1}(s)| \leq C \varepsilon s + C a_{\varepsilon} \int_{0}^{s} |X_{n-1}(\tau) - X_{n-2}(\tau)|\,d \tau$$
        y repetiendo este proceso hasta llegar a la primera diferencia ($X_1 - X_0$), obtenemos
        $$C \varepsilon s + C a_{\varepsilon} \int_{0}^{s} |X_{n-1}(\tau) - X_{n-2}(\tau)|\,d \tau \leq $$
        $$C \varepsilon t + C a_{\varepsilon} \left[ \dfrac{C \varepsilon t^2}{2} + Ca_{\varepsilon} \int_{0}^{t} \int_{0}^{s} 
        \overbrace{|X_{n-1}(\tau) - X_{n-2}(\tau)|}^{C \varepsilon \tau + C a_\varepsilon \int_{0}^{\tau} |X_{n-2}(\sigma) - X_{n-3}(\sigma)|\,d \sigma}\,d \tau \, ds  \right] =$$
        $$C \varepsilon t + C^2 \varepsilon a_\varepsilon \dfrac{t^2}{2} + C^2 a^2_\varepsilon
        \left[ C \varepsilon \dfrac{t^3}{6} + C a_\varepsilon \int_{0}^{t} \int_{0}^{s} \int_{0}^{\tau} |X_{n-2}(\sigma) - X_{n-3}(\sigma)| d \sigma \right] = $$
        $$C \varepsilon t + C^2 \varepsilon a_\varepsilon \dfrac{t^2}{2} + C^3 \varepsilon a^2_\varepsilon \dfrac{t^3}{6} + C^3 a_\varepsilon^3 \int_{0}^{t} \int_{0}^{s} \int_{0}^{\tau}
        |X_{n-2}(\sigma) - X_{n-3}(\sigma)|\,d \sigma$$
        y en general esta expresión se escribe como sigue:
        $$\varepsilon \sum_{k=0}^{n-1} \dfrac{c^{k+1} t^{k+1} a_\varepsilon^k}{(k+1)!} + C^n a^n_\varepsilon \int_{0}^{t} \overbrace{\cdots}^{n)} \int_{0}^{s} |X_1(s) - \overbrace{X_0(s)}^{x}|\,ds$$
        y esta primera integral a su vez se escribe como
        $$|X_1(s) - x| = \int_{0}^{s} \omega(\tau, x)\,d \tau \leq \int_{0}^{T_0} \omega(\tau, x)\, d \tau \leq M_{T_0}$$
        Así pues, recapitulando, tenemos que 
        $$C \varepsilon s + C a_{\varepsilon} \int_{0}^{s} |X_{n-1}(\tau) - X_{n-2}(\tau)|\,d \tau \leq $$
        $$\varepsilon \sum_{k=0}^{n-1} \dfrac{c^{k+1} t^{k+1} a_\varepsilon^k}{(k+1)!} + C^n a^n_\varepsilon \int_{0}^{t} \overbrace{\cdots}^{n)} \int_{0}^{s} |X_1(s) - \overbrace{X_0(s)}^{x}|\,ds <$$
        $$\varepsilon C t \sum_{k=0}^{n-1} \dfrac{(c t a_\varepsilon)^k}{k!} + C^n a_\varepsilon^n M_{T_0} \dfrac{t^n}{n!} \leq$$
        $$\varepsilon C T_0 \sum_{k=0}^{n-1} \dfrac{(c T_0 a_\varepsilon)^k}{k!} + C^n a_\varepsilon^n M_{T_0} \dfrac{T_0^n}{n!}$$
        En resumen, tenemos que 
        \begin{equation}\label{eq:ec8}
            |X_{n+1}(t) - X_n(t)| \leq \varepsilon C T_0 \sum_{k=0}^{n-1} \dfrac{(c T_0 a_\varepsilon)^k}{k!} + 
            C^n a_\varepsilon^n M_{T_0} \dfrac{T_0^n}{n!} \quad 0 < \varepsilon \leq \dfrac{1}{e}
        \end{equation}
        donde $a_\varepsilon = -\log(\varepsilon)$. Ahora, tomamos $\varepsilon = e^{-n}$ y $T_0 < \frac{1}{2Ce}$, y sustituyendo estos valores en el segundo miembro de \eqref{eq:ec8} nos queda\footnote{Usando la monotonía de la 
        integral: $$M_{T_0} = \int_{0}^{T_0} |\omega(s,x)|\,ds < \int_{0}^{\frac{1}{2Ce}} |\omega(s,x)|\,ds \stackrel{not.}{=}M_{\frac{1}{2Ce}}$$}:
        $$|X_{n+1}(t) - X_n(t)| < \dfrac{e^{-n}}{2e} \sum_{k=0}^{n} \dfrac{n^k}{(2e)^k k!} + \dfrac{n^n}{(2e)^n n!} M_{\frac{1}{2Ce}}$$
        Como $1/(2e) < 1$, $\sum_{k=0}^{n} \frac{n^k}{(2e)^k k!} < e^{\frac{n}{2e}}$ tenemos que 
        $$\dfrac{e^{-n}}{2e} \sum_{k=0}^{n} \dfrac{n^k}{(2e)^k k!} + \dfrac{n^n}{(2e)^n n!} M_{\frac{1}{2Ce}} < \overbrace{e^{-n\left(1 - \dfrac{1}{2e} \right)}}^{a_n} + 
        \overbrace{\dfrac{n^n}{(2e)^n n!} M_{\frac{1}{2Ce}}}^{b_n}$$
        Por el Ejercicio \ref{ejercicio:ej23}, la serie de término general $a_n + b_n$ es convergente, y por el criterio de la mayorante de 
        Weierstrass, llegamos a que
        $$X_0(t) + \sum_{k=0}^{\infty} (X_{k+1}(t) - X_k(t)) \quad \text{c.u.} \Longrightarrow$$
        $$X_n(t) = X_0(t) + \sum_{k=0}^{n-1} (X_{k+1}(t) - X_k(t)) \quad \text{c.u.}$$ 
        En este punto tenemos que 
        $$X_{n+1}(t) = x + \int_{0}^{t} \omega(s,X_n(s))\,ds$$
        Tomando límites en cada miembro
        $$X_{n+1}(t) \stackrel{c.u.}{\to} X(t) \text{ continua}$$
        $$x + \int_{0}^{t} \omega(s,X_n(s))\,ds \to x + \lim_{n \nearrow \infty} \int_{0}^{t} \omega(s,X_n(s))\,ds \stackrel{(*)}{=} x + \int_{0}^{t} \omega(s,X(s))\,ds$$
        en $(*)$ se ha usado el Ejercicio \ref{ejercicio:ej24}. \\

        Falta ver que $X$ es solución de \eqref{eq:ec6}. Esto se deduce de que hemos obtenido que:
        $$X(t)=x+\int_0^t\omega(s,X(s))\,ds$$
        En particular, $X(0)=x$ y, por el Teorema Fundamental del Cálculo, $X'(t)=\omega(t,X(t))$ para casi todo
        $t\in[0,T[$. Concluimos así que $X$ es solución de \eqref{eq:ec6}.

        \item[(ii)] Unicidad. Supongamos que $X_1$ y $X_2$ son soluciones distintas de $$X(t) = x + \int_{0}^{t} \omega(s,X(s))\,ds$$ Sea
        $$Y(t) := |X_1(t) - X_2(t)|$$
        queremos probar que $Y(t) = 0 \quad \forall t$. Vemos primero que 
        $$Y(0) = X_1(0) - X_2(0) = 0$$
        ya que, al ser $X_1$ y $X_2$ soluciones, verifican la misma condición inicial. También notamos que $Y$ es continua como composición de funciones continuas.
        Así $Y(t) < 1$ en cierto intervalo $[0, T^{*}[$ por continuidad\footnote{Podría escribirse $T^*=\sup\{t\in[0,T]:Y(s)<1\text{ para todo }s\in[0,t]\}$.}. Además $Y(T^{*}) = 1$. Tenemos
        entonces que (trabajamos en $[0,T^{*}[)$.
        $$Y(t) \leq \int_{0}^{t} |\omega(s,X_1(s)) - \omega(s,X_2(s))|\,ds \stackrel{LL}{\leq} C \int_{0}^{t} \lm^{*}(Y(s))\,ds = $$
        $$C \int_{0}^{t} \lm(Y(s))\,ds$$
        Recordando que $\lm(r) = r(1-\log(r))$ es creciente en $\Rmas$, en particular lo es en $]0,1[$ y $0 < \lm(x) < 1$. \\

        Definimos $\lm_a(t) := C \int_{0}^{t} \lm(Y(s))\,ds + a$ y mayoramos:
        $$C \int_{0}^{t} \overbrace{\lm(Y(s))}^{<1}\,ds + a \leq C T_0 + a < 1 \quad [0,T_0], \quad T_0 < T^{*} \leq T$$
        Donde en la última desigualdad $(<1)$ se han tomado (dado que podemos hacerlo) las siguientes restricciones: $0 < a < 1/2$ y $T_0 < 1/(2C)$.
        Por último, obtenemos 
        $$\int_{a}^{\lm_a(t)} \dfrac{dr}{\lm(r)} \stackrel{r = \lm_a(s)}{=} \int_{0}^{t} \dfrac{\lm_a'(s)}{\lm(\lm_a(s))}\,ds = 
        C\int_{0}^{t} \overbrace{\dfrac{\lm(Y(s))}{\lm(\lm_a(s))}}^{<1}\,ds < Ct$$
        Hacia la izquierda, descomponemos la integral:
        $$
        \int_a^{\lm_a(t)}\frac{dr}{\lm(r)}
        =
        \int_a^1\frac{dr}{\lm(r)}
        -
        \int_{\lm_a(t)}^1\frac{dr}{\lm(r)}.
        $$
        Por tanto, como $Y(t)\leq \lm_a(t) - a <\lm_a(t)$ ya que $a > 0$, obtenemos
        $$
        \begin{aligned}
        \int_a^1\frac{dr}{\lm(r)}
        &<Ct+\int_{\lm_a(t)}^1\frac{dr}{\lm(r)}\\
        &\leq Ct+\int_{Y(t)}^1\frac{dr}{\lm(r)}\\
        &=Ct+\log(1-\log(Y(t))),
        \end{aligned}
        $$
        siempre que $Y(t)>0$ y $t\in[0,T_0]$. \\

        Llegamos ahora a una contradicción. En efecto,
        supongamos que existe $t_*\in[0,T_0]$ tal que
        $Y(t_*)>0$. Entonces,
        \[
        \int_a^1\frac{dr}{\lm(r)}
        =\log(1-\log(a))
        \xrightarrow[a\to0^+]{}+\infty,
        \]
        mientras que el segundo miembro de la desigualdad
        anterior permanece finito, lo cual es una contradicción.
        Por tanto, necesariamente
        \[
        Y(t)=0\qquad\forall t\in[0,T_0].
        \]

        Finalmente, como $X_1(T_0)=X_2(T_0)$, podemos
        repetir el mismo argumento tomando $T_0$ como
        nuevo instante inicial. Procediendo sucesivamente
        en intervalos suficientemente pequeños, obtenemos
        que $Y(t)=0$ para todo $t\in[0,T^*[$.

        Por continuidad, $T^*$ no puede ser un instante
        anterior a $T$ en el que $Y$ alcance el valor $1$.
        En consecuencia, $T^*=T$ y concluimos que
        \[
        X_1(t)=X_2(t)\qquad\forall t\in[0,T[,
        \]
        lo que demuestra la unicidad. \qedhere
    \end{itemize}
\end{proof}

\begin{ejercicio}\label{ejercicio:ej23} % TODO: Hacer
    Estudiar, por el criterio que se prefiera (recomendable el criterio del cociente), la convergencia de las series de término general $a_n$
    y $b_n$. \\

    Haciendo el criterio del cociente sale que ambas series convergen. En particular la serie de término general $a_n + b_n$ converge.
\end{ejercicio}

\begin{ejercicio}\label{ejercicio:ej24} % TODO: Hacer
    Demostrar que 
    $$\lim_{n \nearrow \infty} \int_{0}^{t} \omega(s,X_n(s))\,ds = \int_{0}^{t} \omega(s,X(s))\,ds$$
    Se justifica con que $\omega$ es LL en el segundo argumento, en particular continua.
\end{ejercicio}

% En dos semanas empiezan fluidos. 6.10.2026

\section{Modelo de transporte no lineal: ecuación de Burgers}

% Lo siguiente importante para Parcial 1

Esta es la ecuación más simple de transporte no lineal y ya es complicada. Para simplificar la notación vamos a restringir la variable espacial
al caso unidimensional (la temporal ya era unidimensional). Los gradientes serán simplemente derivadas. 

\begin{equation}\label{eq:ec9}
    \begin{cases}
        \dfrac{\partial u}{\partial t} + u \dfrac{\partial u}{\partial x} = 0 \\
        u(0,x) = u_0(x)
    \end{cases}
\end{equation}

\begin{observacion}
    Ahora $\omega$ va a depender de $u$, y ahí es donde se genera la no linealidad. Para ver que \eqref{eq:ec9} es no lineal, vemos que se puede reescribir como
    \begin{equation}
        \dfrac{\partial u}{\partial t} + \dfrac{\partial}{\partial x} \left( \overbrace{\dfrac{u^2}{2}}^{\text{flujo}} \right) = 0
    \end{equation}
    Comparamos el flujo de $\omega : u$ de \eqref{eq:ec10} con $u^2/2$, que ahora es un flujo cuadrático. \\

    Recordemos también que 
    $$\dfrac{d}{dt}(u(t,X(t;0,x))) = \dfrac{\partial u}{\partial t}(t,X(t;0,x)) + \dfrac{\partial u}{\partial x}(t,X(t;0,x))X'(t)$$
    Y vimos además que cuando la ecuación de transporte era lineal $X'(t) = \omega$, en este caso es el coeficiente que aparece multiplicando la derivada de posición, 
    es decir, en $\eqref{eq:ec9}$ se correspondería con 
    \begin{equation}\label{eq:ec11}
        X'(t) = u(t,X(t)) 
    \end{equation}
    Por tanto, nos queda
    $$\dfrac{\partial u}{\partial t}(t,X(t;0,x)) + \dfrac{\partial u}{\partial x}(t,X(t;0,x))\overbrace{X'(t)}^{u(t,X)} = 0$$
    Cualquier solución de nuestro problema evaluada a lo largo de las trayectorias características no depende del tiempo. Así, su valor en un instante cualquiera
    coincide con su valor en $0$:
    \begin{equation}\label{eq:ec12}
        u(t,X(t;0,x)) = u(0,X(0;0,x)) = u_0(x) \quad \forall t
    \end{equation}
    Retomando \eqref{eq:ec10} y sustituyendo, podemos llegar a la formulación integral ya conocida:
    $$X(t) - x = \int_0^t u_0(x)\,ds = u_0(x) t$$
    por lo que en conclusión
    \begin{equation}\label{eq:ec17}
        X(t) = x + u_0(x)t
    \end{equation}
    es decir: \textbf{las soluciones de la ecuación de Burgers (las trayectorias características) son rectas, donde la pendiente la determina la condición inicial $u_0(x)$}. \\

    Veremos casos de un modelo de Burgers con infinitas soluciones. No conseguimos retener \textbf{nada} de la teoría clásica. Ni se mantiene la regularidad, ni la continuidad, ni la
    unicidad. Vamos a hacer lo poco que se puede hacer.
\end{observacion}

\begin{ejemplo}
    Consideramos el siguiente modelo de Burgers:
    \begin{equation*}
        \begin{cases}
            \dfrac{\partial u}{\partial t} + u \dfrac{\partial u}{\partial x} = 0 \\
            u(0,x) = \dfrac{1}{1+x^2}
        \end{cases}
    \end{equation*}
    Las características asociadas a ese modelo ya sabemos que son líneas rectas, y que vienen determinadas por la condición inicial. Tenemos de momento que
    $$X(t;0,x) = x + \dfrac{t}{1+x^2}$$
    Pongamos distintos valores:
    $$x=0 \Longrightarrow X_1(t;0,0) = t$$
    $$x=1 \Longrightarrow X_2(t;0,1) = 1 + \dfrac{t}{2}$$
    Gráficamente, tenemos las siguientes dos rectas:
    \begin{figure}[H]
        \begin{center}
            \tikzset{every picture/.style={line width=0.75pt}} %set default line width to 0.75pt        

            \begin{tikzpicture}[x=0.75pt,y=0.75pt,yscale=-1,xscale=1]
            %uncomment if require: \path (0,300); %set diagram left start at 0, and has height of 300

            %Shape: Axis 2D [id:dp3242275334401292] 
            \draw  (230.67,150) -- (360,150)(240.67,38) -- (240.67,181) (353,145) -- (360,150) -- (353,155) (235.67,45) -- (240.67,38) -- (245.67,45)  ;
            %Straight Lines [id:da6800909913200827] 
            \draw    (270,144) -- (270,156) ;
            %Straight Lines [id:da504363558290425] 
            \draw    (300,144) -- (300,156) ;
            %Straight Lines [id:da41768238504122124] 
            \draw  [dash pattern={on 0.84pt off 2.51pt}]  (300.5,150.5) -- (300.5,89.5) ;
            %Shape: Circle [id:dp23983039955771324] 
            \draw  [fill={rgb, 255:red, 47; green, 0; blue, 0 }  ,fill opacity=1 ] (299,89.5) .. controls (299,88.67) and (299.67,88) .. (300.5,88) .. controls (301.33,88) and (302,88.67) .. (302,89.5) .. controls (302,90.33) and (301.33,91) .. (300.5,91) .. controls (299.67,91) and (299,90.33) .. (299,89.5) -- cycle ;
            %Straight Lines [id:da8200181075779964] 
            \draw    (240,150) -- (330.67,60) ;
            %Straight Lines [id:da5662793742139913] 
            \draw    (230,120) -- (250,120) ;
            %Straight Lines [id:da2599786825287116] 
            \draw    (239.67,120) -- (363.67,58) ;
            %Straight Lines [id:da18823594983513814] 
            \draw    (230,90) -- (250,90) ;

            % Text Node
            \draw (264,158) node [anchor=north west][inner sep=0.75pt]   [align=left] {$\displaystyle 1$};
            % Text Node
            \draw (219,24) node [anchor=north west][inner sep=0.75pt]   [align=left] {$\displaystyle x$};
            % Text Node
            \draw (301,43) node [anchor=north west][inner sep=0.75pt]   [align=left] {$\displaystyle X_{1}$};
            % Text Node
            \draw (215,104) node [anchor=north west][inner sep=0.75pt]   [align=left] {$\displaystyle 1$};
            % Text Node
            \draw (374,49) node [anchor=north west][inner sep=0.75pt]   [align=left] {$\displaystyle X_{2}$};
            % Text Node
            \draw (295,159) node [anchor=north west][inner sep=0.75pt]   [align=left] {$\displaystyle 2$};
            % Text Node
            \draw (367,150) node [anchor=north west][inner sep=0.75pt]   [align=left] {$\displaystyle t$};
            % Text Node
            \draw (215,74) node [anchor=north west][inner sep=0.75pt]   [align=left] {$\displaystyle 2$};

            \end{tikzpicture}
        \end{center}
    \end{figure}
    Basta igualar y despejar el tiempo para ver en qué instante chocan. En este caso es claro, pues $t = 1 + t/2 \iff t/2 = 1 \iff t = 2$.
    Esta colisión podría derivar en un absurdo a priori:
    $$1 = u_0(0) \stackrel{(*)}{=} u(2,\overbrace{X_1(2;0,0)}^{X_2(2;0,1)}) = u(2,X_2(2;0,1)) = u_0(1) = 1/2$$
    donde se ha usado en $(*)$ que $X_1$ es una característica y \eqref{eq:ec12}. Todo lo anterior es correcto. La única explicación posible es que 
    la colisión entre características se traduce en la discontinuidad de la solución. Donde las características chocan, la solución salta. Realmente 
    estamos evaluando en un punto donde no podemos evaluar, porque los límites laterales no coinciden, a pesar de que el dato inicial $u(0,x) = 1/(1+x^2) \in C^\infty(\R)$.
\end{ejemplo}

\begin{ejercicio}\label{ejercicio:ej26}
    Fijado en \eqref{eq:ec9} un dato inicial $u_0 \in C^1(\R)$ (marco de la teoría clásica), \underline{decreciente}\footnote{Esta será la condición fundamental para que 
    las características colisionen.} y con derivada inferiormente acotada\footnote{Cuestión técnica.}, comprueba\footnote{Se trata de escribir cuándo se tienen garantías de 
    que las características choquen (en tal caso tendremos garantías también de discontinuidad de la solución). Caracterizar cuál es el primer tiempo de choque de dos características 
    cualesquiera. A partir de ahí se produce el salto de la solución y se pierde la continuidad.} que el primer instante en el que dos rectas características 
    cualesquiera colisionan, denotado por $T^*$, viene dado por $$T^{*} = - \dfrac{1}{\inf_{x \in \R} \{u_0'(x)\}}$$

    Para hacer el ejercicio, lo más sencillo es verlo como una aplicación inmediata del Teorema del Valor Medio. Así sale en dos líneas. Tomar dos características cualquiera,
    que sabemos que son rectas, chocamos (vemos en qué instantes chocan), tomamos mínimos, y ahí se aplica el Teorema del Valor Medio.

    % TODO: Hacer
\end{ejercicio}

De esta manera, la solución de Burgers definida en $u : \left[0, T^{*} \right[ \times \R \to \R$ la solución sería $C^1$, y se podría aplicar la teoría clásica. Aún así,
$T^{*}$ puede ser muy pequeño. \\

% Lo que viene es ejercicio de Parcial 1

Visto todo lo anterior, lo natural parece ser trabajar con soluciones débiles, dado que no nos exigen ni continuidad. Lo que resta de estudio de la ecuación de Burgers, que es poco,
porque no se puede hacer mucho más sin entrar en cosas demasiado complicadas, es estudiar cuáles son las soluciones débiles más sencillas del modelo de Burgers. Estas serán de la forma:
\begin{equation}\label{eq:ec13}
    u(t,x) = \begin{cases}
        u^{-}, \quad& x < \lm t \\
        u^{+}, \quad& x > \lm t
    \end{cases}
\end{equation}

Gráficamente\footnote{Se ha representado con cierto $\lm > 0$.}:

\begin{figure}[H]
    \begin{center}
        \tikzset{every picture/.style={line width=0.75pt}} %set default line width to 0.75pt        

        \begin{tikzpicture}[x=0.75pt,y=0.75pt,yscale=-1,xscale=1]
        %uncomment if require: \path (0,300); %set diagram left start at 0, and has height of 300

        %Shape: Axis 2D [id:dp3242275334401292] 
        \draw  (230.67,151) -- (360,151)(239.67,38) -- (239.67,166) (353,146) -- (360,151) -- (353,156) (234.67,45) -- (239.67,38) -- (244.67,45)  ;
        %Straight Lines [id:da8200181075779964] 
        \draw    (240,151) -- (351.67,80) ;

        % Text Node
        \draw (219,24) node [anchor=north west][inner sep=0.75pt]   [align=left] {$\displaystyle x$};
        % Text Node
        \draw (363,62) node [anchor=north west][inner sep=0.75pt]   [align=left] {$\displaystyle x\ =\ \lambda t$};
        % Text Node
        \draw (367,150) node [anchor=north west][inner sep=0.75pt]   [align=left] {$\displaystyle t$};
        % Text Node
        \draw (317,114) node [anchor=north west][inner sep=0.75pt]   [align=left] {$\displaystyle u^{-}$};
        % Text Node
        \draw (286,77) node [anchor=north west][inner sep=0.75pt]   [align=left] {$\displaystyle u^{+}$};


        \end{tikzpicture}
    \end{center}
\end{figure}

Las soluciones tendrán lado constante a la izquierda de $x = \lm t$ y a la derecha. Para trabajar con estas soluciones, el dato inicial del que hay que partir debe ser 

\begin{equation}\label{eq:ec14}
    u_0(x) = \begin{cases}
        u^{-}, \quad& x < 0 \\
        u^{+}, \quad& x > 0
    \end{cases}
\end{equation}


Cuando a la ecuación de Burgers se le asocia un dato inicial con salto ($u^{-}$ por arriba y $u^{+}$ por abajo). Este tipo de dato inicial da pie al siguiente tipo más específico de problema, conocido como \underline{problema de Riemann}:
\begin{definicion}[Problema de Riemann]\label{definicion:d21}
    El problema de Burgers \eqref{eq:ec9} con un dato inicial de la forma \eqref{eq:ec14} se denomina problema de Riemann. \\
\end{definicion}

%\begin{observacion}
%    Lo que hace el modelo de Burgers es ``inclinar'' la recta $x=0$ a $x=\lm t$. \\
%\end{observacion}
% El único valor posible de $\lm = \dfrac{1}{2}(u^{+} + u^{-})$ para que $u$ sea solución débil.
% Veremos el martes bajo qué condiciones esto es un problema de Riemann
% Quizá sugiera que el que entregue el ejercicio propuesto lo cuente como participación.
% Seguramente usaremos el lunes 19 para terminar el bloque, y recomienda llevarse la prueba a la semana siguiente o más. Lo mejor sería retrasarla al menos una semana dice. 
% 13.10.2026
Retomamos el problema \eqref{eq:ec9}. Una solución débil de este problema era \eqref{eq:ec13}, donde $u^{-}, u^{+}, \lm \in \R$, para cierto $\lm$, y el dato inicial era \eqref{eq:ec14}.
Así, el modelo lo que hace es transformar la recta $x=0$ (inicial) en la recta $x = \lm t$, y se mantienen los valores 
originales $u^{-}$ y $u^{+}$. \\

Queremos ver entonces cuándo \eqref{eq:ec13} es solución de \eqref{eq:ec9}, cuando tenemos un dato inicial de la forma \eqref{eq:ec14}. \\

El siguiente resultado será de gran utilidad:

\begin{teo}[Condición de salto de Rankine-Hugoniot]\label{teo:t2}
    Sean\footnote{Se impone que $u^{+} \neq u^{-}$ para que exista un salto. En otro caso la función es constante y no hay salto.} $\lm, u^{+}, u^{-} \in \R$ : $u^{+} \neq u^{-}$. Entonces la función \eqref{eq:ec13} es una solución débil del problema de Riemann asociado a \eqref{eq:ec9}
    si y solamente si $$\lm = \dfrac{\dfrac{1}{2}(u^{+})^2 - \dfrac{1}{2}(u^{-})^2}{u^{+} - u^{-}} = \dfrac{1}{2}(u^{+} + u^{-})$$ 
\end{teo}

\begin{observacion}
    El significado físico de $\lm$ es la velocidad a la que varía la recta. Geométricamente sería la pendiente. Aunque lo vamos a utilizar para 
    Burgers, el Teorema \ref{teo:t2} vale para cualquier modelo no lineal:
    \begin{equation}
        \begin{cases}
            \dfrac{\partial u}{\partial t} + \dfrac{\partial}{\partial x}(f(u)) = 0 \\
            u(0,x) = u_0(x)
        \end{cases}
    \end{equation}
    y en tal caso el valor de $\lm$ es:
    $$\lm = \dfrac{f(u^{+}) - f(u^{-})}{u^{+} - u^{-}}$$
\end{observacion}

Antes de demostrar el Teorema \ref{teo:t2}, veamos qué significa ser solución débil del problema de Burgers \eqref{eq:ec9}:

\begin{definicion}[Solución débil de \eqref{eq:ec9}]
    Buscamos una función $u \in L_{loc}^1(\left[0, T\right[ \times \R)$ tal que
    $$\int_{0}^{T} \int_{-\infty}^{+\infty} \left( \dfrac{\partial u}{\partial t} + \dfrac{\partial}{\partial x} \left( \dfrac{u^2}{2} \right) \right) \varphi \,dx\,dt = 0$$
    Integrando por partes:
    $$\int_{-\infty}^{+\infty} u(0,x) \varphi(0,x)\,dx + \int_{0}^{T} \int_{-\infty}^{+\infty} u \dfrac{\partial \varphi}{\partial t}\,dx\,dt
    + \int_{0}^{T} \int_{-\infty}^{+\infty} \dfrac{u^2}{2} \dfrac{\partial \varphi}{\partial x}\,dx\,dt = 0$$
    Así pues, si $u \in L_{loc}^1(\left[0, T\right[ \times \R)$, entonces 
    \begin{equation}\label{eq:ec15}
        \boxed{\int_{0}^{T} \int_{-\infty}^{+\infty} \left[ u \dfrac{\partial \varphi}{\partial t} + \dfrac{u^2}{2} \dfrac{\partial \varphi}{\partial x} \right]\,dx\,dt = - 
        \int_{-\infty}^{+\infty} u_0(x) \varphi(0,x)\,dx}
    \end{equation}
    para toda $\varphi \in C_c^1(\left[0, T\right[ \times \R)$.
\end{definicion}

Recordamos para la demostración del Teorema \ref{teo:t2} el Teorema de la Divergencia\footnote{También conocido como Teorema de Gauss.}:

\begin{teo}[de la Divergencia]
    Sea $\Omega \subset \R^N$ un abierto cuya frontera $\partial \Omega$ es de clase $C^1$, y sea $A \in C^1(\bar{\Omega}, \R^N)$ un campo vectorial.
    Entonces
    $$\int_{\Omega} \di(A)\,dx = \int_{\partial \Omega} A \cdot \vec{\eta} \, d \sigma$$
    donde $\vec{\eta}$ es el vector normal unitario exterior a $\partial \Omega$.
\end{teo}

\begin{proof}
    La función \eqref{eq:ec13} cumple \eqref{eq:ec15} si 
    \begin{equation}\label{eq:ec16}
        \begin{aligned}
            &\overbrace{\int_{0}^{T} \int_{-\infty}^{\lm t}
            \left[ u^{-} \dfrac{\partial \varphi}{\partial t}
            + \dfrac{(u^{-})^2}{2}
            \dfrac{\partial \varphi}{\partial x} \right]
            \,dx\,dt}^{I_1} \\
            &+ \overbrace{\int_{0}^{T} \int_{\lm t}^{+\infty}
            \left[ u^{+} \dfrac{\partial \varphi}{\partial t}
            + \dfrac{(u^{+})^2}{2}
            \dfrac{\partial \varphi}{\partial x} \right]
            \,dx\,dt}^{I_2} = \\
            &- \int_{-\infty}^{0} u^{-}\varphi(0,x)\,dx
            - \int_{0}^{+\infty} u^{+}\varphi(0,x)\,dx
        \end{aligned}
    \end{equation}
    Para ver esto solo hay que aplicar el Teorema de la Divergencia y sale directamente. Consideramos 
    $$A^{\mp} = \left(u^{\mp} \varphi, \dfrac{(u^{\mp})^2}{2} \varphi \right)$$
    y $$\di_{(t,x)}(A^{\mp}) = \dfrac{\partial u^{\mp} \varphi}{\partial t} + \dfrac{\partial}{\partial x} \left( \dfrac{(u^{\mp})^2}{2} \varphi \right) = $$
    $$u^{\mp} \dfrac{\partial \varphi}{\partial t} + \dfrac{(u^{\mp})^2}{2} \dfrac{\partial \varphi}{\partial x}$$
    El resultado de aplicar el Teorema de la Divergencia es:
    $$I_1 = - \int_{-\infty}^{0} u^{-} \varphi(0,x) \,dx - \int_{0}^{T} \varphi(t, \lm t) \left( \lm u^{-} - \dfrac{(u^{-})^2}{2} \right)$$
    $$I_2 = - \int_{0}^{+\infty} u^{+} \varphi(0,x) \,dx + \int_{0}^{T} \varphi(t, \lm t) \left( \lm u^{+} - \dfrac{(u^{+})^2}{2} \right)$$
    Si sustituimos en \eqref{eq:ec16} el resultado obtenido en $I_1$ e $I_2$, nos queda:
    $$\int_{0}^{T} \varphi(t,\lm t) \left( \lm u^{+} - \dfrac{(u^+)^2}{2} - \lm u^{-} + \dfrac{(u^-)^2}{2} \right) \, dt = 0$$
    para toda función $\varphi$ test, por lo que la única alternativa es que, al igualar a $0$ el otro factor del integrando, y despejando $\lm$, resulta que:
    $$\lm = \dfrac{\dfrac{1}{2}(u^{+})^2 - \dfrac{1}{2}(u^{-})^2}{u^{+} - u^{-}} = \dfrac{1}{2}(u^{+} + u^{-})$$
    como se quería demostrar\footnote{Al final hemos aplicado la versión débil del lema fundamental del cálculo de variaciones.}.
\end{proof}

\begin{ejercicio} % TODO: Hacer
    Comprobar que en efecto salen esos valores de $I_1$ e $I_2$ tras aplicar el Teorema de la Divergencia.
\end{ejercicio}

\begin{observacion}
    Si consideramos $$A^{-} = \left(u^{-} \varphi, \dfrac{(u^-)^2}{2} \varphi \right)$$ y
    $$\Omega = \left[0, T \right[ \times \left]-\infty, \lm t \right[$$
    Gráficamente:
    \begin{figure}[H]
        \begin{center}
            \tikzset{every picture/.style={line width=0.75pt}} %set default line width to 0.75pt        

            \begin{tikzpicture}[x=0.75pt,y=0.75pt,yscale=-1,xscale=1]
            %uncomment if require: \path (0,300); %set diagram left start at 0, and has height of 300

            %Shape: Axis 2D [id:dp3242275334401292] 
            \draw  (230.67,150) -- (360,150)(240,38) -- (240,210) (353,145) -- (360,150) -- (353,155) (235,45) -- (240,38) -- (245,45)  ;
            %Straight Lines [id:da8200181075779964] 
            \draw    (299.83,112.5) -- (351.67,80) ;
            %Straight Lines [id:da10207814113213398] 
            \draw    (300,139) -- (300,161) ;
            %Straight Lines [id:da4016638819857854] 
            \draw  [dash pattern={on 0.84pt off 2.51pt}]  (299.83,112.5) -- (299.83,210.5) ;
            %Straight Lines [id:da7811408021976621] 
            \draw [color={rgb, 255:red, 0; green, 59; blue, 255 }  ,draw opacity=1 ]   (278,128) -- (300,150) ;
            %Straight Lines [id:da45167587101332696] 
            \draw [color={rgb, 255:red, 0; green, 59; blue, 255 }  ,draw opacity=1 ]   (268,136) -- (295,165) ;
            %Straight Lines [id:da8532721754574758] 
            \draw [color={rgb, 255:red, 0; green, 59; blue, 255 }  ,draw opacity=1 ]   (256,146) -- (278,168) ;
            %Straight Lines [id:da21649060037233314] 
            \draw [color={rgb, 255:red, 0; green, 59; blue, 255 }  ,draw opacity=1 ]   (275,165) -- (297,187) ;
            %Straight Lines [id:da7942393059803331] 
            \draw [color={rgb, 255:red, 0; green, 59; blue, 255 }  ,draw opacity=1 ]   (248,162) -- (270,184) ;
            %Straight Lines [id:da6329598689513399] 
            \draw [color={rgb, 255:red, 0; green, 59; blue, 255 }  ,draw opacity=1 ]   (259,173) -- (290,203) ;
            %Straight Lines [id:da5635607272548492] 
            \draw [color={rgb, 255:red, 0; green, 59; blue, 255 }  ,draw opacity=1 ]   (243,174) -- (265,196) ;
            %Straight Lines [id:da6741985877594474] 
            \draw [color={rgb, 255:red, 0; green, 59; blue, 255 }  ,draw opacity=1 ]   (262,193) -- (284,215) ;
            %Curve Lines [id:da5966350701002826] 
            \draw [color={rgb, 255:red, 0; green, 59; blue, 255 }  ,draw opacity=1 ]   (279,229) .. controls (295.66,232.92) and (298.87,235.88) .. (313.11,211.53) ;
            \draw [shift={(314,210)}, rotate = 119.98] [color={rgb, 255:red, 0; green, 59; blue, 255 }  ,draw opacity=1 ][line width=0.75]    (10.93,-3.29) .. controls (6.95,-1.4) and (3.31,-0.3) .. (0,0) .. controls (3.31,0.3) and (6.95,1.4) .. (10.93,3.29)   ;
            %Straight Lines [id:da4584837957354798] 
            \draw [color={rgb, 255:red, 253; green, 0; blue, 0 }  ,draw opacity=1 ]   (240,180) -- (211,180) ;
            \draw [shift={(209,180)}, rotate = 360] [color={rgb, 255:red, 253; green, 0; blue, 0 }  ,draw opacity=1 ][line width=0.75]    (10.93,-3.29) .. controls (6.95,-1.4) and (3.31,-0.3) .. (0,0) .. controls (3.31,0.3) and (6.95,1.4) .. (10.93,3.29)   ;
            %Straight Lines [id:da24433253893145346] 
            \draw [color={rgb, 255:red, 253; green, 0; blue, 0 }  ,draw opacity=1 ]   (300,120) -- (326,120) ;
            \draw [shift={(328,120)}, rotate = 180] [color={rgb, 255:red, 253; green, 0; blue, 0 }  ,draw opacity=1 ][line width=0.75]    (10.93,-3.29) .. controls (6.95,-1.4) and (3.31,-0.3) .. (0,0) .. controls (3.31,0.3) and (6.95,1.4) .. (10.93,3.29)   ;
            %Straight Lines [id:da8693895743082151] 
            \draw [color={rgb, 255:red, 253; green, 0; blue, 0 }  ,draw opacity=1 ]   (271,131) -- (257.16,111.63) ;
            \draw [shift={(256,110)}, rotate = 54.46] [color={rgb, 255:red, 253; green, 0; blue, 0 }  ,draw opacity=1 ][line width=0.75]    (10.93,-3.29) .. controls (6.95,-1.4) and (3.31,-0.3) .. (0,0) .. controls (3.31,0.3) and (6.95,1.4) .. (10.93,3.29)   ;
            %Straight Lines [id:da002150546491489602] 
            \draw  [dash pattern={on 0.84pt off 2.51pt}]  (240,150) -- (299.83,112.5) ;

            % Text Node
            \draw (219,24) node [anchor=north west][inner sep=0.75pt]   [align=left] {$\displaystyle x$};
            % Text Node
            \draw (363,62) node [anchor=north west][inner sep=0.75pt]   [align=left] {$\displaystyle x\ =\ \lambda t$};
            % Text Node
            \draw (367,147) node [anchor=north west][inner sep=0.75pt]   [align=left] {$\displaystyle t$};
            % Text Node
            \draw (303,154) node [anchor=north west][inner sep=0.75pt]   [align=left] {$\displaystyle T$};
            % Text Node
            \draw (316,190) node [anchor=north west][inner sep=0.75pt]   [align=left] {$\displaystyle \textcolor[rgb]{0,0.23,1}{\Omega }$};
            % Text Node
            \draw (159,170) node [anchor=north west][inner sep=0.75pt]   [align=left] {$\displaystyle \textcolor[rgb]{0.99,0,0}{( -1,0)}$};
            % Text Node
            \draw (334,110) node [anchor=north west][inner sep=0.75pt]   [align=left] {$\displaystyle \textcolor[rgb]{0.99,0,0}{(}\textcolor[rgb]{0.99,0,0}{1,0}\textcolor[rgb]{0.99,0,0}{)}$};


            \end{tikzpicture}
        \end{center}
    \end{figure}
\end{observacion}
% Mañana vemos los dos tipos específicos de soluciones de la ecuación de Burgers, que son las ondas de choque, y las ondas de rarefacción.
% El lunes o martes será la última clase, seguro. En principio el lunes se resuelven los ejercicios y terminaríamos el bloque. Ya empezaría
% La parte siguiente.
El choque en las trayectorias características se traduce en un salto en la solución. \\

También se comentó al principio el problema de la no unicidad. Veamos el siguiente ejemplo:
\begin{ejemplo}
    Consideramos el siguiente dato inicial para el problema de Riemann:
    \begin{equation*}
        u_0(x) = \begin{cases}
            1 \quad& x<0 \\
            0 \quad& x>0
        \end{cases}
    \end{equation*}
    Por el Teorema \ref{teo:t2}, la solución es 
    \begin{equation*}
        u(t,x) = \begin{cases}
            1 \quad& x < \lm t \\
            0 \quad& x > \lm t
        \end{cases}
    \end{equation*}
    Para $\lm = 1/2$ (valor dado por la condición de salto de Rankine-Hugoniot). Este tipo de soluciones se denominan \underline{ondas de choque}.
    Aquí el dato inicial era decreciente y sabemos entonces que las trayectorias características se van a encontrar.
\end{ejemplo}

\begin{observacion}
    Si la función tuviera más trozos, habría que comparar los trozos dos a dos, pero el teorema seguiría siendo válido. Veremos un ejemplo con tres trozos.
    Lo importante es que sea \textbf{constante a trozos}. \\

    Además, si el dato inicial fuera de la forma 
    $$u_0(x) = \begin{cases}
        u^{-} \quad& x < a \\
        u^{+} \quad& x > a
    \end{cases}$$
    para cualquier $a \in \R$, se podría considerar el cambio de variable $y = x-a$, aplicar el Teorema \ref{teo:t2}, y deshacer el cambio, por 
    lo que sigue siendo válido para traslaciones de la abscisa en la que se produce el salto en el dato inicial. 
\end{observacion}

\begin{ejemplo}
    En este caso, consideramos el dato inicial \underline{creciente}
    \begin{equation*}
        u_0(x) = \begin{cases}
            0 \quad& x < 0 \\
            1 \quad& x > 0
        \end{cases}
    \end{equation*}
    Nuevamente, por el Teorema \ref{teo:t2}, la solución es 
    \begin{equation*}
        u(t,x) = \begin{cases}
            0 \quad& x < t/2 \\
            1 \quad& x > t/2
        \end{cases}
    \end{equation*}
    En este caso, tenemos infinitas soluciones, ya que, para cada $\alpha \in \left]0, 1 \right[$, tenemos las siguientes soluciones (ondas de choque también):
    \begin{equation*}
        u_\alpha(t,x) = \begin{cases}
            0 \quad& x < \frac{\alpha t}{2} \\
            \alpha \quad& \frac{\alpha t}{2} < x < \frac{(\alpha+1) t}{2} \\
            1 \quad& x > \frac{(\alpha+1) t}{2} \\
        \end{cases}
    \end{equation*}
    Pero vamos más allá. Este problema tiene más soluciones, de una naturaleza \textbf{completamente distinta} a las ondas de choque. Una de estas
    soluciones nuevas es:
    \begin{equation}\label{eq:ec18}
        u(t,x) = \begin{cases}
            0 \quad& x < 0 \\
            \frac{x}{t} \quad& 0 \leq x \leq t \\
            1 \quad& x > t
        \end{cases}
    \end{equation}
    Esta solución es destacable porque es \textbf{continua}. Esto se debe a que el dato inicial es \textbf{creciente}\footnote{Si fuera
    decreciente, sabemos que las características chocan y se pierde la continuidad.}.
    En términos de las trayectorias características de la ecuación de Burgers (que, recordemos, son rectas con origen el punto de arranque 
    y pendiente el dato inicial, de la forma \eqref{eq:ec17}). Gráficamente\footnote{Se han intercambiado los ejes para ver con más claridad el lugar que ocupan las ondas de rarefacción.}:
    \begin{figure}[H]
        \begin{center}
            \tikzset{every picture/.style={line width=0.75pt}} %set default line width to 0.75pt        

            \begin{tikzpicture}[x=0.75pt,y=0.75pt,yscale=-1,xscale=1]
            %uncomment if require: \path (0,300); %set diagram left start at 0, and has height of 300

            %Shape: Axis 2D [id:dp3242275334401292] 
            \draw  (168,150) -- (360,150)(240,38) -- (240,210) (353,145) -- (360,150) -- (353,155) (235,45) -- (240,38) -- (245,45)  ;
            %Straight Lines [id:da8200181075779964] 
            \draw    (240,150) -- (321.67,69.67) ;
            %Straight Lines [id:da7811408021976621] 
            \draw [color={rgb, 255:red, 0; green, 59; blue, 255 }  ,draw opacity=1 ]   (230,129) -- (230,191) ;
            %Straight Lines [id:da9902153940282713] 
            \draw [color={rgb, 255:red, 0; green, 59; blue, 255 }  ,draw opacity=1 ]   (220,129) -- (220,191) ;
            %Straight Lines [id:da46681775109807866] 
            \draw [color={rgb, 255:red, 0; green, 59; blue, 255 }  ,draw opacity=1 ]   (210,129) -- (210,191) ;
            %Straight Lines [id:da2862107334178614] 
            \draw [color={rgb, 255:red, 0; green, 59; blue, 255 }  ,draw opacity=1 ]   (200,129) -- (200,191) ;
            %Straight Lines [id:da891584242145164] 
            \draw [color={rgb, 255:red, 0; green, 59; blue, 255 }  ,draw opacity=1 ]   (240.67,229.67) -- (200.67,229.67) ;
            \draw [shift={(198.67,229.67)}, rotate = 360] [color={rgb, 255:red, 0; green, 59; blue, 255 }  ,draw opacity=1 ][line width=0.75]    (10.93,-3.29) .. controls (6.95,-1.4) and (3.31,-0.3) .. (0,0) .. controls (3.31,0.3) and (6.95,1.4) .. (10.93,3.29)   ;
            %Straight Lines [id:da02222492658014852] 
            \draw [color={rgb, 255:red, 255; green, 0; blue, 0 }  ,draw opacity=1 ]   (280,170) -- (361.67,89.67) ;
            %Straight Lines [id:da4338095240978027] 
            \draw [color={rgb, 255:red, 255; green, 0; blue, 0 }  ,draw opacity=1 ]   (290,190) -- (371.67,109.67) ;
            %Straight Lines [id:da7151116176791149] 
            \draw [color={rgb, 255:red, 255; green, 0; blue, 0 }  ,draw opacity=1 ]   (269,156) -- (350.67,75.67) ;
            %Straight Lines [id:da012982910756638821] 
            \draw [color={rgb, 255:red, 14; green, 255; blue, 0 }  ,draw opacity=1 ]   (240,150) -- (304.67,58.67) ;
            %Straight Lines [id:da44861127255647437] 
            \draw [color={rgb, 255:red, 14; green, 255; blue, 0 }  ,draw opacity=1 ]   (240,150) -- (289.67,57.67) ;
            %Straight Lines [id:da9915204219627409] 
            \draw [color={rgb, 255:red, 14; green, 255; blue, 0 }  ,draw opacity=1 ]   (240,150) -- (271.67,56.67) ;
            %Straight Lines [id:da1270037446521043] 
            \draw [color={rgb, 255:red, 14; green, 255; blue, 0 }  ,draw opacity=1 ]   (240,150) -- (253.67,55.67) ;

            % Text Node
            \draw (367,144) node [anchor=north west][inner sep=0.75pt]   [align=left] {$\displaystyle x$};
            % Text Node
            \draw (328,52) node [anchor=north west][inner sep=0.75pt]   [align=left] {$\displaystyle x=t$};
            % Text Node
            \draw (225,21) node [anchor=north west][inner sep=0.75pt]   [align=left] {$\displaystyle t$};
            % Text Node
            \draw (177,147) node [anchor=north west][inner sep=0.75pt]   [align=left] {\textcolor[rgb]{0,0.23,1}{...}};
            % Text Node
            \draw (199,237.67) node [anchor=north west][inner sep=0.75pt]   [align=left] {$\displaystyle \textcolor[rgb]{0,0.23,1}{x< 0}$};


            \end{tikzpicture}
        \end{center}
    \end{figure}
    donde se muestran las siguientes trayectorias características: 
    \begin{equation*}
        X(t;0,x) = \begin{cases}
            x \quad& x < 0 \\
            x + t \quad& x > 0
        \end{cases}
    \end{equation*}
\end{ejemplo}

Como podemos ver, las trayectorias características dejan un hueco entre $x=0$ y $x=t$. Las rectas verdes con las que se rellena este hueco 
dan pie a la siguiente definición: 

\begin{definicion}[Ondas de Rarefacción]
    Las soluciones al problema de Riemann (\ref{definicion:d21}) de la forma \eqref{eq:ec18} se denominan \underline{ondas de rarefacción}. 
    Nuevamente, es importante recalcar que solo es posible encontrar este tipo de soluciones cuando el dato inicial es \underline{creciente}.
\end{definicion}

\begin{observacion}
    Si cambiamos de escala, es decir, pasamos de $(t,x) \mapsto (\alpha t, \alpha x)$, para cierto $\alpha \in \R$, la solución sigue siendo válida.
    En efecto, tenemos que 
    $$u(t,x) = U \left( \dfrac{x}{t} \right) \rightsquigarrow 
    U'\left( \dfrac{x}{t} \right) \left( - \dfrac{x}{t^2} \right) + U \left( \dfrac{x}{t} \right) U' \left( \dfrac{x}{t} \right) \dfrac{1}{t} = 0$$
    $$U' \left( \dfrac{x}{t} \right) \neq 0 \Longrightarrow - \dfrac{x}{t^2} + U \left( \dfrac{x}{t} \right) \dfrac{1}{t} = 0 \iff$$
    \begin{equation}\label{eq:ec19}
        U \left( \dfrac{x}{t} \right) = \dfrac{x}{t}
    \end{equation}
    Esta última se conoce como solución \underline{autosemejante} o \underline{autosimilar}. Es la que ``rellena el hueco'' que dejan las ondas 
    de rarefacción.
\end{observacion}

\begin{observacion}
    Algunas conclusiones:
    \begin{itemize}
        \item Si el modelo de transporte no podemos aspirar a recuperar las propiedades ``buenas''.
        \item No podemos decir que el problema esté bien planteado en el sentido de la existencia y unicidad.
        \item Es relevante la monotonía del dato inicial.
        \begin{itemize}
            \item Si es decreciente, las características se encontrarán y tendremos un salto.
            \item Si es creciente, aparte de las ondas de choque, podemos encontrar perfiles autosimilares del tipo \eqref{eq:ec19}, generando 
            ondas de rarefacción.
        \end{itemize}
    \end{itemize}
\end{observacion}

Algunos ejercicios para concluir el bloque:

\begin{ejercicio}
    Resolver los siguientes problemas (la ecuación es la de Burgers siempre, y solo cambia el dato inicial):
    $$\dfrac{\partial u}{\partial t} + u \dfrac{\partial u}{\partial x} = 0$$
    para los siguientes datos iniciales, y escribir cómo son las trayectorias características:
    \begin{itemize}
        \item $u_0(x) = \begin{cases}
            2 \quad& x < 0 \\
            1 \quad& 0 < x < 2 \\
            0 \quad& x < 2
        \end{cases}$. \\

        En primer lugar obtenemos las trayectorias características:
        $$X(t; 0, x) = x + u_0(x) t = \begin{cases}
            x + 2t \quad& x < 0 \\
            x + t \quad& 0 < x < 2 \\
            x \quad& x > 2
        \end{cases}$$
        Gráficamente\footnote{El dibujo siempre ayuda en estos casos.}:
        \begin{figure}[H]
            \begin{center}
                \tikzset{every picture/.style={line width=0.75pt}} %set default line width to 0.75pt        
                \begin{tikzpicture}[x=0.75pt,y=0.75pt,yscale=-1,xscale=1]
                    %uncomment if require: \path (0,300); %set diagram left start at 0, and has height of 300

                    %Shape: Axis 2D [id:dp3242275334401292] 
                    \draw  (130,150) -- (398,150)(240,38) -- (240,210) (391,145) -- (398,150) -- (391,155) (235,45) -- (240,38) -- (245,45)  ;
                    %Straight Lines [id:da2862107334178614] 
                    \draw [color={rgb, 255:red, 0; green, 59; blue, 255 }  ,draw opacity=1 ]   (396,55) -- (185,166) ;
                    %Straight Lines [id:da02222492658014852] 
                    \draw [color={rgb, 255:red, 255; green, 0; blue, 0 }  ,draw opacity=1 ]   (319.91,191.02) -- (320.46,78.46) ;
                    %Straight Lines [id:da4338095240978027] 
                    \draw [color={rgb, 255:red, 255; green, 0; blue, 0 }  ,draw opacity=1 ]   (340.21,191.83) -- (339.76,79.27) ;
                    %Straight Lines [id:da7151116176791149] 
                    \draw [color={rgb, 255:red, 255; green, 0; blue, 0 }  ,draw opacity=1 ]   (299.91,190.39) -- (300.45,78.84) ;
                    %Straight Lines [id:da012982910756638821] 
                    \draw [color={rgb, 255:red, 14; green, 255; blue, 0 }  ,draw opacity=1 ]   (193,196) -- (374,18) ;
                    %Straight Lines [id:da8509471356906648] 
                    \draw    (300,139) -- (300,162) ;
                    %Straight Lines [id:da3794350886499228] 
                    \draw [color={rgb, 255:red, 255; green, 0; blue, 0 }  ,draw opacity=1 ]   (330.91,190.02) -- (330.46,77.46) ;
                    %Straight Lines [id:da6561401855095113] 
                    \draw [color={rgb, 255:red, 255; green, 0; blue, 0 }  ,draw opacity=1 ]   (351.21,190.83) -- (349.76,78.27) ;
                    %Straight Lines [id:da27040052523613545] 
                    \draw [color={rgb, 255:red, 255; green, 0; blue, 0 }  ,draw opacity=1 ]   (310.91,189.39) -- (310.45,77.84) ;
                    %Straight Lines [id:da3632659939494831] 
                    \draw [color={rgb, 255:red, 14; green, 255; blue, 0 }  ,draw opacity=1 ]   (203,208) -- (373,36) ;
                    %Straight Lines [id:da97824068996717] 
                    \draw [color={rgb, 255:red, 14; green, 255; blue, 0 }  ,draw opacity=1 ]   (243.1,208) -- (393,56) ;
                    %Straight Lines [id:da12914183632548215] 
                    \draw [color={rgb, 255:red, 14; green, 255; blue, 0 }  ,draw opacity=1 ]   (230.1,201) -- (380,49) ;
                    %Straight Lines [id:da5712018545643891] 
                    \draw [color={rgb, 255:red, 0; green, 59; blue, 255 }  ,draw opacity=1 ]   (390,48) -- (179,159) ;
                    %Straight Lines [id:da8348964328891428] 
                    \draw [color={rgb, 255:red, 0; green, 59; blue, 255 }  ,draw opacity=1 ]   (373,45) -- (162,156) ;

                    % Text Node
                    \draw (406,141) node [anchor=north west][inner sep=0.75pt]   [align=left] {$\displaystyle x$};
                    % Text Node
                    \draw (225,21) node [anchor=north west][inner sep=0.75pt]   [align=left] {$\displaystyle t$};
                    % Text Node
                    \draw (286,153.5) node [anchor=north west][inner sep=0.75pt]   [align=left] {$\displaystyle 2$};
                    % Text Node
                    \draw (357,113) node [anchor=north west][inner sep=0.75pt]   [align=left] {\textcolor[rgb]{1,0,0}{...}};
                    % Text Node
                    \draw (257.69,78.97) node [anchor=north west][inner sep=0.75pt]  [rotate=-63.09] [align=left] {\textcolor[rgb]{0,0.23,1}{...}};

                \end{tikzpicture}
            \end{center}
        \end{figure}
        Distinguimos dos a dos: el primer dato inicial que consideramos es 
        $$u_1(0,x) = \begin{cases}
            2 \quad& x < 0 \\
            1 \quad& 0 < x < 2
        \end{cases}$$
        Y aplicando la condición de salto de Rankine-Hugoniot (Teorema \ref{teo:t2}) obtenemos que 
        $$\lm = \dfrac{1}{2}(u_1^{+} + u_1^{-}) = \dfrac{3}{2}$$
        Por lo que una primera solución sería 
        $$u_1(t,x) = \begin{cases}
            2 \quad& x < \frac{3t}{2} \\
            1 \quad& x > \frac{3t}{2}
        \end{cases}$$
        Análogamente, el segundo inicial considerado es 
        $$u_2(0,x) = \begin{cases}
            1 \quad& 0 < x < 2 \\
            0 \quad& x > 2
        \end{cases}$$
        Y aplicando la condición de salto de Rankine-Hugoniot (Teorema \ref{teo:t2}) obtenemos que 
        $$\lm = \dfrac{1}{2}(u_2^{+} + u_2^{-}) = \dfrac{1}{2}$$
        Por lo que una segunda solución sería 
        $$u_2(t,x) = \begin{cases}
            1 \quad& x < \frac{t}{2} + 2 \\
            0 \quad& x > \frac{t}{2} + 2
        \end{cases}$$
        En conclusión, la solución conjunta sería 
        $$u(t,x) = \begin{cases}
            2 \quad& x < \frac{3t}{2} \\
            1 \quad& \frac{3t}{2} < x < \frac{t}{2} + 2 \\
            0 \quad& x > \frac{t}{2} + 2
        \end{cases}$$
        Veamos ahora donde corta. Gráficamente:
        \begin{figure}[H]
            \begin{center}
                \tikzset{every picture/.style={line width=0.75pt}} %set default line width to 0.75pt        

                \begin{tikzpicture}[x=0.75pt,y=0.75pt,yscale=-1,xscale=1]
                %uncomment if require: \path (0,300); %set diagram left start at 0, and has height of 300

                %Shape: Axis 2D [id:dp3242275334401292] 
                \draw  (230.67,150) -- (360,150)(240.67,38) -- (240.67,181) (353,145) -- (360,150) -- (353,155) (235.67,45) -- (240.67,38) -- (245.67,45)  ;
                %Straight Lines [id:da504363558290425] 
                \draw    (300,144) -- (300,156) ;
                %Straight Lines [id:da41768238504122124] 
                \draw  [dash pattern={on 0.84pt off 2.51pt}]  (300.5,150.5) -- (299.5,61.5) ;
                %Shape: Circle [id:dp23983039955771324] 
                \draw  [fill={rgb, 255:red, 47; green, 0; blue, 0 }  ,fill opacity=1 ] (298,60) .. controls (298,59.17) and (298.67,58.5) .. (299.5,58.5) .. controls (300.33,58.5) and (301,59.17) .. (301,60) .. controls (301,60.83) and (300.33,61.5) .. (299.5,61.5) .. controls (298.67,61.5) and (298,60.83) .. (298,60) -- cycle ;
                %Straight Lines [id:da8200181075779964] 
                \draw    (220,180) -- (321,28) ;
                %Straight Lines [id:da2599786825287116] 
                \draw    (241,90) -- (340,38) ;
                %Straight Lines [id:da18823594983513814] 
                \draw    (230,90) -- (250,90) ;
                %Straight Lines [id:da9505158034700008] 
                \draw  [dash pattern={on 0.84pt off 2.51pt}]  (241,60) -- (299.5,60) ;

                % Text Node
                \draw (219,24) node [anchor=north west][inner sep=0.75pt]   [align=left] {$\displaystyle x$};
                % Text Node
                \draw (153,184) node [anchor=north west][inner sep=0.75pt]   [align=left] {$\displaystyle x=\ 3t/2$};
                % Text Node
                \draw (277,155) node [anchor=north west][inner sep=0.75pt]   [align=left] {$\displaystyle t^{*} =2$};
                % Text Node
                \draw (367,150) node [anchor=north west][inner sep=0.75pt]   [align=left] {$\displaystyle t$};
                % Text Node
                \draw (215,74) node [anchor=north west][inner sep=0.75pt]   [align=left] {$\displaystyle 2$};
                % Text Node
                \draw (353,28) node [anchor=north west][inner sep=0.75pt]   [align=left] {$\displaystyle x=\ t/2\ +\ 2$};
                % Text Node
                \draw (187,51) node [anchor=north west][inner sep=0.75pt]   [align=left] {$\displaystyle x^{*} =3$};


                \end{tikzpicture}
            \end{center}
        \end{figure}
        Así pues: 
        $$u(2,x) = \begin{cases}
            2 \quad& x < 3 \\
            0 \quad& x > 3
        \end{cases} \stackrel{R.H}{\Longrightarrow} \lm = 1$$
        Si $x = a + \lm t$, evaluando en $(t,x) = (t^*, x^*) = (2,3)$ se tiene que $a = 1$
        y la solución quedaría 
        $$u(t,x) = \begin{cases}
            2 \quad& x < 1 + t \\
            0 \quad& x > 1 + t
        \end{cases}$$
        donde $(t,x) \in \left[2, +\infty \right[ \times \R$.
        \item $u_0(x) = \begin{cases}
            1 \quad& x < -1 \\
            0 \quad& -1 < x < 1 \\
            -1 \quad& x > 1
        \end{cases}$. Este se resuelve análogo al anterior.
        \item $u_0(x) = \begin{cases}
            1 \quad& x < 0 \\
            1-x \quad& 0 \leq x \leq 1 \\
            0 \quad& x > 1
        \end{cases}$. No lo resuelve.
    \end{itemize}
    Para resolver el primer (y segundo) apartado. El primer ``lote'' de trayectorias características chocará con el segundo, y el segundo con el tercero.
    Tendremos dos rectas de discontinuidad distintas. En algún punto ambas rectas chocarán entre sí. Una vez escrita, tendremos un choque con 
    tres tramos que estará bien definido en el sentido de Rankine-Hugoniot hasta el punto en que choquen ambas rectas de discontinuidad.
    Una vez chocan estas dos rectas, al calcularlo veremos que el tramo central de la solución que hemos escrito, colapsa. Así, los tres
    tramos se convierten en dos, también constantes a trozos. Ahí colocamos la nueva condición inicial para saber qué pasa desde el instante de choque
    en adelante. Como ahora el dato inicial solo tiene dos tramos, aplicamos Rankine-Hugoniot una vez. \\

    Como los datos iniciales son decrecientes, no podemos esperar encontrar soluciones continuas. \\

    Para el tercero, que \textbf{no} es un problema de Riemann, porque el dato inicial no es constante a trozos, usamos una propiedad fundamental,
    que es de lo poco que se ``hereda'' del modelo lineal:
    $$\dfrac{d}{dt} (u(t,X(t;0,x))) = 0 \Longrightarrow u(t,X(t;0,x)) = u_0(x)$$
    Lo que se sugiere es aplicar esta propiedad en cada uno de los tramos. En el primero y en el tercero es trivial. La cosa es aplicar en el 
    tramo central a ver qué sale.
\end{ejercicio}

\begin{ejercicio}[Modelos de Tráfico]
    Denotamos $\rho(t,x)$ el número de vehículos por unidad de longitud a lo largo de una carretera en el instante $t$, \textbf{sin contar el propio}.
    Llamamos $u(t,x)$ a la distribución de velocidades de los vehículos. Ténganse además en cuenta las siguientes consideraciones\footnote{Para simplificar el modelo.}:
    \begin{itemize}
        \item Todos los vehículos son del mismo tipo.
        \item Si no hay más vehículos que el propio ($\rho(t,x) = 0$), entonces la velocidad a la que circulo es la máxima posible, pongamos $u = u_{max}$.
        \item Si el número de vehículos es el máximo posible, pongamos $\rho = \rho_{max}$, entonces no nos podemos mover porque la vía está colapsada,
        es decir, $u = 0$.
    \end{itemize}
    En tales circunstancias, el flujo viene descrito por 
    $$f(\rho) = \rho u$$
    Recordemos que en modelos de transporte el flujo era $f(u)$ en:
    \begin{equation}\label{eq:ec20}
        \dfrac{\partial u}{\partial t} + \dfrac{\partial}{\partial x} \left( f(u) \right) = 0
    \end{equation}
    Entonces 
    $$u(\rho) = u_{max} \left( 1 - \dfrac{\rho}{\rho_{max}} \right)$$
    Se pide:
    \begin{enumerate}
        \item Construir el modelo de transporte no lineal para estos datos (es decir, obtener una ecuación de la forma \eqref{eq:ec20}), y verificar que el modelo resultante
        (que será una ecuación lineal que no podemos identificar con nada (lo que nos saldrá no es Burgers, que es el único modelo lineal que no conocemos))
        es reducible a un modelo de Burgers a través del siguiente cambio de variable:
        \begin{equation}\label{eq:ec21}
            R(t,x) = u_{max} - 2 \dfrac{u_{max}}{\rho_{max}} \rho(t,x)
        \end{equation}
        y obtener una ecuación de la forma 
        $$\dfrac{\partial R}{\partial t} + R \dfrac{\partial R}{\partial x} = 0$$
        De esta manera, podemos sacar conclusiones aplicando la teoría conocida de Burgers, y luego volver al modelo original deshaciendo el cambio. \\

        Resolvemos por orden de aparición. Primeramente tenemos que 
        $$\dfrac{\partial \rho}{\partial t} + \dfrac{\partial}{\partial x} f(\rho) = 0$$
        y desarrollando
        $$\dfrac{\partial \rho}{\partial t} + \dfrac{\partial}{\partial x} (\rho u) = 
        \dfrac{\partial \rho}{\partial t} + \dfrac{\partial \rho}{\partial x} u_{max} \left( 1 - \dfrac{\rho}{\rho_{max}} \right) - 
        \rho \dfrac{\partial \rho}{\partial x} \dfrac{u_{max}}{\rho_{max}} =$$
        $$\boxed{\dfrac{\partial \rho}{\partial t} + u_{max} \dfrac{\partial \rho}{\partial x} - 2 \dfrac{u_{max}}{\rho_{max}} \rho \dfrac{\partial \rho}{\partial x} = 0}$$

        Antes de seguir podríamos haber sacado factor común $\frac{\partial \rho}{\partial x}$ en esta última expresión y tendríamos el cambio de variable multiplicando directamente.
        Si no nos damos cuenta, reducimos a un modelo de Burgers a través del cambio \eqref{eq:ec21}. Despejando $\rho(t,x)$, nos queda:

        $$\rho(t,x) = (u_{max} - R(t,x)) \dfrac{\rho_{max}}{2 u_{max}} \Longrightarrow$$
        $$\dfrac{\partial R}{\partial t} = - 2 \dfrac{u_{max}}{\rho_{max}} \dfrac{\partial \rho}{\partial t} \Longrightarrow$$
        $$\boxed{\dfrac{\partial \rho}{\partial t} = - \dfrac{\rho_{max}}{2 u_{max}} \dfrac{\partial R}{\partial t}}$$
        $$\boxed{\dfrac{\partial \rho}{\partial x} = - \dfrac{\rho_{max}}{2 u_{max}} \dfrac{\partial R}{\partial x}}$$

        Y si seguimos desarrollando:

        $$- \dfrac{\rho_{max}}{2 u_{max}} \dfrac{\partial R}{\partial t} - \dfrac{\rho_{max}}{2} \dfrac{\partial R}{\partial x} + 
        \dfrac{\partial R}{\partial x} \dfrac{\rho_{max}}{2 u_{max}} (u_{max} - R) = 0$$
        $$- \dfrac{\rho_{max}}{2 u_{max}} \dfrac{\partial R}{\partial t} - \dfrac{\rho_{max}}{2} \dfrac{\partial R}{\partial x} + 
        \dfrac{\rho_{max}}{2} \dfrac{\partial R}{\partial x} - \dfrac{\rho_{max}}{2 u_{max}} R \dfrac{\partial R}{\partial x} = 0 \Longrightarrow $$
        $$\boxed{\dfrac{\partial R}{\partial t} + R \dfrac{\partial R}{\partial x} = 0}$$

        que es un modelo de Burgers.

        \item Encontrar las ondas de choque y/o de rarefacción en cada uno de los siguientes casos, considerando cada uno de los siguientes datos iniciales:
        \begin{itemize}
            \item $\rho_0(x) = \begin{cases}
                \frac{1}{2} \rho_{max} \quad& x < 0 \\
                \rho_{max} \quad& x > 0
            \end{cases}$
            Como el dato inicial es creciente\footnote{No necesariamente debe ser $1/2$, podría usarse cualquier otra fracción entre $0$ y $1$}, tendremos ondas de rarefacción con Rankine-Hugoniot, y ondas de rarefacción. Este dato inicial 
            se corresponde con una situación de atasco. \\

            Utilizando el cambio anterior, tenemos 
            $$R_0(x) = \begin{cases}
                0 \quad& x < 0 \\
                -u_{max} \quad& x > 0
            \end{cases}$$
            y aplicando el Teorema \ref{teo:t2}, nos queda 
            $$R(t,x) = \begin{cases}
                0 \quad& x < - \frac{u_{max}}{2} t \\
                -u_{max} \quad& x > - \frac{u_{max}}{2} t
            \end{cases}$$
            Deshaciendo el cambio, en función de $\rho$ nos queda:
            $$\rho(t,x) = \begin{cases}
                \dfrac{\rho_{max}}{2} \quad& x < - \frac{u_{max}}{2} t \\
                \rho_{max} \quad& x > - \frac{u_{max}}{2} t
            \end{cases}$$
            Falta encontrar las ondas de rarefacción. Las características en este caso son:
            $$X(t;0,x) = \begin{cases}
                x - u_{max} t \quad& x < 0 \\
                x \quad& x > 0
            \end{cases}$$
            Gráficamente:
            \begin{figure}[H]
                \begin{center}
                    \tikzset{every picture/.style={line width=0.75pt}} %set default line width to 0.75pt        

                    \begin{tikzpicture}[x=0.75pt,y=0.75pt,yscale=-1,xscale=1]
                    %uncomment if require: \path (0,300); %set diagram left start at 0, and has height of 300

                    %Shape: Axis 2D [id:dp3242275334401292] 
                    \draw  (130,150) -- (398,150)(240,38) -- (240,210) (391,145) -- (398,150) -- (391,155) (235,45) -- (240,38) -- (245,45)  ;
                    %Straight Lines [id:da02222492658014852] 
                    \draw [color={rgb, 255:red, 255; green, 0; blue, 0 }  ,draw opacity=1 ]   (319.91,191.02) -- (320.46,78.46) ;
                    %Straight Lines [id:da4338095240978027] 
                    \draw [color={rgb, 255:red, 255; green, 0; blue, 0 }  ,draw opacity=1 ]   (340.21,191.83) -- (339.76,79.27) ;
                    %Straight Lines [id:da7151116176791149] 
                    \draw [color={rgb, 255:red, 255; green, 0; blue, 0 }  ,draw opacity=1 ]   (299.91,190.39) -- (300.45,78.84) ;
                    %Straight Lines [id:da3794350886499228] 
                    \draw [color={rgb, 255:red, 255; green, 0; blue, 0 }  ,draw opacity=1 ]   (330.91,190.02) -- (330.46,77.46) ;
                    %Straight Lines [id:da6561401855095113] 
                    \draw [color={rgb, 255:red, 255; green, 0; blue, 0 }  ,draw opacity=1 ]   (351.21,190.83) -- (349.76,78.27) ;
                    %Straight Lines [id:da27040052523613545] 
                    \draw [color={rgb, 255:red, 255; green, 0; blue, 0 }  ,draw opacity=1 ]   (310.91,189.39) -- (310.45,77.84) ;
                    %Straight Lines [id:da43272068532874564] 
                    \draw [color={rgb, 255:red, 255; green, 0; blue, 0 }  ,draw opacity=1 ]   (270.91,191.02) -- (271.46,78.46) ;
                    %Straight Lines [id:da3447299153844853] 
                    \draw [color={rgb, 255:red, 255; green, 0; blue, 0 }  ,draw opacity=1 ]   (291.21,191.83) -- (290.76,79.27) ;
                    %Straight Lines [id:da09487745849680329] 
                    \draw [color={rgb, 255:red, 255; green, 0; blue, 0 }  ,draw opacity=1 ]   (250.91,190.39) -- (251.45,78.84) ;
                    %Straight Lines [id:da6111461950814078] 
                    \draw [color={rgb, 255:red, 255; green, 0; blue, 0 }  ,draw opacity=1 ]   (281.91,190.02) -- (281.46,77.46) ;
                    %Straight Lines [id:da9155729958220498] 
                    \draw [color={rgb, 255:red, 255; green, 0; blue, 0 }  ,draw opacity=1 ]   (261.91,189.39) -- (261.45,77.84) ;
                    %Straight Lines [id:da10353936092794236] 
                    \draw    (138,66) -- (240,150) ;
                    %Straight Lines [id:da8684112446016625] 
                    \draw [color={rgb, 255:red, 255; green, 0; blue, 0 }  ,draw opacity=1 ]   (208.4,174.72) -- (123.51,100.81) ;
                    %Straight Lines [id:da5352469293291663] 
                    \draw [color={rgb, 255:red, 255; green, 0; blue, 0 }  ,draw opacity=1 ]   (222.26,159.88) -- (136.73,86.73) ;
                    %Straight Lines [id:da5211199646319873] 
                    \draw [color={rgb, 255:red, 255; green, 0; blue, 0 }  ,draw opacity=1 ]   (214.82,165.74) -- (129.28,92.59) ;
                    %Straight Lines [id:da9407595491423938] 
                    \draw [color={rgb, 255:red, 255; green, 0; blue, 0 }  ,draw opacity=1 ]   (228.69,150.9) -- (142.5,78.5) ;
                    %Straight Lines [id:da12980266937388085] 
                    \draw [color={rgb, 255:red, 255; green, 0; blue, 0 }  ,draw opacity=1 ]   (201.29,180.49) -- (116.51,107.98) ;
                    %Straight Lines [id:da08886092608553076] 
                    \draw [color={rgb, 255:red, 255; green, 255; blue, 0 }  ,draw opacity=1 ]   (161,63) -- (240,150) ;
                    %Straight Lines [id:da1604666647621502] 
                    \draw [color={rgb, 255:red, 255; green, 255; blue, 0 }  ,draw opacity=1 ]   (181,58) -- (240,150) ;
                    %Straight Lines [id:da9931053519803237] 
                    \draw [color={rgb, 255:red, 255; green, 255; blue, 0 }  ,draw opacity=1 ]   (202,50) -- (240,150) ;
                    %Straight Lines [id:da02051653712381518] 
                    \draw [color={rgb, 255:red, 255; green, 255; blue, 0 }  ,draw opacity=1 ]   (223,49) -- (240,150) ;

                    % Text Node
                    \draw (406,141) node [anchor=north west][inner sep=0.75pt]   [align=left] {$\displaystyle x$};
                    % Text Node
                    \draw (225,21) node [anchor=north west][inner sep=0.75pt]   [align=left] {$\displaystyle t$};
                    % Text Node
                    \draw (357,113) node [anchor=north west][inner sep=0.75pt]   [align=left] {\textcolor[rgb]{1,0,0}{...}};
                    % Text Node
                    \draw (69,41) node [anchor=north west][inner sep=0.75pt]   [align=left] {$\displaystyle x\ =-u_{max} t$};


                    \end{tikzpicture}
                \end{center}
            \end{figure}
            Las líneas amarillas muestran dónde se pueden meter las ondas de rarefacción. La frontera está delimitada por 
            $x = 0$ y por $x = - u_{max} t$. De esta manera, los tramos son $x<-u_{max}t$, $-u_{max} t < x < 0$ y $x > 0$.
            De esta manera, las ondas de rarefacción son:
            $$\rho(t,x) = \begin{cases}
                \frac{1}{2} \rho_{max} \quad& x < -u_{max} t \\
                \frac{x}{t} \quad & -u_{max} t \leq x < 0 \\
                \rho_{max} \quad& x > 0
            \end{cases}$$
            \item $\rho_0(x) = \begin{cases}
                \rho_{max} \quad& x < 0 \\
                \frac{1}{2} \rho_{max} \quad& x > 0
            \end{cases}$
            El dato inicial es en este caso decreciente. Este dato inicial 
            se corresponde con una situación que se suele denominar como de luz verde (el semáforo pasa de rojo a verde). \\

            Es completamente análogo al anterior pero cambiando los sentidos de las desigualdades. En este caso no hay ondas de rarefacción:
            Utilizando el cambio anterior, tenemos 
            $$R_0(x) = \begin{cases}
                0 \quad& x > 0 \\
                -u_{max} \quad& x < 0
            \end{cases}$$
            y aplicando el Teorema \ref{teo:t2}, nos queda 
            $$R(t,x) = \begin{cases}
                0 \quad& x > - \frac{u_{max}}{2} t \\
                -u_{max} \quad& x < - \frac{u_{max}}{2} t
            \end{cases}$$
            Deshaciendo el cambio, en función de $\rho$ nos queda:
            $$\rho(t,x) = \begin{cases}
                \dfrac{\rho_{max}}{2} \quad& x > - \frac{u_{max}}{2} t \\
                \rho_{max} \quad& x < - \frac{u_{max}}{2} t
            \end{cases}$$
            \item $\rho_0(x) = \begin{cases}
                \rho_{max} \quad& x > 0 \\
                \frac{1}{2} \rho_{max} \quad& x < 0
            \end{cases}$
        \end{itemize}
    \end{enumerate}
\end{ejercicio}

% El lunes resolveremos los dos ejercicios. El que es un poco más complicado es el c)





    \chapter{Mecánica de Fluidos}
\noindent

\section{Introducción}

El cálculo de variaciones visto en Modelos Matemáticos II se consideraba ``matemática pura''. Pasó de los matemáticos puros, al ámbito 
de los físicos. Ahora en física hay algo en cuarto de mecánica analítica (que es fácil para ellos). Los ingenieros están usando también
la mecánica, y ahora está volviendo a la matemática. \\

El primero fue Euler. Lo que vamos a ver ahora son las ecuaciones de Euler. En París a principios del siglo XIX. Había que canalizar el agua,
ver cómo se podía distribuir. El surgimiento vino con Navier y Stokes. La mecánica de fluidos ha quedado en manos de las matemáticas. Veremos 
tres o cuatro corrientes que surgen de la mecánica de fluidos. \\

Los únicos que le dedican un poco de atención son los ingenieros aeronáuticos. Los que diseñan el coche de Fernando Alonso que se están equivocando. 
Detrás del diseño de los coches y aviones, en general están matemáticos. \\

Otro lugar en el que los fluidos son importantes es en el ámbito del clima. Hoy en día se predice 2 o 3 días antes. Empiezan diciendo, según
nuestros modelos (los modelos son los mismos), lo que cambia es cómo se resuelve el modelo, y cómo resuelve Navier-Stokes. Hay modelos que 
consideran la orografía, pero, en general, cuanto más realista es el modelo, más complicado es de resolver. \\

Vamos a comenzar por el caso $2$-dimensional, y con las ecuaciones de Euler. Tenemos un plano y se trata de describir cómo se mueve un fluido.
Gráficamente:
\begin{figure}[H]
    \begin{center}
        \tikzset{every picture/.style={line width=0.75pt}} %set default line width to 0.75pt        
            \begin{tikzpicture}[x=0.75pt,y=0.75pt,yscale=-1,xscale=1]
            %uncomment if require: \path (0,300); %set diagram left start at 0, and has height of 300

            %Shape: Rectangle [id:dp7209051313285793] 
            \draw   (294.35,76.5) -- (455.26,76.5) -- (310.34,145.5) -- (149.43,145.5) -- cycle ;
            %Straight Lines [id:da648079438822599] 
            \draw    (280,111) -- (280,34) ;
            \draw [shift={(280,32)}, rotate = 90] [color={rgb, 255:red, 0; green, 0; blue, 0 }  ][line width=0.75]    (10.93,-3.29) .. controls (6.95,-1.4) and (3.31,-0.3) .. (0,0) .. controls (3.31,0.3) and (6.95,1.4) .. (10.93,3.29)   ;
            %Curve Lines [id:da02355300563515772] 
            \draw    (264,62) .. controls (305.37,63.97) and (301.14,46.54) .. (261.82,52.7) ;
            \draw [shift={(260,53)}, rotate = 350.31] [color={rgb, 255:red, 0; green, 0; blue, 0 }  ][line width=0.75]    (10.93,-3.29) .. controls (6.95,-1.4) and (3.31,-0.3) .. (0,0) .. controls (3.31,0.3) and (6.95,1.4) .. (10.93,3.29)   ;
            %Shape: Circle [id:dp28371925872083015] 
            \draw  [color={rgb, 255:red, 0; green, 0; blue, 0 }  ,draw opacity=1 ][fill={rgb, 255:red, 0; green, 0; blue, 0 }  ,fill opacity=1 ] (281.83,112.83) .. controls (281.83,111.82) and (281.01,111) .. (280,111) .. controls (278.99,111) and (278.17,111.82) .. (278.17,112.83) .. controls (278.17,113.85) and (278.99,114.67) .. (280,114.67) .. controls (281.01,114.67) and (281.83,113.85) .. (281.83,112.83) -- cycle ;
            %Straight Lines [id:da556275463205884] 
            \draw    (280,113) -- (295.92,88.35) ;
            \draw [shift={(297,86.67)}, rotate = 122.85] [color={rgb, 255:red, 0; green, 0; blue, 0 }  ][line width=0.75]    (10.93,-3.29) .. controls (6.95,-1.4) and (3.31,-0.3) .. (0,0) .. controls (3.31,0.3) and (6.95,1.4) .. (10.93,3.29)   ;
            %Curve Lines [id:da8523005951397323] 
            \draw    (292,105.67) .. controls (299.92,128.11) and (349.98,157.08) .. (389.8,128.55) ;
            \draw [shift={(391,127.67)}, rotate = 143.13] [color={rgb, 255:red, 0; green, 0; blue, 0 }  ][line width=0.75]    (10.93,-3.29) .. controls (6.95,-1.4) and (3.31,-0.3) .. (0,0) .. controls (3.31,0.3) and (6.95,1.4) .. (10.93,3.29)   ;

            % Text Node
            \draw (295,32) node [anchor=north west][inner sep=0.75pt]   [align=left] {$\displaystyle \omega =\ \text{rot} \ u$};
            % Text Node
            \draw (221.89,113) node [anchor=north west][inner sep=0.75pt]   [align=left] {$\displaystyle ( x_{1} ,x_{2})$};
            % Text Node
            \draw (305,80) node [anchor=north west][inner sep=0.75pt]   [align=left] {$\displaystyle u\ =\ ( u_{1} ,u_{2})$};
            % Text Node
            \draw (393,110) node [anchor=north west][inner sep=0.75pt]   [align=left] {$\displaystyle x( t) \ =\ ( x_{1} ,x_{2})( t)$};
        \end{tikzpicture}
    \end{center}
\end{figure}
Ya vimos en \ref{section:rot} la definición del rotacional, en inglés \textit{curl}. Cada punto está afectado por una presión. Se trata de ver cuáles son las ecuaciones que verifican el modelo anterior.
Consideramos\footnote{A priori no sabemos si será $[0,T]$.}
\begin{equation}\label{eq:ec24}
    u : \left[0, T \right[ \times \R^2 \to \R^2
\end{equation}
    
y la ecuación 

\begin{equation}\label{eq:ec22}
    \dfrac{\partial u}{\partial t} + (u \cdot \nabla) u + \nabla p = 0
\end{equation}
¿Qué significa \eqref{eq:ec22}? Analizamos sumando a sumando:
\begin{itemize}
    \item En primer lugar, $$u \cdot \nabla = \sum_{j=1}^{2} u_j \dfrac{\partial}{\partial x_j}$$
    \item Como $p$ es un escalar $$\nabla p = \left( \dfrac{\partial p}{\partial x_1}, \dfrac{\partial p}{\partial x_2} \right)$$ 
\end{itemize}

Hay muchos tipos de fluido. El aire que tenemos aquí es un fluido. Si tuviéramos una jeringa, podríamos comprimir el aire. Nosotros vamos a 
suponer ahora que estudiamos fluidos incompresibles, es decir, $\di(u) = 0$, donde
$$\di(u) = \dfrac{\partial u_1}{\partial x_1} + \dfrac{\partial u_2}{\partial x_2}$$
Vemos ahora que a $u$ se le asocia la ecuación
$$\dfrac{\partial u}{\partial t} + (a \cdot \nabla) u = 0$$
$$\dfrac{\partial X}{\partial t} = \di(a)$$
Vemos que en lugar de $u \cdot \nabla$ ponemos $a \cdot \nabla$. \\

Una condición natural es que $|u(t,x)| \to 0$ cuando $|x| \to \infty$. Ahora nos preguntamos qué pasa en el tiempo inicial, para lo cual
tendremos en cuenta una condición inicial: $u(t=0,x) = u_0(x)$. Recopilando todo:

\begin{definicion}[Ecuaciones de Euler]
    Las ecuaciones de Euler en $\R^N$, con $N=2,3$, son:
    \begin{equation}\label{eq:ec23}
        \begin{cases}
            \dfrac{\partial u}{\partial t} + (u \cdot \nabla) u + \nabla p = 0 \\
            \di(u) = 0 \\
            |u(t,x)| \to 0 \quad |x| \to +\infty \\
            u(t=0,x) = u_0(x)
        \end{cases}
    \end{equation}
\end{definicion}

\begin{observacion}
    Vemos que las ecuaciones son las mismas tanto para $N=2$ como para $N=3$. Para $N=4$ puede tener sentido, por ejemplo, en el ámbito de la gravitación.
    En el caso bidimensional, podemos tomar componentes en la primera ecuación de \eqref{eq:ec23}, y obtener
    \begin{equation}
        \dfrac{\partial u_i}{\partial t} + (u \cdot \nabla) u_i + \dfrac{\partial p}{\partial x_i} = 0 \quad i=1,2
    \end{equation}
\end{observacion}

Otro vector importante a tener en cuenta es el vector vorticidad $\omega$. ¿Por qué? Porque muchas veces lo que es importante del fluido 
no es la velocidad en sí, o la presión, sino cuándo hay una turbulencia (rotación con mucha fuerza), y esto último se puede medir bien en los 
mapas del tiempo. \\

¿Qué es la vorticidad?:

\begin{definicion}[Vorticidad]
    Definimos $\omega$ como el rotacional del vector de velocidades $u$. Es decir
    $$\omega = \rot(u)$$
    donde\footnote{Consultar \ref{section:rot} si se quiere recordar más en profundidad.} 
    $$\rot(u) = \nabla \times u = \begin{vmatrix}
        i & \frac{\partial}{\partial x_1} & u_1(t, x_1, x_2) \\\\
        j & \frac{\partial}{\partial x_2} & u_2(t, x_1, x_2) \\\\
        k & \frac{\partial}{\partial x_3} & u_3(t, x_1, x_2)
    \end{vmatrix} = \left(0, 0, \overbrace{\dfrac{\partial u_2}{\partial x_1} - \dfrac{\partial u_1}{\partial x_2}}^{\omega} \right)$$
\end{definicion}

\begin{observacion}
    No tiene sentido, a priori, definirlo fuera de $\R^3$.
\end{observacion}

\begin{ejercicio} % TODO: Hacer para mañana
    Comprobar que la siguiente expresión es cierta:
    $$\rot((u \cdot \nabla)u) = (u \cdot \nabla) \omega + \di(u) \omega$$
    Podemos hacerlo componente a componente y sumar o como sea. \\

    Como el fluido es incompresible, $\di(u) \omega = 0$, y se simplifica la expresión anterior.
\end{ejercicio}

Otra pregunta es si podemos aplicar el rotacional de la primera ecuación de \eqref{eq:ec23}. Pues a priori no, pero vamos a hacerlo,
y luego analizamos cuándo tiene sentido esta expresión. \\

\begin{equation}
    \dfrac{\partial \omega}{\partial t} + (u \cdot \nabla) \omega + \rot(\nabla p) = 0
\end{equation}

\begin{observacion}
    En dimensión 3 aparece el siguiente término:
    $$(u \cdot \nabla) \omega - \boxed{(\omega \cdot \nabla) u} \di(u) \omega$$
    Este sustraendo señalado es el que representa la \textit{deformación de la vorticidad}, y es el culpable de muchas cosas que pasan
    en el caso tridimensional que en el bidimensional no ocurren (por ejemplo, el estudio de la regularidad de $\omega$).
\end{observacion}

Veamos cuánto vale $\rot(\nabla p)$: 

$$\rot(\nabla p) = \dfrac{\partial \dfrac{\partial p}{\partial x_2}}{\partial x_1} - \dfrac{\partial \dfrac{\partial p}{\partial x_1}}{\partial x_2} = 0$$

Por lo tanto tenemos hasta el momento la siguiente ecuación:

\begin{equation}\label{eq:ec25}
    \dfrac{\partial \omega}{\partial t} + (u \cdot \nabla) \omega = 0
\end{equation}

Siguen siendo tres incógnitas, pero hemos pasado de una ecuación vectorial a escalar ($\omega$ es el escalar). Otras ecuaciones que debemos considerar son
$$\omega = \rot(u)$$
$$\di(u) = 0$$
$$|u| \to 0 \quad |x| \to 0$$
$$\omega(t=0,x) = \omega_0(x)$$

No hemos simplificado gran cosa hasta ahora, pero nos disponemos a hacerlo. \\

Usaremos el siguiente lema de Poincaré:

\begin{lema}[de Poincaré]
    Sea $u$ un campo de velocidades en $\R^2$ como \eqref{eq:ec24}. Entonces se tiene que 
    $$\di(u) = 0 \iff \exists \psi : u = \rot(\psi)$$
\end{lema}

\begin{observacion}
    Para una formulación más técnica habría que estudiar $1$-formas (cosa que no vamos a hacer). La implicación hacia la derecha es sencilla,
    ya que basta aplicar rotacional. Así:
    $$\rot(\psi) = \left(
        \dfrac{\partial \psi}{\partial x_2}, - \dfrac{\partial \psi}{\partial x_1}, 0
    \right)$$
    Como $$\di\left( \dfrac{\partial \psi}{\partial x_2}, - \dfrac{\partial \psi}{\partial x_1} \right) = 
    \dfrac{\partial}{\partial x_1} \left( \dfrac{\partial \psi}{\partial x_2} \right) + 
    \dfrac{\partial}{\partial x_2} \left( - \dfrac{\partial \psi}{\partial x_1} \right) = 0$$
\end{observacion}

Tenemos entonces que $\omega = \rot(\rot(\psi))$\footnote{Hemos sustituido $u = \rot(\psi)$ en $\omega = \rot(u)$.}
y nos queda 
$$\omega = \left(0, 0, - \dfrac{\partial^2 \psi}{\partial x_1^2}, - \dfrac{\partial^2 \psi}{\partial x_2^2} \right)$$
o sea, $\omega = - \Delta \psi$ en $\R^2$.
A veces se representa también $\rot \psi$ como $\nabla^\perp \psi$.
Considerando la condición $|\nabla \psi| \to 0$ cuando $|x| \to \infty$, nos queda una ecuación como la estudiada en modelos.
La solución fundamental del laplaciano (vista en Modelos Matemáticos II) es
$$\Gamma_2 = G = - \dfrac{1}{2 \pi} \ln |x|$$
Lo bueno con esto es que para cualquier función solución de $-\Delta \psi = \omega$ la solución es 
$$\psi = G \ast \omega$$
donde $\ast$ representa el operador convolución. Y ahora podemos aplicar $u = \rot \psi$. Recordemos que cuando 
derivábamos una convolución podíamos derivar donde quisiéramos. El problema de la derivabilidad lo introduce el logaritmo. Suponiendo 
que podemos hacerlo, resulta que 
$$u = \rot \psi = \dfrac{1}{2 \pi} \left( - \dfrac{\dfrac{x_2}{|x|}}{|x|}, \dfrac{\dfrac{x_1}{|x|}}{|x|} \right) \ast \omega$$
y al producto de los primeros dos factores:
\begin{equation}\label{eq:ec26}
    \dfrac{1}{2 \pi} \left( - \dfrac{\dfrac{x_2}{|x|}}{|x|}, \dfrac{\dfrac{x_1}{|x|}}{|x|} \right) = \dfrac{(-x_2, x_1)}{2 \pi |x|^2} = K
\end{equation}
De esta manera $u = K \ast \omega$. \\

% Veremos que esta ecuación no es como la de Burgers.

Recordamos que en dimensión 2D, la ecuación es \eqref{eq:ec25}. Al pasar a 3D aparece otro sustraendo, que en $2D$ es 0:
$$\dfrac{\partial \omega}{\partial t} + (u \cdot \nabla) \omega - (\omega \cdot \nabla)u = 0$$
Se puede añadir incluso $f \rightarrow \nu \Delta \omega$(una difusión).
$$\dfrac{\partial \omega}{\partial t} = \nu \Delta \omega$$
Al añadir esa $f$ la ecuación se le denomina de Navier-Stokes. Respecto al problema del milenio, al principio se consideraba sin fuerza externa $f$,
pero hace 10-15 años se permitió demostrar también con fuerza externa (se consideró también válida una hipotética demostración). Aunque sin la $f$ 
también es Navier-Stokes. \\

En el caso $3D$ la solución fundamental vale 
$$\Gamma_3 = \dfrac{1}{\sigma_3 |x|}$$
donde $\sigma_3$ es el área de la esfera unidad en dimensión 3. 
$\Gamma_2$ es más singular (más fuerte) en $0$, que $\Gamma_3$. \\

Retomando entonces el valor de $K$ de \eqref{eq:ec26}, nos queda que 

\begin{equation}\label{eq:ec27}
    u(t,x) = \int_{\R^2} \dfrac{(-y_2, y_1)}{2 \pi |y|^2} \omega(t,x-y)\,dy
\end{equation}

\begin{observacion}
    Recordamos que la convolución es simétrica. Tiene que ver con el orden en que se escribe $x$ y $y$.
\end{observacion}

Esto es una integral singular. Toda una rama del análisis. A veces a la ecuación de Burgers le llaman la ecuación de los fluidos en una dimensión.
Quien dice eso no sabe lo que dice. La ecuación \eqref{eq:ec25} también se puede escribir como
\begin{equation}\label{eq:ec28}
    \dfrac{\partial \omega}{\partial t} + \overbrace{\di(u \omega)}^{(u \cdot \nabla) \omega} = 0
\end{equation}

porque $\di(u) = 0$ (al suponer que es incompresible). Sustituyendo $u = K \ast \omega$ en \eqref{eq:ec28}, nos queda 

\begin{equation}\label{eq:ec29}
    \dfrac{\partial \omega}{\partial t} + \di((K \ast \omega) \omega) = 0
\end{equation}

y podríamos pensar en comparar \eqref{eq:ec29} con la expresión de Burgers

$$\dfrac{\partial \omega}{\partial t} + \dfrac{\partial}{\partial x} \left( \dfrac{\omega^2}{2} \right) = 0$$

Sin embargo, la ecuación de Euler no desarrolla choque, a diferencia de la ecuación de Burgers, debido a la convolución\footnote{Que es como 
una especie de ``media'' alrededor del punto.}. \\

Consideremos $\chi_{\Omega}$ la función característica de $\Omega$, dada por 
\Func{\chi_\Omega}{\R^2}{\{0,1\}}{x}{\chi_\Omega(x) = \begin{cases} 1 \quad& x \in \Omega \\ 0 \quad& x \not\in \Omega \end{cases} }

Vemos gráficamente su uso:

\begin{figure}[H]
    \begin{center}
        \tikzset{every picture/.style={line width=0.75pt}} %set default line width to 0.75pt        

        \begin{tikzpicture}[x=0.75pt,y=0.75pt,yscale=-1,xscale=1]
        %uncomment if require: \path (0,300); %set diagram left start at 0, and has height of 300

        %Shape: Polygon Curved [id:ds3053046118097098] 
        \draw   (101,86) .. controls (121,76) and (154,117) .. (134,137) .. controls (114,157) and (43,107) .. (63,137) .. controls (83,167) and (-2,156) .. (17,103) .. controls (36,50) and (81,96) .. (101,86) -- cycle ;
        %Shape: Circle [id:dp45660914167393973] 
        \draw   (53,192) .. controls (53,178.19) and (64.19,167) .. (78,167) .. controls (91.81,167) and (103,178.19) .. (103,192) .. controls (103,205.81) and (91.81,217) .. (78,217) .. controls (64.19,217) and (53,205.81) .. (53,192) -- cycle ;
        %Curve Lines [id:da39259218054114486] 
        \draw    (168,87) .. controls (187.6,74.26) and (209.12,74.01) .. (233.5,85.29) ;
        \draw [shift={(235,86)}, rotate = 205.64] [color={rgb, 255:red, 0; green, 0; blue, 0 }  ][line width=0.75]    (10.93,-3.29) .. controls (6.95,-1.4) and (3.31,-0.3) .. (0,0) .. controls (3.31,0.3) and (6.95,1.4) .. (10.93,3.29)   ;
        %Shape: Polygon Curved [id:ds4134756710331108] 
        \draw   (335,86) .. controls (355,76) and (388,117) .. (368,137) .. controls (348,157) and (277,107) .. (297,137) .. controls (317,167) and (232,156) .. (251,103) .. controls (270,50) and (315,96) .. (335,86) -- cycle ;
        %Shape: Circle [id:dp8570779698638371] 
        \draw   (287,192) .. controls (287,178.19) and (298.19,167) .. (312,167) .. controls (325.81,167) and (337,178.19) .. (337,192) .. controls (337,205.81) and (325.81,217) .. (312,217) .. controls (298.19,217) and (287,205.81) .. (287,192) -- cycle ;
        %Straight Lines [id:da4669948322720253] 
        \draw    (431,131) -- (495,130.03) ;
        \draw [shift={(497,130)}, rotate = 179.13] [color={rgb, 255:red, 0; green, 0; blue, 0 }  ][line width=0.75]    (10.93,-3.29) .. controls (6.95,-1.4) and (3.31,-0.3) .. (0,0) .. controls (3.31,0.3) and (6.95,1.4) .. (10.93,3.29)   ;
        %Curve Lines [id:da8186342329158359] 
        \draw    (547,170) .. controls (532,148) and (531,94) .. (596,93) ;
        %Curve Lines [id:da5025705863547494] 
        \draw    (582,173) .. controls (572,161) and (555,163) .. (547,170) ;
        %Curve Lines [id:da6580378891015024] 
        \draw    (615,145) .. controls (615,160) and (597,174) .. (582,173) ;
        %Curve Lines [id:da9498702609506636] 
        \draw    (615,145) .. controls (619,119) and (580,142) .. (572,152) ;
        %Curve Lines [id:da9148030838417588] 
        \draw    (559,119) .. controls (541,138) and (558,168) .. (572,152) ;
        %Curve Lines [id:da7310709300556156] 
        \draw    (559,119) .. controls (570,105) and (575,102) .. (596,93) ;
        %Curve Lines [id:da255938325673975] 
        \draw    (296,189) .. controls (290.09,213.63) and (333.66,218.85) .. (326.37,181.73) ;
        \draw [shift={(326,180)}, rotate = 77.01] [color={rgb, 255:red, 0; green, 0; blue, 0 }  ][line width=0.75]    (10.93,-3.29) .. controls (6.95,-1.4) and (3.31,-0.3) .. (0,0) .. controls (3.31,0.3) and (6.95,1.4) .. (10.93,3.29)   ;

        % Text Node
        \draw (71,100) node [anchor=north west][inner sep=0.75pt]   [align=left] {$\displaystyle 1$};
        % Text Node
        \draw (69,49) node [anchor=north west][inner sep=0.75pt]   [align=left] {$\displaystyle 0$};
        % Text Node
        \draw (119,53) node [anchor=north west][inner sep=0.75pt]   [align=left] {$\displaystyle \omega _{0} \ =\ \chi _{\Omega }$};
        % Text Node
        \draw (152,121) node [anchor=north west][inner sep=0.75pt]   [align=left] {$\displaystyle \Omega $};
        % Text Node
        \draw (305,100) node [anchor=north west][inner sep=0.75pt]   [align=left] {$\displaystyle 1$};
        % Text Node
        \draw (303,49) node [anchor=north west][inner sep=0.75pt]   [align=left] {$\displaystyle 0$};
        % Text Node
        \draw (369,55) node [anchor=north west][inner sep=0.75pt]   [align=left] {$\displaystyle \omega ( t,x)$};
        % Text Node
        \draw (386,121) node [anchor=north west][inner sep=0.75pt]   [align=left] {$\displaystyle \Omega $};
        % Text Node
        \draw (433,86) node [anchor=north west][inner sep=0.75pt]   [align=left] {Tras cierto \\tiempo};


        \end{tikzpicture}
        \caption{Ejemplo de parches de vorticidad}
        \label{fig:f31}
    \end{center}
\end{figure}

Este tipo de datos iniciales de la vorticidad\footnote{Aquellos en los que la vorticidad es función característica de un conjunto $\Omega$.} se denominan parches (o paquetes) de vorticidad\footnote{En inglés, \textit{vortex patch}.}. 
Si fuera verdad la Figura \ref{fig:f31}, lo que nos interesa es el contorno. Buscar cómo evolucionan los contornos es una rama que se denomina dinámica de contornos.
En este caso tiene soporte compacto y es integrable. \\

Para formalizar más, el espacio en que estas soluciones interesantes viven es $\omega(t=0,x) = \omega_0(x) \in L^\infty(\R^2) \cap L^1(\R^2)$. Casi no 
haría falta que fueran integrables si le añadimos soporte compacto. \\

\begin{figure}[H]
    \begin{center}
        \tikzset{every picture/.style={line width=0.75pt}} %set default line width to 0.75pt        

        \begin{tikzpicture}[x=0.75pt,y=0.75pt,yscale=-1,xscale=1]
        %uncomment if require: \path (0,300); %set diagram left start at 0, and has height of 300

        %Straight Lines [id:da2812470710295212] 
        \draw    (31,51) -- (79,51) ;
        \draw [shift={(81,51)}, rotate = 180] [color={rgb, 255:red, 0; green, 0; blue, 0 }  ][line width=0.75]    (10.93,-3.29) .. controls (6.95,-1.4) and (3.31,-0.3) .. (0,0) .. controls (3.31,0.3) and (6.95,1.4) .. (10.93,3.29)   ;
        %Straight Lines [id:da47566060676251687] 
        \draw    (31,70) -- (79,70) ;
        \draw [shift={(81,70)}, rotate = 180] [color={rgb, 255:red, 0; green, 0; blue, 0 }  ][line width=0.75]    (10.93,-3.29) .. controls (6.95,-1.4) and (3.31,-0.3) .. (0,0) .. controls (3.31,0.3) and (6.95,1.4) .. (10.93,3.29)   ;
        %Straight Lines [id:da09091776510160365] 
        \draw    (30,90) -- (78,90) ;
        \draw [shift={(80,90)}, rotate = 180] [color={rgb, 255:red, 0; green, 0; blue, 0 }  ][line width=0.75]    (10.93,-3.29) .. controls (6.95,-1.4) and (3.31,-0.3) .. (0,0) .. controls (3.31,0.3) and (6.95,1.4) .. (10.93,3.29)   ;
        %Straight Lines [id:da15369811841719883] 
        \draw    (31,161) -- (79,161) ;
        \draw [shift={(81,161)}, rotate = 180] [color={rgb, 255:red, 0; green, 0; blue, 0 }  ][line width=0.75]    (10.93,-3.29) .. controls (6.95,-1.4) and (3.31,-0.3) .. (0,0) .. controls (3.31,0.3) and (6.95,1.4) .. (10.93,3.29)   ;
        %Straight Lines [id:da16898794784094828] 
        \draw    (31,180) -- (79,180) ;
        \draw [shift={(81,180)}, rotate = 180] [color={rgb, 255:red, 0; green, 0; blue, 0 }  ][line width=0.75]    (10.93,-3.29) .. controls (6.95,-1.4) and (3.31,-0.3) .. (0,0) .. controls (3.31,0.3) and (6.95,1.4) .. (10.93,3.29)   ;
        %Straight Lines [id:da4063777233796856] 
        \draw    (30,200) -- (78,200) ;
        \draw [shift={(80,200)}, rotate = 180] [color={rgb, 255:red, 0; green, 0; blue, 0 }  ][line width=0.75]    (10.93,-3.29) .. controls (6.95,-1.4) and (3.31,-0.3) .. (0,0) .. controls (3.31,0.3) and (6.95,1.4) .. (10.93,3.29)   ;
        %Straight Lines [id:da4020866803904283] 
        \draw    (140,90) -- (140,119) ;
        \draw [shift={(140,121)}, rotate = 270] [color={rgb, 255:red, 0; green, 0; blue, 0 }  ][line width=0.75]    (10.93,-3.29) .. controls (6.95,-1.4) and (3.31,-0.3) .. (0,0) .. controls (3.31,0.3) and (6.95,1.4) .. (10.93,3.29)   ;
        %Curve Lines [id:da15908134021767917] 
        \draw    (226,63) .. controls (252.2,49.6) and (304.2,57.6) .. (317,92) ;
        %Curve Lines [id:da21164709689584338] 
        \draw    (226,63) .. controls (209.2,72.6) and (218.2,85.6) .. (230.2,96.6) ;
        %Curve Lines [id:da8087665006399677] 
        \draw    (230.2,96.6) .. controls (230.2,81.6) and (294.2,61.6) .. (317,92) ;
        %Straight Lines [id:da9647454151853694] 
        \draw    (224.2,69.6) -- (234,82) ;
        %Straight Lines [id:da8782904375213573] 
        \draw    (236.2,65.6) -- (246,78) ;
        %Straight Lines [id:da09315412580354421] 
        \draw    (250.2,61.6) -- (260,74) ;
        %Straight Lines [id:da7099557171234606] 
        \draw    (267.2,61.6) -- (277,74) ;
        %Straight Lines [id:da32765479150297305] 
        \draw    (287.2,66.6) -- (296,78) ;
        %Straight Lines [id:da5742289166766529] 
        \draw    (227.2,54.6) -- (239.18,34.32) ;
        \draw [shift={(240.2,32.6)}, rotate = 120.58] [color={rgb, 255:red, 0; green, 0; blue, 0 }  ][line width=0.75]    (10.93,-3.29) .. controls (6.95,-1.4) and (3.31,-0.3) .. (0,0) .. controls (3.31,0.3) and (6.95,1.4) .. (10.93,3.29)   ;
        %Shape: Circle [id:dp9743827661546515] 
        \draw  [color={rgb, 255:red, 0; green, 0; blue, 0 }  ,draw opacity=1 ][fill={rgb, 255:red, 0; green, 0; blue, 0 }  ,fill opacity=1 ] (314.7,89.7) .. controls (314.7,88.43) and (315.73,87.4) .. (317,87.4) .. controls (318.27,87.4) and (319.3,88.43) .. (319.3,89.7) .. controls (319.3,90.97) and (318.27,92) .. (317,92) .. controls (315.73,92) and (314.7,90.97) .. (314.7,89.7) -- cycle ;
        %Straight Lines [id:da44093713674943735] 
        \draw    (320.2,85.6) -- (332.18,65.32) ;
        \draw [shift={(333.2,63.6)}, rotate = 120.58] [color={rgb, 255:red, 0; green, 0; blue, 0 }  ][line width=0.75]    (10.93,-3.29) .. controls (6.95,-1.4) and (3.31,-0.3) .. (0,0) .. controls (3.31,0.3) and (6.95,1.4) .. (10.93,3.29)   ;
        %Curve Lines [id:da9782551750351288] 
        \draw    (221.45,127.21) .. controls (250.85,125.95) and (294.78,154.9) .. (292.07,191.5) ;
        %Curve Lines [id:da6457406154474422] 
        \draw    (221.45,127.21) .. controls (202.17,128.93) and (204.93,144.5) .. (211.25,159.5) ;
        %Curve Lines [id:da6592127058421853] 
        \draw    (211.25,159.5) .. controls (217.51,145.86) and (284.02,154.37) .. (292.07,191.5) ;
        %Straight Lines [id:da8164395489389693] 
        \draw    (217.06,132.45) -- (220.8,147.81) ;
        %Straight Lines [id:da5997991788831862] 
        \draw    (229.63,133.82) -- (233.37,149.18) ;
        %Straight Lines [id:da21725399929697986] 
        \draw    (244.03,136.02) -- (247.76,151.38) ;
        %Straight Lines [id:da37594237267956354] 
        \draw    (259.48,143.11) -- (263.22,158.47) ;
        %Straight Lines [id:da5660065818913599] 
        \draw    (275.57,155.99) -- (278.82,170.02) ;
        %Curve Lines [id:da4348210794241574] 
        \draw    (160.2,132.6) .. controls (170.14,141.96) and (190.17,136.87) .. (205.99,118.07) ;
        \draw [shift={(207.2,116.6)}, rotate = 128.66] [color={rgb, 255:red, 0; green, 0; blue, 0 }  ][line width=0.75]    (10.93,-3.29) .. controls (6.95,-1.4) and (3.31,-0.3) .. (0,0) .. controls (3.31,0.3) and (6.95,1.4) .. (10.93,3.29)   ;
        %Curve Lines [id:da773331916689267] 
        \draw    (159.2,156.6) .. controls (171.02,152.66) and (183.81,145.81) .. (204.26,172.36) ;
        \draw [shift={(205.2,173.6)}, rotate = 233.13] [color={rgb, 255:red, 0; green, 0; blue, 0 }  ][line width=0.75]    (10.93,-3.29) .. controls (6.95,-1.4) and (3.31,-0.3) .. (0,0) .. controls (3.31,0.3) and (6.95,1.4) .. (10.93,3.29)   ;
        %Curve Lines [id:da6631647361637638] 
        \draw    (224.46,178.63) .. controls (236.65,199.82) and (249.91,153.53) .. (236.73,165.38) ;
        %Curve Lines [id:da6563990817342066] 
        \draw    (236.73,165.38) .. controls (234.1,168.32) and (234.18,173.34) .. (237.41,171.1) ;
        %Curve Lines [id:da1508397637449559] 
        \draw    (230.2,110.42) .. controls (246.34,88.9) and (265.31,134.82) .. (247.15,123.34) ;
        %Curve Lines [id:da5937672725257489] 
        \draw    (247.15,123.34) .. controls (243.52,120.47) and (243.52,115.45) .. (247.95,117.6) ;
        %Curve Lines [id:da9324690646615994] 
        \draw    (268.2,131.42) .. controls (284.34,109.9) and (303.31,155.82) .. (285.15,144.34) ;
        %Curve Lines [id:da9285295925185382] 
        \draw    (285.15,144.34) .. controls (281.52,141.47) and (281.52,136.45) .. (285.95,138.6) ;
        %Curve Lines [id:da800505442183486] 
        \draw    (303.2,166.6) .. controls (321.2,149.6) and (289.2,162.6) .. (304.2,142.6) ;
        %Curve Lines [id:da8942786587901006] 
        \draw    (323.2,182.6) .. controls (341.2,165.6) and (309.2,178.6) .. (324.2,158.6) ;
        %Straight Lines [id:da19147972215759101] 
        \draw    (275.2,118.6) -- (319.21,113.82) ;
        \draw [shift={(321.2,113.6)}, rotate = 173.8] [color={rgb, 255:red, 0; green, 0; blue, 0 }  ][line width=0.75]    (10.93,-3.29) .. controls (6.95,-1.4) and (3.31,-0.3) .. (0,0) .. controls (3.31,0.3) and (6.95,1.4) .. (10.93,3.29)   ;
        %Curve Lines [id:da7954650578120924] 
        \draw    (235.2,192.6) .. controls (253.2,184.6) and (268.2,185.6) .. (261.2,163.6) ;
        %Curve Lines [id:da9433565556338858] 
        \draw    (247.82,201.02) .. controls (265.82,193.02) and (280.82,194.02) .. (273.82,172.02) ;
        %Curve Lines [id:da48799610894319867] 
        \draw    (267.2,208.6) .. controls (307.2,178.6) and (260.2,230.6) .. (300.2,200.6) ;
        %Curve Lines [id:da2469256654584021] 
        \draw    (303.2,216.6) .. controls (343.2,186.6) and (277.2,253.6) .. (317.2,223.6) ;
        %Curve Lines [id:da31900950697469144] 
        \draw [color={rgb, 255:red, 0; green, 0; blue, 255 }  ,draw opacity=1 ]   (292.07,191.5) .. controls (325,187.6) and (327,211.6) .. (368,207.6) ;
        %Straight Lines [id:da7667827219254041] 
        \draw    (360,198) -- (403.45,162.86) ;
        \draw [shift={(405,161.6)}, rotate = 141.03] [color={rgb, 255:red, 0; green, 0; blue, 0 }  ][line width=0.75]    (10.93,-3.29) .. controls (6.95,-1.4) and (3.31,-0.3) .. (0,0) .. controls (3.31,0.3) and (6.95,1.4) .. (10.93,3.29)   ;
        %Straight Lines [id:da4619382821140623] 
        \draw    (348,220.6) -- (386.38,248.43) ;
        \draw [shift={(388,249.6)}, rotate = 215.94] [color={rgb, 255:red, 0; green, 0; blue, 0 }  ][line width=0.75]    (10.93,-3.29) .. controls (6.95,-1.4) and (3.31,-0.3) .. (0,0) .. controls (3.31,0.3) and (6.95,1.4) .. (10.93,3.29)   ;
        %Straight Lines [id:da4973385330210537] 
        \draw    (410,209.6) -- (507.02,195.88) ;
        \draw [shift={(509,195.6)}, rotate = 171.95] [color={rgb, 255:red, 0; green, 0; blue, 0 }  ][line width=0.75]    (10.93,-3.29) .. controls (6.95,-1.4) and (3.31,-0.3) .. (0,0) .. controls (3.31,0.3) and (6.95,1.4) .. (10.93,3.29)   ;
        %Curve Lines [id:da9763631526140586] 
        \draw [color={rgb, 255:red, 0; green, 0; blue, 255 }  ,draw opacity=1 ]   (326,194.6) .. controls (333.56,188.93) and (343.8,194.87) .. (352.5,200.94) ;
        \draw [shift={(354,202)}, rotate = 215.42] [color={rgb, 255:red, 0; green, 0; blue, 255 }  ,draw opacity=1 ][line width=0.75]    (10.93,-3.29) .. controls (6.95,-1.4) and (3.31,-0.3) .. (0,0) .. controls (3.31,0.3) and (6.95,1.4) .. (10.93,3.29)   ;

        % Text Node
        \draw (238,15) node [anchor=north west][inner sep=0.75pt]   [align=left] {Sólido};
        % Text Node
        \draw (337,43) node [anchor=north west][inner sep=0.75pt]   [align=left] {Pequeña singularidad.\\Pasan más cosas};
        % Text Node
        \draw (326,100) node [anchor=north west][inner sep=0.75pt]   [align=left] {Pequeñas \\turbulencias};
        % Text Node
        \draw (414,137) node [anchor=north west][inner sep=0.75pt]   [align=left] {Componente\\tangencial \\a $\displaystyle u$, $\displaystyle u\cdot \tau $};
        % Text Node
        \draw (373,198) node [anchor=north west][inner sep=0.75pt]   [align=left] {$\displaystyle \textcolor[rgb]{0,0,1}{\delta _{F}}$};
        % Text Node
        \draw (395,228) node [anchor=north west][inner sep=0.75pt]   [align=left] {Componente normal\\a $\displaystyle u$, $\displaystyle u\cdot n$, no tiene por \\qué ser regular};
        % Text Node
        \draw (519,162) node [anchor=north west][inner sep=0.75pt]   [align=left] {Superficie si \\Dim 3. Línea si \\Dim 2};
        % Text Node
        \draw (194,244) node [anchor=north west][inner sep=0.75pt]   [align=left] {\textit{Vortex Sheet}};


        \end{tikzpicture}
    \caption{Formación de una hoja de vorticidad.}
    \label{fig:f34}
    \end{center}
\end{figure}

Este tipo de fenómenos se dan cuando, por ejemplo, despega un avión. Si no se controla la hoja de vórtice\footnote{También conocida como lámina o capa de vórtices. En inglés como \textit{vortex sheet}.}, 
tiene tanta fuerza que parte el ala del avión (da igual que esté hecho de acero). El $\delta_F$ que aparece en la Figura \ref{fig:f34} es una delta de Dirac, que veremos más adelante. \\

Estos ejemplos motivan el estudio de la regularidad de la velocidad y de la vorticidad. En particular, nos interesa determinar bajo qué hipótesis sobre 
$\omega$ podemos garantizar que el campo de velocidades $(u=K*\omega)$ sea acotado y continuo. Volviendo a las ecuaciones, $u = K \ast \omega$ no depende del tiempo. 
Burgers no regulariza. Esta ecuación ($u = K \ast \omega$) es de tipo hiperbólico\footnote{Vista en Modelos Matemáticos II.}. La solución no se espera que sea continua. \\

Veamos el siguiente resultado\footnote{De momento paramos el tiempo, y luego estudiamos qué pasa con él.}:

\begin{lema}
    ·Si $\omega \in L^1(\R^2) \cap L^\infty(\R^2) \Longrightarrow u \in L^\infty$. La segunda parte nos dice que, no solo es acotada, sino que es continua.
    Además, $\norma{u}_{L^\infty} \leq \norma{\omega}_{L^1}^{1/2} \norma{\omega}_{L^\infty}^{1/2}$.
\end{lema}

\begin{proof}
    Consideramos 
    $$u(x) = \int_{\mathbb{R}^2} K(x-y)\omega(y)\,dy =
    \int\limits_{\substack{y\in\mathbb{R}^2\\|x-y|\leq d }} K(x-y)\omega(y)\,dy +
    \int\limits_{\substack{y\in\mathbb{R}^2\\|x-y| > d }} K(x-y)\omega(y)\,dy = I_1 + I_2$$
    donde $d \in \R$ es arbitrario. Como queremos ver que está acotado:
    $$|u(x)| \leq
    \int\limits_{\substack{y\in\mathbb{R}^2\\|x-y|\leq d}} |K(x-y)\omega(y)|\,dy + 
    \int\limits_{\substack{y\in\mathbb{R}^2\\|x-y| > d}} |K(x-y)\omega(y)|\,dy$$
    Tenemos que 
    $$I_1 = \int_{|x-y| \leq d} |K(x-y) \omega(y)|\,dy \leq \int_{|x-y| \leq d} \dfrac{|\omega(y)|}{2 \pi |x-y|}\,dy \stackrel{(*)}{\leq} 
    \dfrac{\norma{\omega}_{L^\infty}}{2 \pi} \int_{|x-y| \leq d} \dfrac{1}{|x-y|}\,dy$$
    donde en $(*)$ se ha usado que $\omega$ está uniformemente acotada. Para resolver la última integral que nos ha quedado pasamos a coordenadas polares.
    $$(x-y)_1 = r \cos \theta \quad (x-y)_2 = r \sen \theta$$
    y
    $$\dfrac{\norma{\omega}_{L^\infty}}{2 \pi} \int_{0}^{2 \pi} \int_{0}^{d} \dfrac{1}{r} r \,dr\,d \theta = d \norma{\omega}_{L^\infty}$$
    Veamos ahora la segunda integral:
    $$I_2 = \int_{|x-y| > d} \dfrac{|\omega(y)|}{2 \pi |x-y|}\,dy \leq \dfrac{1}{2 \pi d} \int_{|x-y| > d} |\omega(y)|\,dy \leq 
    \dfrac{\norma{\omega}_{L^1}}{2 \pi d}$$
    Juntando todo lo anterior, resulta que
    $$|u(x)| \leq d \norma{\omega}_{L^\infty} + \dfrac{\norma{\omega}_{L^1}}{2 \pi d}$$
    Y ahora buscamos el valor óptimo de $d$. Para ello, consideramos la función
    $$f(d) = d \norma{\omega}_{L^\infty} + \dfrac{\norma{\omega}_{L^1}}{2 \pi d}$$
    y estudiamos su monotonía. Como
    $$f'(d) = \norma{\omega}_{L^\infty} - \dfrac{\norma{\omega}_{L^1}}{2 \pi d^2}$$
    $$f'' > 0$$
    deducimos que 
    $$d = \sqrt{\dfrac{\norma{\omega}_{L^1}}{2 \pi \norma{\omega}_{L^\infty}}}$$
    es el valor óptimo. Por tanto,
    $$|u(x)| \leq d \norma{\omega}_{L^\infty} + \dfrac{\norma{\omega}_{L^1}}{2 \pi d} = $$
    $$\dfrac{1}{\sqrt{2 \pi}} \norma{\omega}_{L^1}^{1/2} \norma{\omega}_{L^\infty}^{1/2} + 
    \dfrac{1}{\sqrt{2 \pi}} \norma{\omega}_{L^1}^{1/2} \norma{\omega}_{L^\infty}^{1/2} = \dfrac{2}{\sqrt{2 \pi}} \norma{\omega}_{L^1}^{1/2} \norma{\omega}_{L^\infty}^{1/2} $$ 
\end{proof}

\begin{ejercicio}
    ¿Qué ocurre si $\omega \in L^p(\R^2) \cap L^q(\R^2)$?, con $1 \leq p < 2 \leq q \leq +\infty$. Demostrar que $u \in L^\infty$ ($u$ está acotada). Se hace igual, 
    pero en vez de hacer directamente acotación, hay que aplicar Hölder primero. 
\end{ejercicio}

% 26.10.2026

Tenemos también el siguiente lema:

\begin{lema}\label{lema:l1}
    Si $\omega \in L^1 \cap L^\infty$, entonces:
    $$\begin{cases}
        1)& \quad |u(t,x+y) - u(t,x)| \leq C |y| \{1 + |\ln|y||\} \\
        2)& \quad |u(t,x)| \to 0 \quad |x| \to \infty
    \end{cases}$$
    donde $C = (C(\norma{\omega}_{L^1}, \norma{\omega}_{L^\infty}))$
\end{lema}

\begin{observacion}
    La segunda propiedad se deduce de que $u = K \ast \omega$. \\
\end{observacion}

Para estudiar el lema anterior vemos primero algunas propiedades de la función $K$:

\begin{ejercicio}\label{ejercicio:ej31}
    Demostrar que 
    $$|K(a) - K(b)| \leq C \left\{ \dfrac{1}{|a|} + \dfrac{1}{|b|} \right\} \quad \forall a,b \in \R^2 \setminus \{0\}$$
    $$|K(a) - K(b)| = \dfrac{1}{2 \pi} \dfrac{|b-a|}{|a||b|}$$

    La primera propiedad es sencilla de comprobar, y la segunda se deducirá a partir de la propia definición de $K$:
    $$K(x) = \dfrac{(-x_2, x_1)}{2 \pi |x|^2}$$
\end{ejercicio}

Demostramos el Lema \ref{lema:l1}:

\begin{proof}
    Para la primera propiedad, vemos que 
    $$u(x+y) - u(x) = \int_{\R^2} (K(x+y-z) - K(x-z)) \omega(z)\,dz \Longrightarrow$$
    $$|u(x+y) - u(x)| \leq \int_{\R^2} |(K(x+y-z) - K(x-z)) \omega(z)|\,dz \leq \int_{|x-z| \leq d} + \int_{|x-z| > d} = I_1 + I_2$$
    Hacemos ya la elección $d = 2|y| << 1$. Usamos ahora la primera propiedad del Ejercicio \ref{ejercicio:ej31} para $I_1$, y nos queda
    $$I_1 \leq \dfrac{1}{2 \pi} \int_{|x-z| \leq d} \dfrac{|\omega(z)|}{|x+y-z|} + \dfrac{1}{2 \pi} \int_{|x-z| \leq d} \dfrac{|\omega(z)|}{|x-z|}$$
    Estudiamos el denominador de la primera integral que acota a $I_1$:
    $$|x+y-z| \stackrel{|x-z| \leq d}{\leq} |x-z| + |y| \leq d + \dfrac{d}{2} = \dfrac{3d}{2}$$
    y se tiene que 
    $$\dfrac{1}{2 \pi} \int_{|x-z| \leq d} \dfrac{|\omega(z)|}{|x+y-z|} + \dfrac{1}{2 \pi} \int_{|x-z| \leq d} \dfrac{|\omega(z)|}{|x-z|} \leq $$
    $$\dfrac{1}{2 \pi} \int_{|x+y-z| \leq 3d/2} \dfrac{|\omega(z)|}{|x+y-z|} + \dfrac{1}{2 \pi} \int_{|x-z| \leq d} \dfrac{|\omega(z)|}{|x-z|}$$
    Ahora aplicamos los siguientes cambios a polares:
    $$(x+y-z)_i = r \cos \theta \quad (x+y-z)_j = r \sen \theta$$
    $$(x-z)_i = r \cos \theta \quad (x-z)_j = r \sen \theta$$
    obteniendo que 
    $$\dfrac{1}{2 \pi} \int_{|x+y-z| \leq 3d/2} \dfrac{|\omega(z)|}{|x+y-z|} + \dfrac{1}{2 \pi} \int_{|x-z| \leq d} \dfrac{|\omega(z)|}{|x-z|} \leq 
    C d \norma{\omega}_{L^\infty}$$
    Repetimos este esquema con $I_2$:
    $$I_2 \leq \dfrac{1}{2 \pi} \int_{|x-z| > d} \dfrac{|y| |\omega(z)|}{|x+y-z| |x-z|} \stackrel{(*)}{=} \int_{d \leq |x-z| \leq 1} + \int_{|x-z| > 1}$$
    donde en $(*)$ se ha usado que esta división en dos partes de la integral es la que funciona.
    En este caso el denominador $|x+y-z|$ está acotado inferiormente:
    $$|x+y-z| \geq |x-z| - |y| = |x-z| - \dfrac{d}{2} \geq |x-2| - \dfrac{|x-2|}{2} = \dfrac{|x-2|}{2}$$
    de donde
    $$\int_{1 \leq |x-z| \leq d} + \int_{|x-z| > 1} \leq \dfrac{|y|}{2 \pi} \int_{d \leq |x-z| \leq 1} \dfrac{2 |\omega(z)|}{|x-z|^2} + \int_{|x-z| > 1}$$
    Integramos en cada uno de los dos trozos. En el tramo $d \leq |x-z| \leq 1$:
    $$\dfrac{|y|}{\pi} \norma{\omega}_{L^\infty} \dfrac{1}{|x-2|^2} = \dfrac{|y|}{\pi} \norma{\omega}_{L^\infty} \int_{0}^{2 \pi} \int_{d}^{1} \dfrac{1}{r^2}
    r\,dr\,d \theta = 2 |y|\norma{\omega}_{L^\infty}(-\ln \overbrace{d}^{2 |y|}) \stackrel{d < 1}{=} $$
    $$2 |y| \norma{\omega}_{L^\infty}(|\ln d|) = 2 |y| \norma{\omega}_{L^\infty} (\ln 2 + |\ln|y||)$$
    Acotando la suma anterior de integrales, nos queda:
    $$\dfrac{|y|}{2 \pi} \int_{d \leq |x-z| \leq 1} \dfrac{2 |\omega(z)|}{|x-z|^2} + \int_{|x-z| > 1} \leq \dfrac{|y|}{\pi} \norma{\omega}_{L^1}$$
    y usando entonces las cotas que tenemos para $I_1$ e $I_2$, se deduce que 
    $$|u(x+y) - u(x)| \leq \int_{\R^2} |(K(x+y-z) - K(x-z)) \omega(z)|\,dz \leq I_1 + I_2 \leq$$
    $$|y| \left\{2 C \norma{\omega}_{L^\infty} + \dfrac{\norma{\omega}_{L^1}}{\pi} + 2 \ln 2 \norma{\omega}_{L^\infty} + 
    2 \norma{\omega}_{L^\infty} |\ln|y|| \right\}$$
\end{proof}

\begin{observacion}
    Justo el hecho de que se tenga $L^1$ y $L^\infty$ nos da la unicidad. Si buscamos exámenes de años anteriores, esto se puede preguntar en el examen. 
    Siempre intenta preguntar alguna cosa de teoría y algo que haya que hacer alguna cuenta. No hay mucho que preguntar. \\
\end{observacion}

Ahora, consideramos $f \in L^1$, y, dado $\varepsilon > 0$, existe $R > 0$ con\footnote{Se lee como ``$R$ suficientemente grande''.} $R >> 1$
tal que $$\int_{|y| > R} |f(y)|\,dy < \varepsilon$$
Esto vale también para casi cualquier $L^p$, ya que vale para $1 < p < \infty$, y entonces se tiene que
$$\int_{|y| > R} |f(y)|^p\,dy < \varepsilon$$
Para una integral lo suficientemente grande, tienen tendencia a irse a $0$ para un radio lo suficientemente grande. \\

Para comprobar esto, podemos verlo como
$$\int_{|y| > R} |f(y)|\,dy = \int_{\R^2} |f(y)|\,dy - \int_{|y| < R} |f(y)|\,dy \to 0 \quad (R \to \infty)$$
Consideramos ahora $|x| > 2R$. Vemos que 
$$|u(x)| \leq \int \dfrac{|\omega(y)|}{2 \pi |x-y|}\,dy = \int_{|y| > R} + \int_{|y| < R}$$
Ahora, si $|y| < R < |x| / 2$, donde la última desigualdad se debe a que $|x| < 2R$ por hipótesis, entonces
$$|x-y| \geq |x| - |y| = |x| \left\{ 1 - \dfrac{|y|}{|x|} \right\} > \dfrac{|x|}{2}$$
y
$$\dfrac{1}{2 \pi} \int_{|y| < R} \dfrac{|\omega(y)|}{|x-y|} \leq \dfrac{1}{2 \pi} \dfrac{2}{|x|} \int_{|y| < R} |\omega(y)|$$
Nuevamente, dividimos
$$
\int_{|y|>R}
=
\int_{\substack{|y|>R\\|x-y|\leq d}}
+
\underbrace{
\int_{\substack{|y|>R\\|x-y|>d}}
\frac{|\omega(y)|}{2\pi|x-y|}
}_{(*)}
$$
De esta manera, recuperando lo anterior:
$$|u(x)| \leq \int \dfrac{|\omega(y)|}{2 \pi |x-y|}\,dy = \int_{|y| > R} + \int_{|y| < R} = \dfrac{1}{\pi} \dfrac{1}{|x|} \norma{\omega_{L^1}} + 
\overbrace{d}^{\sqrt{\varepsilon}} \norma{\omega}_{L^\infty} + \dfrac{1}{2 \pi} \dfrac{\varepsilon}{d}$$
La integral $(*)$ nos da:
$$\dfrac{1}{2 \pi} \dfrac{1}{d} \int_{|y| > R} |\omega(y)| < \varepsilon$$
y no menor que $\norma{\omega}_{L^1}$, ya que se nos iría a infinito. \\

El siguiente resultado nos resultará bastante útil en lo que sigue:
\begin{lema}
    Sean $\varphi, \psi$ dos cambios de variable en $\R^2$ con Jacobianos inversos acotados, $\omega \in L^1 \cap L^\infty$ y 
    $$\left| \int_{\R^2} (K(x - \psi(y)) - K(x - \varphi(y)))\omega(y) \right| \leq C \norma{J_\psi^{-1}}_{L^\infty}
    \norma{J_\varphi^{-1}}_{L^\infty} \norma{\psi - \varphi}_{L^\infty} \{1 + |\ln \norma{\psi - \varphi}_{L^\infty}|\}$$
    Como sugerencia estudiar el caso $\varphi = 1$.
\end{lema}

\begin{observacion}
    Este resultado es una generalización de la primera propiedad del Lema \ref{lema:l1}.
\end{observacion}

% 27.10.2026

Vamos a estudiar a continuación las siguientes ecuaciones:
\begin{equation}\label{eq:ec30}
    \begin{cases}
        \dfrac{\partial \omega}{\partial t} + (u \cdot \nabla) \omega = 0 \\
        u = K \ast \omega \\
        \omega(t=0,x) = \omega_0(x) \in L^1 \cap L^\infty
    \end{cases}
\end{equation}

Y los siguientes pasos:
\begin{enumerate}
    \item Regularización (linealización), $\varepsilon \to (u^\varepsilon, \omega^\varepsilon)$ (soluciones de un problema aproximado).
    \item Estimaciones ``a priori''.
    \item Compacidad $(u^\varepsilon, \omega^\varepsilon) \stackrel{?}{\to} (u,\omega)$.
    \item $(u, \omega)$ verifique \eqref{eq:ec30}.
\end{enumerate}

Así pues, consideramos $\varepsilon > 0$, y una función 
$$\rho \in C_0^\infty(\R^2) = \begin{cases}
    c e^{\dfrac{1}{|x|^2 - 1}} \quad& |x| \leq 1 \\
    0 \quad& |x| > 1
\end{cases}; \quad c = \left( \int_{B(0,1)} e^{\dfrac{1}{|x|^2 - 1}}  \right)^{-1} \Longrightarrow \int \rho = 1$$

Con estas funciones vamos a construir las funciones
$$C_0^\infty(B(0,\varepsilon)) \ni \rho_\varepsilon = \varepsilon^{-2} \rho \left( \dfrac{x}{\varepsilon} \right) \Longrightarrow \omega_0 \ast \rho_\varepsilon = \omega_0^\varepsilon = \omega^0$$
Gráficamente:
\begin{figure}[H]
    \begin{center}
        \tikzset{every picture/.style={line width=0.75pt}} %set default line width to 0.75pt        

        \begin{tikzpicture}[x=0.75pt,y=0.75pt,yscale=-1,xscale=1]
        %uncomment if require: \path (0,300); %set diagram left start at 0, and has height of 300

        %Straight Lines [id:da44261844071929823] 
        \draw    (179,160) -- (411,160) ;
        %Straight Lines [id:da6567160598379451] 
        \draw    (290,149) -- (290,172) ;
        %Straight Lines [id:da7983049400227261] 
        \draw    (320,149) -- (320,172) ;
        %Straight Lines [id:da22603462141458197] 
        \draw    (260,150) -- (260,173) ;
        %Straight Lines [id:da19826657884532428] 
        \draw    (230,148) -- (230,171) ;
        %Straight Lines [id:da21094039468564174] 
        \draw    (350,148) -- (350,171) ;
        %Curve Lines [id:da1863316553924157] 
        \draw    (195,160) .. controls (260,160) and (242,80) .. (291,79) ;
        %Curve Lines [id:da3633895598898311] 
        \draw    (381,160) .. controls (320,160) and (342,80) .. (291,79) ;
        %Curve Lines [id:da9557797216212501] 
        \draw [color={rgb, 255:red, 0; green, 166; blue, 255 }  ,draw opacity=1 ]   (260,160) .. controls (280.97,160) and (275.16,28.64) .. (290.97,27) ;
        %Curve Lines [id:da5860357687341343] 
        \draw [color={rgb, 255:red, 0; green, 166; blue, 255 }  ,draw opacity=1 ]   (320,160) .. controls (300.32,160) and (307.42,28.64) .. (290.97,27) ;
        %Curve Lines [id:da4852955872046737] 
        \draw [color={rgb, 255:red, 255; green, 255; blue, 0 }  ,draw opacity=1 ]   (229,160) .. controls (270.94,160) and (259.32,64.2) .. (290.94,63) ;
        %Curve Lines [id:da8261174027182207] 
        \draw [color={rgb, 255:red, 255; green, 255; blue, 0 }  ,draw opacity=1 ]   (349,160) .. controls (309.65,160) and (323.84,64.2) .. (290.94,63) ;

        % Text Node
        \draw (285,173) node [anchor=north west][inner sep=0.75pt]   [align=left] {$\displaystyle 0$};
        % Text Node
        \draw (314,176) node [anchor=north west][inner sep=0.75pt]   [align=left] {$\displaystyle \varepsilon $};
        % Text Node
        \draw (246,173) node [anchor=north west][inner sep=0.75pt]   [align=left] {$\displaystyle -\varepsilon $};
        % Text Node
        \draw (211,172) node [anchor=north west][inner sep=0.75pt]   [align=left] {$\displaystyle -2\varepsilon $};
        % Text Node
        \draw (341,173) node [anchor=north west][inner sep=0.75pt]   [align=left] {$\displaystyle 2\varepsilon $};


        \end{tikzpicture}
    \end{center}
\end{figure}
Respecto a la definición de $\rho_\varepsilon$, tenemos que 
$$\varphi \in C_0^0 \Longrightarrow \prodescalar{\rho_\varepsilon}{\varphi} = \int \rho_\varepsilon \varphi \Longrightarrow \lim\limits_{\varepsilon \to 0} \int_{\R^2} \rho_\varepsilon(x) \varphi(x)\,dx = \varphi(0) = \prodescalar{\delta_{x = 0}}{\varphi}$$
donde $\delta$ es la delta de Dirac\footnote{La delta de Dirac es una distribución caracterizada por $\langle \delta_0,\varphi\rangle=\varphi(0)$ para toda función test $\varphi$.}.
Sabemos que 
$$\norma{\omega_0^\varepsilon - \omega_0}_{L^p} \to 0 \quad 1 \leq p < +\infty$$
Si consideramos el intervalo $[0, \varepsilon]$\footnote{Tiempo muy pequeño.}, tenemos el siguiente problema:
\begin{equation}
    (P_1) \begin{cases}
        \dfrac{\partial \omega^1}{\partial t} + (u^0 \cdot \nabla) \omega^1 = 0 \\
        u^0 = K \ast \omega^0 \text{ (dado)}\\
        \omega^1(t=0,x) = \omega^0(x)
    \end{cases}
\end{equation}

Lo importante es que $u^0$ no depende del tiempo, y las funciones con superíndice $0$ están dadas. La velocidad es la convolución con la vorticidad.
Lo único que hemos hecho es retrasar la velocidad un poco en el tiempo. Si esto funciona, la idea es continuar en los siguientes pasos. Ahora consideramos
el intervalo $[\varepsilon, 2 \varepsilon]$, y tenemos en este caso el problema:
\begin{equation}
    (P_2) \begin{cases}
        \dfrac{\partial \omega^2}{\partial t} + (u^1 \cdot \nabla) \omega^2 = 0 \\
        u^1 = K \ast \omega^1(\varepsilon, \cdot)\\
        \omega^2(\varepsilon,x) = \omega^1(\varepsilon, x)
    \end{cases}
\end{equation}

Repitiendo hasta $n$ veces, y situándonos entonces en el intervalo $[(n-1)\varepsilon, n \varepsilon]$ llegamos al problema:
\begin{equation}
    (P_n) \begin{cases}
        \dfrac{\partial \omega^n}{\partial t} + (u^{n-1} \cdot \nabla) \omega^n = 0 \\
        u^{n-1} = K \ast \omega^{n-1}((n-1)\varepsilon, \cdot)\\
        \omega^n((n-1)\varepsilon,x) = \omega^{n-1}((n-1)\varepsilon, x)
    \end{cases}
\end{equation}

De aquí llegamos a las siguientes ecuaciones:
\begin{equation}\label{eq:ec31}
    \begin{cases}
        \dfrac{d X^0(t)}{dt} = u^0(X^0(t)) \\
        X^0(0) = x
    \end{cases} \quad t \in [0,\varepsilon]
\end{equation}

Ahora, dado que $\omega_0 \in L^1 \cap L^\infty$, podemos ver si en norma está acotado:
$$\norma{\omega^0}_{L^1} = \norma{\omega_0 \ast \rho_\varepsilon}_{L^1} \leq \norma{\omega_0}_{L^1} \overbrace{\norma{\rho_\varepsilon}_{L^1}}^{1}$$
$$\norma{\omega^0}_{L^\infty} = \norma{\omega_0 \ast \rho_\varepsilon}_{L^\infty} \leq \norma{\omega_0}_{L^\infty} 
\overbrace{\norma{\rho_\varepsilon}_{L^1}}^{1}$$
Por lo tanto, $\omega^0 \in L^1 \cap L^\infty$ independientemente de $\varepsilon$. \\

Esto lo resumimos en el siguiente lema:
\begin{lema}
    Sea $\omega^0 \in L^1 \cap L^\infty$. Entonces se tiene que 
    \begin{equation*}
        \begin{cases}
            u^0 \in L^\infty \\
            |u^0(x+y) - u^0(x)| \leq C |y| \{1 + |\ln|y||\} \\
            |u^0(x)| \to 0 \quad |x| \to \infty \\
            \di(u^0) = 0
        \end{cases} \quad C = C(\norma{\omega^0}_{L^1}, \norma{\omega^0}_{L^\infty})
    \end{equation*}
    Y además, existe una única solución $X^0$ de \eqref{eq:ec31}.
\end{lema}

\begin{observacion}
    $\di(u^0) = 0$, porque $K$ es un rotacional, y vimos en el Ejercicio \ref{ejercicio:ej12} apartado (i) que la divergencia del rotacional es $0$.
\end{observacion}

\begin{proof}
    Vemos que 
    $$\dfrac{d}{dt} \omega^1(t, X^0(t)) = \left( \dfrac{\partial \omega^1}{\partial t} + \left( \dfrac{\partial X^0}{\partial t} \cdot \nabla \right)
    \omega^1 \right)(t, X^0(t)) = 0$$
    Y ya tenemos la solución de este problema, pues 
    $$\omega^1(t, X^0(t)) = \omega^1(0, X^0(0)) = \omega^0(x)$$
\end{proof}

\begin{observacion}
    Para cada $\varepsilon$ fijo, $u^0$ ya es infinitamente derivable en \eqref{eq:ec30}. Además, $X$ es un difeomorfismo, y podemos hacer un 
    cambio de variable si fuera necesario. Dicho cambio sería el siguiente:
    $$y = X^0(t; 0, x) = X^0(t) \Longrightarrow x = X^0(0; t,y)$$
    Así, aplicando el cambio:
    $$\omega^1(t,y) = \omega^0(X^0(0;t,y))$$
\end{observacion}

Lo anterior vale para un tiempo muy pequeño $[0, \varepsilon]$. Ahora nos preguntamos por el valor de la norma de $\omega^1$, 
con objeto de seguir estudiando propiedades:
$$\norma{\omega^1(t, \cdot)}_{L^1} = \norma{\omega^0(X^0(0; t, y))}_{L^1} = \int_{\R^2} |\omega^0(X^0(0;t,y))|\,dy = $$
$$\int_{\R^2} |\omega^0(z)| \overbrace{J_{x^0}}^{1}\,dz = \norma{\omega^0}_{L^1}$$
Lo anterior vale para la norma en $L^1$, pero en $L^\infty$ es análogo, pues:
$$\norma{\omega^1(t, \cdot)}_{L^\infty} = \sup_{y \in \R^2} |\omega^0(X^0(0;t,y))| = \sup_{z \in \R^2} |\omega^0(z)| = \norma{\omega^0}_{L^\infty}$$
es decir, las normas se preservan. \\

Este procedimiento además se puede implementar con numérico. En realidad lo que hay que resolver es \eqref{eq:ec30}. Lo demás viene como 
consecuencia de esta resolución. Podemos en este punto enunciar el siguiente resultado:

\begin{notacion}
    Dado el par $(u^\varepsilon, \omega^\varepsilon)(t,x)$, si $t \in [(n-1)\varepsilon, n \varepsilon]$, con $n = 1, \ldots, [T/\varepsilon] = N$,
    donde $[\cdot]$ denota a la parte entera, entonces 
    $$(u^\varepsilon, \omega^\varepsilon) \stackrel{not}{\equiv} (u^{n-1}, \omega^n)$$
    Asimismo, denotamos 
    \begin{equation*}
        \begin{cases}
            \dfrac{\partial \omega^\varepsilon}{\partial t} + (u^\varepsilon \cdot \nabla) \omega^\varepsilon = 0 \\ 
            u^\varepsilon = K \ast \omega^\varepsilon
        \end{cases}
        \iff 
        \begin{cases}
            \dfrac{\partial \omega^n}{\partial t} + (u^{n-1} \cdot \nabla) \omega^n = 0 \\ 
            u^{n-1} = K \ast \omega^{n}(n\varepsilon)
        \end{cases}
    \end{equation*}
\end{notacion}

\begin{observacion}
    No ponemos la condición inicial porque se va obteniendo iterativamente, pero la primera condición inicial es $\omega^\varepsilon(t=0,x) =
    \omega^0(x)$, y entonces tenemos:
    \begin{equation*}
        \begin{cases}
            \dfrac{\partial \omega^\varepsilon}{\partial t} + (u^\varepsilon \cdot \nabla) \omega^\varepsilon = 0 \\ 
            u^\varepsilon = K \ast \omega^\varepsilon \\
            \omega^\varepsilon(t=0,x) = \omega^0(x)
        \end{cases}
    \end{equation*}
    y los pares $\{(u^\varepsilon, \omega^\varepsilon)\}_{\varepsilon > 0}$.
\end{observacion}

Tenemos el siguiente lema:

\begin{lema}\label{lema:l31}
    Sea $\omega_0 \in L^1 \cap L^\infty$. Se verifica que existe $C (\in \R)$ independiente de $\varepsilon$ tal que 
    \begin{enumerate}
        \item $\norma{u^\varepsilon}_{L^\infty} \leq C$
        \item $$\norma{\omega^\varepsilon(t, \cdot)}_{L^1} = \norma{\omega^0}_{L^1} \quad (\norma{\omega^{n+1}(t, \cdot)}_{L^1} = \norma{\omega^n}_{L^1})$$
        $$\norma{\omega^\varepsilon(t, \cdot)}_{L^\infty} = \norma{\omega^0}_{L^\infty} \quad (\norma{\omega^{n+1}(t, \cdot)}_{L^\infty} = \norma{\omega^n}_{L^\infty})$$
        \item $|u^\varepsilon(x) - u^\varepsilon(y)| \leq C|x-y|\{1 + |\ln|x-y||\}$
        \item $|u^p(x) - u^q(x)| \leq C |p \varepsilon - q \varepsilon| \{1 + |\ln|p \varepsilon - q \varepsilon||\}$
    \end{enumerate}
\end{lema}

\begin{observacion}
    En el punto 4., al poner $p$ como superíndice, estamos definiendo $u^p$ en el intervalo $[(p-2)\varepsilon, (p-1)\varepsilon]$, y análogo con $q$.
\end{observacion}

La Figura \ref{fig:f32} nos muestra que en cada intervalo $\left[(n-1)\varepsilon, n\varepsilon \right]$, el campo de velocidades aproximado $u^\varepsilon$ depende de la variable espacial, pero permanece constante respecto del tiempo.

\begin{figure}[H]
    \begin{center}
        \tikzset{every picture/.style={line width=0.75pt}} %set default line width to 0.75pt        

        \begin{tikzpicture}[x=0.75pt,y=0.75pt,yscale=-1,xscale=1]
        %uncomment if require: \path (0,300); %set diagram left start at 0, and has height of 300

        %Straight Lines [id:da44261844071929823] 
        \draw    (150,160) -- (452,160) ;
        %Straight Lines [id:da6567160598379451] 
        \draw    (161,149) -- (161,172) ;
        %Straight Lines [id:da7983049400227261] 
        \draw    (201,149) -- (201,172) ;
        %Straight Lines [id:da21094039468564174] 
        \draw    (241,148) -- (241,171) ;
        %Straight Lines [id:da6065593239600682] 
        \draw    (351,149) -- (351,172) ;
        %Straight Lines [id:da8139488212809114] 
        \draw    (391,149) -- (391,172) ;
        %Straight Lines [id:da33702440515964105] 
        \draw    (430,148) -- (430,171) ;
        %Straight Lines [id:da8105430731321772] 
        \draw    (161,130) -- (201,130) ;
        %Straight Lines [id:da1927654285774716] 
        \draw    (201,100) -- (241,100) ;
        %Straight Lines [id:da6364192335175295] 
        \draw    (351,130) -- (391,130) ;
        %Straight Lines [id:da37849782165622403] 
        \draw    (391,100) -- (431,100) ;

        % Text Node
        \draw (156,173) node [anchor=north west][inner sep=0.75pt]   [align=left] {$\displaystyle 0$};
        % Text Node
        \draw (195,176) node [anchor=north west][inner sep=0.75pt]   [align=left] {$\displaystyle \varepsilon $};
        % Text Node
        \draw (232,173) node [anchor=north west][inner sep=0.75pt]   [align=left] {$\displaystyle 2\varepsilon $};
        % Text Node
        \draw (319,171) node [anchor=north west][inner sep=0.75pt]   [align=left] {$\displaystyle ( n-1) \varepsilon $};
        % Text Node
        \draw (382,176) node [anchor=north west][inner sep=0.75pt]   [align=left] {$\displaystyle n\varepsilon $};
        % Text Node
        \draw (406,171) node [anchor=north west][inner sep=0.75pt]   [align=left] {$\displaystyle ( n+1) \varepsilon $};
        % Text Node
        \draw (280,172) node [anchor=north west][inner sep=0.75pt]   [align=left] {...};

        \end{tikzpicture}
        \caption{Evolución temporal de $u^\varepsilon$.}
        \label{fig:f32}
    \end{center}
\end{figure}

Usamos 4 del lema anterior con $q = n$, y $p = n+1$. Tenemos que 
$$u^{n+1}(x) - u^n(x) = \int_{\R^2} K(x-y) \left( 
\underbrace{\omega^{n+1} ((n+1)\varepsilon, y)}_{\omega^{n+1}(n \varepsilon, X^n(n \varepsilon, (n+1)\varepsilon, y))} - \omega^{n} (n\varepsilon, y) \right)\,dy$$
y a su vez $\omega^{n+1}(n \varepsilon, X^n(n \varepsilon, (n+1)\varepsilon, y)) = \omega^{n}(n \varepsilon, X^n(n \varepsilon, (n+1)\varepsilon, y))$. Y ahora nos preguntamos 
por qué no podemos hacer la diferencia entre $\omega^{n}(n \varepsilon, X^n(n \varepsilon, (n+1)\varepsilon, y))$ y $\omega^n(n\varepsilon, y)$.
Se debe a que estamos en un problema hiperbólico y solo se transportan las normas, pero no podemos controlar la diferencia. Hacemos el siguiente 
cambio de variable: $z = X^n(n \varepsilon, (n+1)\varepsilon, y)$. Nos queda 
$$\int_{\R^2} K(x - X^n(n \varepsilon, (n+1)\varepsilon, z)) \omega^n(n \varepsilon, z) \overbrace{J^{-1}_{X^n}}^{1}\,dz - 
\int_{\R^2} K(x - \overbrace{y}^{z}) \omega^n(n \varepsilon, \overbrace{y}^{z})\,d\overbrace{y}^{z} = $$
$$\int_{\R^2} [K(x - \overbrace{X^n(n \varepsilon, (n+1)\varepsilon, z)}^{\psi = I_d}) - K(x-z)] \omega^n(n \varepsilon, z) \,dz$$
y aplicando el resultado anterior y conviertiendo las igualdades en desigualdades (por monotonía de la integral):
$$|u^{n+1}(x) - u^n(x)| \leq \int_{\R^2} 
\left| [K(x - \overbrace{X^n(n \varepsilon, (n+1)\varepsilon, z)}^{\psi = I_d}) - K(x-z)] \omega^n(n \varepsilon, z) \right|\,dz \leq$$
$$C \max_{z} |X^n((n+1)\varepsilon, n \varepsilon, z) - z| \{1 + |\ln \max_z |X^n((n+1)\varepsilon, n \varepsilon, z) - z||\}$$
Para ver cuánto vale ese máximo, reescribimos el argumento:
$$X^n((n+1)\varepsilon, n \varepsilon, z) - \overbrace{z}^{X^n(n \varepsilon, n \varepsilon, z)} \stackrel{TFC}{=} 
\int_{n \varepsilon}^{(n+1)\varepsilon} \dfrac{\partial}{\partial t} X^n(s, n \varepsilon, z)\,ds = 
\int_{n \varepsilon}^{(n+1)\varepsilon} u^n(s, X(s, n \varepsilon, z))\,ds$$
de donde
$$|X^n((n+1)\varepsilon, n \varepsilon, z) - z| \stackrel{(*)}{\leq} \int_{n \varepsilon}^{(n+1)\varepsilon} C \,ds = C \varepsilon$$
donde en $(*)$ se ha usado la propiedad $1$ del lema anterior. \\

\begin{observacion}
    Esta demostración es típica de examen.
\end{observacion}

% Demostración típica de examen esta dice. Entender lo que estamos haciendo.

Como $$\norma{\omega^0}_{L^1} = \norma{\omega_0 \ast \rho_\varepsilon}_{L^1} \leq \norma{\omega}_{L^1} \overbrace{\norma{\rho_\varepsilon}_{L^1}}^{1}$$
Y se puede cambiar $L^1$ por $L^\infty$. \\

Vamos a intentar aplicar ahora el Teorema de Ascoli-Arzelà, que recordemos que nos decía, dada una sucesión de funciones continuas
$\{f_\varepsilon\}_{\varepsilon \geq 0}$, que está uniformemente acotada y es equicontinua, entonces existe una parcial (que denotamos 
de la misma forma) tal que $\{f_\varepsilon\} \stackrel{c.u.}{\to} f$ sobre compactos. \\

Para nuestra sucesión, debe ser continua \textbf{en las dos variables}. En espacio lo es, pero no en tiempo, por lo que no podríamos 
aplicar el Teorema anterior directamente. Sí que será equicontinua, pero en tiempo no. Para poder aplicarlo, vamos a cambiar la regularización 
que hemos hecho. ¿Cómo podemos cambiar la $u^\varepsilon$ que tenemos hasta ahora para que sea continua? Lo más sencillo sería haciendo una 
poligonal entre puntos medios de cada intervalo $[n \varepsilon, (n+1)\varepsilon]$. Gráficamente hacemos lo que nos indica la Figura \ref{fig:f33}:

\begin{figure}[H]
    \begin{center}
        \tikzset{every picture/.style={line width=0.75pt}} %set default line width to 0.75pt        

        \begin{tikzpicture}[x=0.75pt,y=0.75pt,yscale=-1,xscale=1]
            %uncomment if require: \path (0,300); %set diagram left start at 0, and has height of 300

            %Straight Lines [id:da44261844071929823] 
            \draw    (150,160) -- (452,160) ;
            %Straight Lines [id:da6567160598379451] 
            \draw    (161,149) -- (161,172) ;
            %Straight Lines [id:da7983049400227261] 
            \draw    (201,149) -- (201,172) ;
            %Straight Lines [id:da21094039468564174] 
            \draw    (241,148) -- (241,171) ;
            %Straight Lines [id:da6065593239600682] 
            \draw    (351,149) -- (351,172) ;
            %Straight Lines [id:da8139488212809114] 
            \draw    (391,149) -- (391,172) ;
            %Straight Lines [id:da33702440515964105] 
            \draw    (430,148) -- (430,171) ;
            %Straight Lines [id:da8105430731321772] 
            \draw    (161,130) -- (201,130) ;
            %Straight Lines [id:da1927654285774716] 
            \draw    (201,100) -- (241,100) ;
            %Straight Lines [id:da6364192335175295] 
            \draw    (351,130) -- (391,130) ;
            %Straight Lines [id:da37849782165622403] 
            \draw    (391,100) -- (431,100) ;
            %Straight Lines [id:da9413423297139589] 
            \draw [color={rgb, 255:red, 255; green, 255; blue, 0 }  ,draw opacity=1 ]   (181,130) -- (221,100) ;
            %Straight Lines [id:da8346763639936287] 
            \draw [color={rgb, 255:red, 255; green, 255; blue, 0 }  ,draw opacity=1 ]   (221,100) -- (260,130) ;
            %Straight Lines [id:da8941574642140017] 
            \draw [color={rgb, 255:red, 255; green, 255; blue, 0 }  ,draw opacity=1 ]   (410.91,99.65) -- (369.01,129.93) ;
            %Straight Lines [id:da354821223796264] 
            \draw [color={rgb, 255:red, 255; green, 255; blue, 0 }  ,draw opacity=1 ]   (369.01,129.93) -- (332.08,97.42) ;
            %Straight Lines [id:da41995666922666564] 
            \draw [color={rgb, 255:red, 255; green, 255; blue, 0 }  ,draw opacity=1 ]   (411,100) -- (450,130) ;

            % Text Node
            \draw (156,173) node [anchor=north west][inner sep=0.75pt]   [align=left] {$\displaystyle 0$};
            % Text Node
            \draw (195,176) node [anchor=north west][inner sep=0.75pt]   [align=left] {$\displaystyle \varepsilon $};
            % Text Node
            \draw (232,173) node [anchor=north west][inner sep=0.75pt]   [align=left] {$\displaystyle 2\varepsilon $};
            % Text Node
            \draw (319,171) node [anchor=north west][inner sep=0.75pt]   [align=left] {$\displaystyle ( n-1) \varepsilon $};
            % Text Node
            \draw (382,176) node [anchor=north west][inner sep=0.75pt]   [align=left] {$\displaystyle n\varepsilon $};
            % Text Node
            \draw (406,171) node [anchor=north west][inner sep=0.75pt]   [align=left] {$\displaystyle ( n+1) \varepsilon $};
            % Text Node
            \draw (280,172) node [anchor=north west][inner sep=0.75pt]   [align=left] {...};
        \end{tikzpicture}
        \caption{Poligonales para tratar de obtener la continuidad de $u^\varepsilon$.}
        \label{fig:f33}
    \end{center}
\end{figure}

Así, la nueva definición de $u^\varepsilon$ será la siguiente:
\begin{equation*}
    \bar{u}^\varepsilon(t,x) = \begin{cases}
        u^0(x) \quad& t \in \left[0, \dfrac{\varepsilon}{2} \right]\\\\
        u^0(x) + \dfrac{t - \varepsilon/2}{\varepsilon} \left(u^1(x) - u^0(x) \right) \quad& t \in \left[\dfrac{\varepsilon}{2}, \dfrac{3\varepsilon}{2} \right]  \\\\
        \vdots \\
        u^j(x) + \dfrac{t - (j + \frac{1}{2})\varepsilon}{\varepsilon} \left(u^{j+1}(x) - u^{j}(x) \right) \quad& t \in \left[\left( j + \dfrac{1}{2} \right)\varepsilon, \left( j + \dfrac{3}{2} \right)\varepsilon \right]
    \end{cases}
\end{equation*}

\begin{lema}
    Sea $\omega_0 \in L^1 \cap L^\infty$, entonces 
    $$\norma{u^\varepsilon(t,x) - \bar{u}^{\varepsilon}(t,x)}_{L^\infty} \leq C \varepsilon |\ln \varepsilon| \stackrel{\varepsilon \to 0}{\to} 0$$
\end{lema}

Hemos cambiado $(u^\varepsilon, \omega^\varepsilon) \to (\bar{u}^\varepsilon, \omega^\varepsilon)$, y a esta sí se le puede aplicar Ascoli-Arzelà.

\begin{ejercicio}
    Comprobar el lema anterior.
\end{ejercicio}

\begin{ejercicio}
    Comprobar la equicontinuidad de $\bar{u}^\varepsilon$.
    $$|\bar{u}^\varepsilon(t,x) - \bar{u}^\varepsilon(s,x)| \leq C |t-s| \{1 + |\ln|t-s||\}$$
\end{ejercicio}

% 3.11.2026

% TODO: Revisar en lo de hoy cuándo la C es mayúscula y cuándo c es minúscula.

% TODO: Meter Dibujo 1 EDMB 3.11.2026

Teníamos que, para $t \in \left[ \left(j + \dfrac{1}{2}\right)\varepsilon, \left(j + \dfrac{3}{2}\right)\varepsilon \right]$, entonces:
% TODO: Añadir módulos para la acotación
$$u^j \left(j + \dfrac{3}{2} \right) - u^{j+1} \left(j + \dfrac{1}{2} \right) + t \left( \dfrac{u^{j+1} - u^j}{\varepsilon} \right) - $$
$$
\begin{cases}
    u^j \quad &t \in \left[ j \varepsilon, \left(j + \dfrac{1}{2} \right)\varepsilon \right[ \\\\
    u^{j+1} \quad &t \in \left] \left(j + \dfrac{1}{2} \right)\varepsilon, \left(j + 1 \right)\varepsilon \right]
\end{cases} = 
\begin{cases}
    (u^{j} - u^{j+1}) \left( j + \dfrac{1}{2} \right) + t \dfrac{u^{j+1} - u^j}{\varepsilon} \\\\
    (u^{j} - u^{j+1}) \left( j + \dfrac{3}{2} \right) + t \dfrac{u^{j+1} - u^j}{\varepsilon}
\end{cases} = $$
$$
\begin{cases}
    (u^{j+1} - u^{j}) \left( \dfrac{t}{\varepsilon} - \left( j + \dfrac{1}{2} \right) \right) \\\\
    (u^{j+1} - u^{j}) \left( \dfrac{t}{\varepsilon} - \left( j + \dfrac{3}{2} \right) \right)
\end{cases} \leq 
\begin{cases}
    \overbrace{C}^{2(j + 1/2)} |u^{j+1} - u^j| \\
    \overbrace{C}^{j + 2} |u^{j+1} - u^j|
\end{cases} < C \varepsilon(1 + |\ln \varepsilon|)$$
y podemos aplicar 4 del Lema \ref{lema:l31}. Podemos demostrar ahora que 
$$|\bar{u}^\varepsilon(t,x) - \bar{u}^\varepsilon(s,x)| \leq C |t-s| \{1 + |\ln|t-s||\}$$
Esta diferencia nos interesa calcularla cuando $t$ y $s$ estén cerca. Si están lejos nos va a dar igual. Esto se debe a que queremos 
un módulo de continuidad de $\bar{u}^\varepsilon$. Supongamos entonces que 
$$t,s \in \left[ \left(j + \dfrac{1}{2}\right)\varepsilon, \left(j + \dfrac{3}{2}\right)\varepsilon \right]$$
Entonces:
$$\left|u^j \left(j + \dfrac{3}{2} \right) - u^{j+1} \left(j + \dfrac{1}{2} \right) + t \left( \dfrac{u^{j+1} - u^j}{\varepsilon} \right) - 
\left[ u^{j} \left( j + \dfrac{3}{2} \right) - u^{j+1} \left( j + \dfrac{1}{2} \right) + s \left( \dfrac{u^{j+1} - u^j}{\varepsilon} \right) \right] \right| = $$
$$\left| (t-s) \left( \dfrac{u^{j+1} - u^j}{\varepsilon} \right) \right| \leq C |t-s| \dfrac{\varepsilon \{1 + |\ln \varepsilon|\}}{\varepsilon} \leq
C |t-s| \{1 + |\ln|t-s||\}$$

% TODO: Completar el segundo caso en que están en distintos tramos (2) del dibujo.

% TODO: Meter Dibujo 1 EDMB 3.11.2026

En un resumen de lo que conocemos hasta ahora, tenemos la siguiente proposición:

\begin{prop}
    Sea $\omega_0 \in L^1 \cap L^\infty$. Podemos construir, para cada $\varepsilon > 0$, el par $(\bar{u}^\varepsilon, \omega^\varepsilon)$ tal que:
    \begin{enumerate}
        \item $\norma{\omega^\varepsilon(t, \cdot)}_{L^1} = \norma{\omega^0}_{L^1} \leq \norma{\omega_0}_{L^1}$, y 
        $$\norma{\omega^\varepsilon(t, \cdot)}_{L^\infty} = \norma{\omega^0}_{L^\infty} \leq \norma{\omega_0}_{L^\infty}$$
        \item $\norma{\bar{u}^\varepsilon}_{L^\infty} \leq C$
        \item $|\bar{u}^\varepsilon(t,x) - \bar{u}^\varepsilon(t,y)| \leq c|x-y| \{1 + |\ln|x-y||\}$.
        \item $|\bar{u}^\varepsilon(t,x) - \bar{u}^\varepsilon(s,x)| \leq C|x-y| \{1 + |\ln|t-s||\}$.
    \end{enumerate} 
    donde $C = C(\norma{\omega_0}_{L^1}, \norma{\omega_0}_{L^\infty})$ es independiente de $\varepsilon$.
\end{prop}

\begin{observacion}
    Estas cuatro propiedades son estimaciones \underline{a priori}, debido a que las soluciones son aproximadas.
\end{observacion}

Estudiamos ahora la compacidad, con objeto a su vez de ver cuánto vale el límite de dicha sucesión. Para ello, vamos a aplicarle 
el Teorema de Ascoli-Arzelà a $\bar{u}^\varepsilon$. Recordemos que la sucesión $\{f_\varepsilon\}$ debe ser uniformemente acotada y 
equicontinua\footnote{Una sucesión de funciones $\{f_\varepsilon\}$, pongamos, $f_\varepsilon : I \to \R^N$, con $I$ un intervalo, es equicontinua si, para todo $\alpha > 0$, 
existe $\delta > 0$, con $|x-y| < \delta$ y $$|f_\varepsilon(x) - f_\varepsilon(y)| \leq \alpha \quad \forall x,y \in I$$}.
Entonces dicho teorema nos dice que existe una parcial $\{f_{\varepsilon'}\}$ tal que $f_{\varepsilon'} \stackrel{c.u.}{to} f$ en compactos. \\

Así, veamos primero que $\bar{u}^\varepsilon$ está uniformemente acotada. Dado $\alpha > 0$, nos preguntamos si existe $\delta > 0$ tal que $|x-y| < \delta$ y
$$|\bar{u}^\varepsilon(t,x) - \bar{u}^\varepsilon(t,y)| \leq C |x-y|\{1 + |\ln|x-y||\} < \alpha$$
Además, en la primera variable:
$$|\bar{u}^\varepsilon(t,x) - \bar{u}^\varepsilon(s,x)| \leq c |t-s|\{1 + |\ln|t-s||\} < \alpha$$

De esta manera, existe $\delta$ tal que $|(t,x) - (s,y)| < \delta$, y entonces se tenga la acotación anterior (en la primera variable). \\

Tenemos en principio una parcial, que denotaremos igual, y una existe una función $u$, tal que $\bar{u}^\varepsilon \stackrel{c.u.}{\to} u$
en compactos de $\left[0, T\right[ \times \R^2$. \\

Supongamos ahora que tenemos $\{f_\varepsilon\}$ tal que $\norma{f_\varepsilon}_{L^p} < C$, es decir, $f_\varepsilon$ está uniformemente acotada 
en $L^p$. Entonces, existe una parcial, $\{f_{\varepsilon'}\}$, y una función $f$, tal que
$$f_{\varepsilon'} \stackrel{L^p}{\rightharpoonup}_{\text{débil}} f$$ % TODO: Poner L^p debajo la flecha y justo debajo también débil
si $1 < p < \infty$ (dado que estos espacios $L^p$ son reflexivos). Recordemos que $\rightharpoonup$ denotaba la convergencia débil\footnote{Estos 
conceptos se han visto en Análisis Funcional.}. El caso $p=1$ no lo vamos a estudiar por ser más delicado. Si $p = +\infty$, entonces
$$\{f_{\varepsilon'} \stackrel{L^\infty}{\rightharpoonup} f\}$$ % TODO: Lo mismo que en la otra pero debajo débil*
Entonces para ampliar esto a la convergencia débil
$$\int f_{\varepsilon'} h \rightarrow \int fg \quad \forall g \in L^1$$
Todo esto lo usamos de la siguiente manera. Para nuestra $\omega^\varepsilon$, existe $\omega \in L^p$ tal que 
$$\omega^{\varepsilon'} \stackrel{L^p}{\rightharpoonup} \omega$$ % TODO: Lo mismo con débil debajo
$$\omega^{\varepsilon'} \stackrel{L^\infty}{\rightharpoonup} \omega$$ % TODO: Lo mismo con débil* debajo

\begin{ejercicio}
    ¿Son la misma $\omega$? Una converge débil y la otra converge débil*. De momento serán la misma $\omega$.
\end{ejercicio}

El siguiente paso consiste en demostrar que el par $(u, \omega)$ es solución de Euler 2D. No van a ser soluciones en el sentido clásico\footnote{
    porque $\omega$ no es derivable, ni $u$ tampoco. Veremos que el par $(u, \omega)$ es solución en el sentido débil}. Tenemos 
que demostrar lo siguiente:
\begin{equation}\label{eq:ec33}
    \begin{cases}
        \dfrac{\partial \omega}{\partial t} + (u \cdot \nabla) \omega = 0 \\
        u = K \ast \omega \\
        \omega(t=0,x) = \omega_0(x)
    \end{cases}
\end{equation}

El sentido débil consiste en tomar cualquier $\varphi \in C^1_0(\left[0, T\right[ \times \R^2)$. Ahora hacemos la integración por partes 
igual que en Modelos Matemáticos II, y nos queda:
$$\int_{0}^{T} \dfrac{\partial \omega}{\partial t} \varphi <= - \int_{0}^{T} \omega \dfrac{\partial \varphi}{\partial t} + 
\overbrace{\omega \varphi]_{0}^{T}}^{\omega(t=0,x) \varphi(0,x)}$$
y con esto se tiene que 
$$- \int_{0}^{T} \int_{\R^2} \left( \dfrac{\partial \varphi}{\partial t} \right) \omega = \int_{\R^2} \omega_0(x) \varphi(0,x)\,dx$$
En la primera ecuación de \eqref{eq:ec33}, podemos ver $(u \cdot \nabla)\omega = \di(u \omega)$, y la formulación débil de Euler es:
$$- \int_{0}^{T} \int_{\R^2} \left( \dfrac{\partial \varphi}{\partial t} + (u \cdot \nabla) \varphi \right) \omega\,dx\,dt = \int_{\R^2} \omega_0(x) \varphi(0,x)\,dx$$

% Seguimos la semana que viene
    %\include{Tema4.tex}
    %\include{Tema5.tex}

    %\fancyhead[R]{\helv \nouppercase{\rightmark}}
    %\chapter{Relaciones de Problemas}
    %\input{Relacion1}
    %\include{Relacion2}
\end{document}
